\pdfoutput=1
\ifdefined\pdftrailerid\pdftrailerid{}\fi
\documentclass[11pt]{article}
\usepackage[T1]{fontenc}
\usepackage{lmodern}
\usepackage[margin=.9in]{geometry}
\usepackage{amsmath,amssymb,amsthm,mathtools,microtype,xcolor}
\usepackage[colorlinks=true,allcolors=blue!50!black]{hyperref}
\newtheorem{theorem}{Theorem}[section]
\newtheorem{lemma}[theorem]{Lemma}
\newtheorem{proposition}[theorem]{Proposition}
\newtheorem{corollary}[theorem]{Corollary}
\theoremstyle{definition}
\newtheorem{definition}[theorem]{Definition}
\newtheorem{example}[theorem]{Example}
\theoremstyle{remark}
\newtheorem{remark}[theorem]{Remark}
\allowdisplaybreaks[1]
\hypersetup{pdftitle={Singular semiclassical sewing for paired-cap Witten--Reshetikhin--Turaev invariants}}
\title{Singular semiclassical sewing for paired-cap\\Witten--Reshetikhin--Turaev invariants}
\author{}
\date{}
\begin{document}
\maketitle
\begin{abstract}
We study the closed \(\mathrm{SU}(2)\) Witten--Reshetikhin--Turaev invariants of tree plumbings capped at each vertex by opposite order-two fillings. For every fixed tree with at least two vertices we prove a two-action power--logarithm expansion of the exact invariant and its smooth joint fibre-trace marks, with errors to any prescribed algebraic order. The two action sectors are selected using bounded finite-level cap insertions and have coefficients supported on the corresponding projected flat loci. On the classical side, a matching decomposition of the actual flat-character space gives the exact reduced Chern--Simons action and absolute torsion normalization on the irreducible leading Hall charts. The resulting leading coefficient agrees with the closed invariant. Fixed, compatibly routed Wilson probes in the dormant part of these charts obey an ungated leading joint angular law, including on trees with macroscopically active branches. A separate signed transition law records the relative scales near overlap strata. The all-orders statement concerns the radial action sectors; the angular comparison determines leading weights and specified overlap limits.
\end{abstract}
\tableofcontents
\section{Introduction and results}\label{intro:main}
Non-isolated flat connections are a basic obstruction to interpreting a semiclassical expansion as a sum of ordinary Gaussian saddle contributions. In a plumbing assembled from simple torus pieces, this issue already occurs in a sharp form: closed invariants can have leading logarithms, and the logarithmic order can grow with the number of pieces. The main question here is how classical critical geometry, absolute weight, and phase fit together in an actual finite quantum sum.

We restrict to connected finite untwisted trees with at least two vertices, with \(+2,-2\) caps at every vertex and the orientation-reversing section--fibre exchange on each backbone edge. The primitive section is the base boundary curve of the punctured-sphere product block. The twist/framing conventions agree with the earlier capped-tree calculations \cite{P48,P49,P49B,P4849}. We use \(r=k+2\geq 3\), \(S_{ab}=\sqrt{2/r}\sin(\pi ab/r)\), \(Z_r(S^2\times S^1)=1\), and \(\tau_r=Z_r/S_{11}\), so \(\tau_r(S^3)=1\). For a fibre colour \(a\), its trace is \(2\cos(\pi a/r)\), and its negative-centre deficit is \((r-a)/r\). Geometric cutoffs and Wilson insertions always refer to this specified marking and ribbon framing.

\paragraph{The classical family and its scale faces.}
The exact flat relations force noncentral fibres to form an independent set. They also give a combinatorial description of the maximal-dimensional actual solution branches. The vertices exposed by some maximum matching and their neighbours are removed first; the remaining forest is perfectly matched. This residual forest controls the number of leading logarithms. Nonzero scaled adjacent charge products belong to an ambient near-flat lift, rather than to the exact flat space itself. The geometric description retains both these scales and the conjugacy information that radial traces alone forget.

\paragraph{The exact invariant by action.}
For the finite modular contraction, the signed angular alias sums reduce to finite rational-cone sums. The forest incidence matrices provide the integrality needed for uniform expansions across intersections of additive walls. Cap-core insertions then yield bounded phase filters and identify the two actual Chern--Simons actions, \(0\) and \(1/2\) modulo \(1\). Their expansion coefficients are distributions in the original fibre traces. Each phase has a power--logarithm expansion without any further level-residue oscillation.

\paragraph{The absolute comparison.}
The reduced classical action in the relevant irreducible charts is obtained from the flat critical equations and the cyclic comparison of the finite deformation complexes; the latter is needed to identify the function, not merely its critical set. Primitive cohomology bases, the stabilizer, and the transverse Gaussian fix the normalization independently of the finite sum. The leading degeneracy localizes the matched angle dots at zero. On the quantum side the actual low-rank ribbon transfer, including its signed tails, then yields the same conditional joint law for the supported Wilson routes, including words connecting distinct dormant components. At fixed transition products the signed kernel is essential. Independent positive radial cutoffs would change the coefficient.

The paper keeps the two levels of conclusion separate: absolute classical weights and joint dormant-angle comparisons at leading order, together with selected transition laws; complete higher radial expansion by action with supported coefficient distributions. The latter leaves a further problem of assigning higher coefficients to individual angular components with the same radial data. This distinction is used in the statements and examples throughout.


\paragraph{Guide to the proof.}
There are three calculations with different roles. Sections~\ref{geo:sec} and~\ref{classic:sec:action} concern actual flat connections and the reduced classical action. The radial and phase sections start from finite modular sums and determine the signed expansions by action. The angular and Hall sections compare a stated class of genuine ribbon probes to the independently normalized classical model. The leading comparison does not use a conjectured expansion of a continuum functional integral.

\section{Global flat strata and the marked resolution of their overlaps}
\label{geo:sec}

Let \(G\) be the fixed backbone tree with \(N\geq2\) vertices. The notation in this section uses the positive Hopf peripheral convention and untwisted backbone. Let \(A_v,B_v\) be, respectively, the marked \(+2,-2\) cap meridians, and \(h_v\) the regular fibre. The product
\(P_v=\prod_{w\sim v}h_w\) is in the inherited section order. The local Seifert relations are \(A_v^2=h_v^{-1}\), \(B_v^2=h_v\), \(A_vB_vP_v=1\); the section--fibre exchange additionally identifies commuting edge pairs. In a tree every base transport can be chosen in the same spanning tree, so there are no stable letters, and simultaneous conjugation is the residual equivalence. In particular these are peripheral relations for the actual manifold, not its abelian plumbing presentation.

\subsection{Exact characters, mixed signs, and incidence}
Eliminating \(B_v=A_v^{-1}P_v^{-1}\), the full equations are
\begin{equation}
 h_v=A_v^{-2},\qquad [h_v,h_w]=1\quad(vw\in E(G)),\qquad
 \operatorname{Ad}_{A_v}P_v=P_v^{-1}.
 \label{geo:flat}
\end{equation}
Indeed the last equation is equivalent to the second cap square; the other requisite block commutators follow since neighboring fibres commute with \(h_v\). Use unit quaternions \(A_v=x_v+\mathbf a_v\). Then \(h_v=(2x_v^2-1)-2x_v\mathbf a_v\). At \(h_v=+1\), one has \(A_v=\pm1\) and \(P_v\) must be central. At \(h_v=-1\), one has \(A_v=u_v\in S^2\), and the last equation is
\(u_v\cdot \operatorname{Im}P_v=0\), including the absence of a condition when \(P_v\) is central. If \(h_v\) is noncentral then the last equation forces \(P_v\) central: each neighboring fibre and \(A_v\) lie in the centralizer circle of \(h_v\). In this case the two cap roots have signs \(\operatorname{sgn}x_v\) and are smooth functions of \(h_v\) on the two sign sheets.

\begin{theorem}[Global character atlas]\label{geo:character-atlas}
The set \(S\) of noncentral fibres is independent. Fix \(S\), the signs of the central fibres, and the cap-root choices at \(S\) and at positive central sites. Put \(T=N_G(S)\), let \(P\) be the positive central sites, and choose which negative sites of \(T\) have central rather than noncentral product; call that subset \(C\). Assign the product sign at every site in
\[
                        L=(P\cap T)\cup C.
\]
The labels are realizable exactly when each site of \(L\) has at least two neighbors in \(S\). Before quotient by conjugation their realization is a connected smooth manifold of dimension
\begin{equation}
 D=2N+|S|-|T|-2|P|-2|P\cap T|-2|C|.
 \label{geo:stratumdimension}
\end{equation}
The stabilizer is \(SU(2)\), a maximal torus, or the centre, according as the rank of the imaginary cap vectors is zero, one, or at least two. These rank refinements and their closure incidences are finite semialgebraic data. Within fixed fibre types, each negative active central product is a smooth codimension-two lock, and simultaneous such locks are independent.
\end{theorem}

\begin{proof}
A leaf of a nontrivial component of \(G[S]\) would have exactly one noncentral neighbor. Its product could not be central. This proves independence without excluding parallel axes.

For the realized labels use the incidence forest on variables \(h_s\), \(s\in S\), and equations \(P_v=\delta_v\), \(v\in L\). Root each component at an \(S\) variable, and at each equation vertex choose one child \(S\) variable. It exists because the equation has at least two incident variables. Solve uniquely for that child in the ordered product. This is a global word recursion on \(SU(2)^f\), \(f=|S|-|L|\). Every solved word contains its component root exactly once, with exponent \(+1\) or \(-1\); the remaining factors do not involve that root. It follows that each word is a submersion. An unlocked negative site in \(T\) meets distinct components of this incidence forest, by acyclicity, so its product is also a submersion by varying one component root. Remove all preimages of the two central values for the \(h_s\)'s and for these unlocked products. This removes a finite union of codimension-three closed submanifolds, leaving a nonempty path-connected base. For \(S=\varnothing\) the base is a point. An equation with exactly one noncentral input is impossible, establishing necessity of the label test.

Over the base attach a circle for each unlocked negative site in \(T\), the perpendicular cap equator, and a sphere at every other negative site. This is a smooth locally trivial connected bundle. Counting its dimensions gives \eqref{geo:stratumdimension}. Since all generators are words in the \(A_v\)'s, the common stabilizer is determined by their imaginary span. To see the lock model, allow independently variable targets in the same word recursion at the sites of \(C\). Near the assigned centre write \(p_v=\operatorname{Im}(\delta_vP_v)\). The actual product equation is \(p_v\in u_v^\perp\); its lock is the zero section in this real two-plane bundle. These target coordinates are independent by the recursion. The central limits of fibre variables, unlike fixed-fibre product locks, need not satisfy the incidence rule obtained by merely weakening strict inequalities.
\end{proof}

Here is an exact way to retain all closing directions when fibre types change. Over a marked quaternion tuple use the simultaneous rotation invariants
\[
              x_i,\qquad d_{ij}=\mathbf a_i\cdot\mathbf a_j,\qquad
              t_{ijk}=\det(\mathbf a_i,\mathbf a_j,\mathbf a_k).
\]
Gram data give an isometry of the spans; a marked nonzero triple fixes its orientation in rank three, and in lower rank one may extend the isometry with positive orientation. Thus these data embed the compact character quotient. They are literal Wilson words since
\begin{align}
 \tfrac12\operatorname{Tr}(A_iA_j)&=x_ix_j-d_{ij},\nonumber\\
 \tfrac12\operatorname{Tr}(A_iA_jA_k)
 &=x_ix_jx_k-x_i d_{jk}-x_j d_{ik}-x_k d_{ij}-t_{ijk}.
 \label{geo:wilsontraces}
\end{align}
For a fixed pair of nonparallel cap pivots \(i,j\), let
\(a=|\mathbf a_i|\), \(b=d_{ij}/a\), \(c=(|\mathbf a_j|^2-b^2)^{1/2}\).
The three coordinates of \(\mathbf a_\ell\) in the common oriented pivot frame are
\[
                  d_{i\ell}/a,\quad
                  (d_{j\ell}-b d_{i\ell}/a)/c,\quad
                  t_{ij\ell}/(ac).
\]
They use the same frame on every approaching branch; the denominators need not vanish when a different fibre enters a centre. At a reducible with a nonzero cap axis, use one pivot instead: the transverse complex coordinates \(z_\ell\) are recovered modulo their common phase by the Hermitian Gram \(\overline z_\ell z_m\), whose real part is the perpendicular dot and whose imaginary part is \(t_{i\ell m}/a\). The rank-zero vertex of this cone is retained.

The raw equations and all label/rank tests are polynomial in the lifted quaternion coordinates. Thus membership in the closure of a sector with predicate \(F_\lambda\), at an invariant tuple \(q\), is exactly the real quantifier test
\[
 \forall\delta>0\ \exists A:\quad F_\lambda(A),\quad
                          |(x,d,t)(A)-q|^2<\delta^2 .
\]
Quantifier elimination gives a finite semialgebraic incidence test; common closure and sector overlap are intersections of these tests. One can refine them to smooth frontier strata. This is a decision procedure, not a claim of a closed combinatorial expression for every pointwise closure.

For example, on \(P_4\) take \(S=\{1,3\}\), with \(h_2=h_4=-1\), and hold the two active radii and their cap signs fixed at equal noncentral angles. Taking \(h_3=h_1^{-1}\) makes \(P_2=+1\); tilting \(h_3\) makes \(P_2\) noncentral and permits a perpendicular cap at vertex 2. The local cap and fibre traces agree even between these different sectors. The literal ordered-product trace \(\operatorname{Tr}P_2\) detects the lock; pair and triple marks also distinguish its angular limits. Thus pointwise angular localization is not furnished by radial and cap-root labels alone.

\begin{theorem}[Abelian anchors and components]\label{geo:components}
Every sector in Theorem~\ref{geo:character-atlas} can approach every central fibre sign assignment compatible with its locked products and fixed root signs, and every connected component meets an abelian character. The components are exactly those of the following finite graph. A vertex is an all-central label at each backbone site: \(N\) for negative, \(P_+\) or \(P_-\) for positive with the indicated cap sign. An edge toggles \(N\leftrightarrow P_\pm\) at a nonempty independent set of vertices and is permitted precisely when every unchanged positive vertex has an even number of toggled neighbors.
\end{theorem}

\begin{proof}
In the word recursion, choosing each free holonomy central solves all allowed central sign assignments. Approach a chosen central tuple by \(g_j(t)=g_j(0)\exp(tX_j)\) with generic \(X_j\). Each excluded word has surjective derivative, so every such imaginary derivative can be chosen nonzero simultaneously. The open base then approaches the desired tuple, with limiting emitted directions and hence limiting perpendicular cap axes. Connect to the arc from the original point in the connected base and lift in the circle/sphere bundles. At an all-central tuple negative cap axes are free; align them continuously to attain an abelian character.

Two central assignments of one sector differ on an independent subset of \(S\); at every unchanged positive site the product parity must agree. Conversely a permitted toggle set gives a sector with those fibres noncentral, the required locked positive products, and unlocked negative products. Its admissibility follows from Theorem~\ref{geo:character-atlas}, so its closure joins the two indicated central labels. An actual incidence of sectors preserves membership in this graph component: the limiting noncentral set shrinks, original positive roots and original locked products persist, and a formerly noncentral root that becomes positive has its inherited sign. Therefore every allowed central endpoint of the limit sector is also an allowed endpoint of the original one. Partitioning the character space by graph component of its endpoint labels gives closed sets with closed complements by the finite incidence argument. Each is connected by the sector closures and the connected all-central sets of fixed labels. This proves the assertion.
\end{proof}

\subsection{Intrinsic Hall matching and the tangent deficit}
Let \(I=G\setminus(S\cup T)\) be the inactive forest, and \(p=|P\cap T|\). At a negative centre let \(V_v=\mathbb R u_v\), and at a positive centre let \(V_v=\mathbb R^3\). The block adjacency \(\mathcal B_I\) has, on an edge \(vw\), the linear map \(V_w\hookrightarrow\mathbb R^3\to V_v\) given by inclusion and orthogonal projection. We write \(h^j=\dim H^j(M_G;\operatorname{ad}\rho)\). In particular \(h^0\) is the Lie dimension of the stabilizer.

\begin{theorem}[Tangent rank at every character]\label{geo:tangent}
At every character, including locks, collinear axes, and continuous stabilizers,
\begin{equation}
            h^1-h^0=3N-3-2|T|-4p-\operatorname{rank}\mathcal B_I .
\label{geo:h1}
\end{equation}
Generically the rank in \eqref{geo:h1} is \(2\nu_b(I)\), where \(\nu_b\) is the maximum integral edge matching with vertex capacities \(b_v=3\) at positive centres and \(b_v=1\) at negative centres.
\end{theorem}

\begin{proof}
Differentiate \eqref{geo:flat} in right-trivialized coordinates, stripping central signs from fibre derivatives. The map \(A_s\mapsto h_s\) at \(s\in S\) has invertible derivative, giving three variables per noncentral fibre. At a central vertex the fibre variables lie in \(V_v\), with two additional cap-axis variables precisely at negative vertices.

An active-centre fibre variable lies along the axis of each noncentral neighbor by the edge commutator. A possibly nonzero such column therefore has all these axes on one line, and at a negative centre also its own cap axis on that line. Every such column has degree at least two in the incidence forest with row vertices \(S\): degree one is excluded at a positive site by its central-product condition, and at a negative one the cap axis would be perpendicular to its single noncentral neighbor. At \(s\in S\) the linearization forces its ordered-product derivative zero, because \(\operatorname{Ad}_{A_s}+I\) is invertible for a noncentral fibre. This is a scalar sum of the incident collinear columns. A nonempty support of nonzero columns in the resulting forest has a leaf on the row side, which would force the sole column there to vanish. Hence all active-centre fibre variables vanish.

The \(p\) active-positive products impose three equations apiece. An active-negative product uses one cap-axis variable if noncentral, and otherwise uses one component of its product derivative. These latter derivatives together with all active-positive ones are independent, by the word recursion with freely varied targets. The inactive radial equations are exactly \(\mathcal B_I\xi=0\), since their neighbors outside \(I\) have zero fibre derivative and they have no noncentral neighbors. Thus, with \(n_-\) the number of negative centres, the cocycle count is \(3|S|-3p+2n_--(|T|-p)+\dim\ker\mathcal B_I\). Subtract \(3-h^0\) for infinitesimal conjugations to obtain \eqref{geo:h1}.

For the generic rank, split \(\mathcal B_I\) bipartitely. In its rectangular block a nonzero determinant term gives a matching of coordinate slots, with the stated capacities. Conversely color the copies of edges in a capacity-maximizing multiset by three colors, greedily down the forest with distinct colors at every vertex. Specialize each selected negative axis to its assigned coordinate vector and use the matching coordinate vectors at positive vertices. The selected minor splits into one forest per color; each has a perfect matching and no second one, so its determinant is nonzero. Thus the polynomial minor is nonzero generically in the independent inactive axes. The same argument applies to rectangular forest maps with local domains or ranges of dimensions one and three, used below.
\end{proof}

The capacitated matching number is the minimum over covers \(A\) of \(I\) of
\(|A|+2|A\cap P|\), by integral bipartite max-flow/min-cut \cite{FordFulkerson1956}. Let \(\nu=\nu(G)\), \(\delta=N-2\nu\), and
\[
                         p_G=(3N-2\nu-3)/2
\]
be the power in the physical \(Z_r\) leading theorem. The generic tangent deficit is exactly
\begin{equation}
 p_G-\tfrac12(h^1-h^0)
 = (|T|+\nu(I)-\nu(G))+2p+\bigl(\nu_b(I)-\nu(I)\bigr)\geq0.
 \label{geo:genericdeficit}
\end{equation}
A matching uses at most \(|T|\) edges incident to \(T\), and after deleting \(T\) only \(I\) contributes; this proves the first nonnegativity. The last term vanishes precisely when an ordinary minimum cover of \(I\) avoids all positive vertices. At exceptional axes the exact rank in \eqref{geo:h1} replaces its generic value. Neither an enlarged tangent space nor generic equality in \eqref{geo:genericdeficit} is a stationary-phase amplitude.

\begin{theorem}[Leading-dimensional smooth sectors and their closure]\label{geo:matching}
A smooth fixed-label sector attains the maximal dimension minus Lie stabilizer \(2p_G\) only when
\begin{equation}
 P=C=\varnothing,\qquad |S|-|N_G(S)|=\delta .
 \label{geo:hall}
\end{equation}
The conditions are sufficient on a generic piece, which is smooth in the full character space. For any fixed-label sector, a generic point lies in the closure of such leading-dimensional smooth sectors if and only if
\begin{equation}
 p=0,\qquad \nu(G)=|T|+\nu(I),\qquad
 \text{a minimum vertex cover of \(I\) avoids \(P\)}.
 \label{geo:hallclosure}
\end{equation}
If any test in \eqref{geo:hallclosure} fails, no point of that label is in the closure, including exceptional tangent rank drops. Genericity in this theorem concerns only a finite list of rank tests on the inactive negative axes.
\end{theorem}

\begin{proof}
Fix a maximum matching. The matched members of independent \(S\) inject into \(T\) by their partners, so \(|S|-|T|\leq\delta\), with equality exactly when \(S\) contains all unmatched vertices and \(T\) consists of the matched partners of its matched members. In that case the matching is perfect on \(I\). Formula \eqref{geo:stratumdimension}, followed by quotient at fixed stabilizer, proves the clean dimension necessity. If \eqref{geo:hall} holds, \(\mathcal B_I\) is generically invertible by Theorem~\ref{geo:tangent}; the fixed sector already exhausts the tangent count. A maximal-rank subset of the defining equations cuts out a local implicit-function manifold of this dimension, in which the smooth fixed sector of the same dimension is open. Thus no other local branch is lost. These sectors exist: the complement of a minimum cover is independent, every vertex of the cover meets its complement, and the matching equality follows. For nonempty \(S\) an unlocked active negative product forces irreducibility; for \(S=\varnothing\) generic negative axes are irreducible.

For closure necessity an approaching leading sector has a fixed independent set \(S\cup J\) after passing to a subsequence, with \(J\subset I\) and \(P\subset J\). Thus \(p=0\). The equality in \eqref{geo:hall} gives
\[
 \delta=|S|-|T|+|J|-|N_I(J)|
       \ \leq\ |S|-|T|+|I|-2\nu(I)\ \leq\ \delta .
\]
Both inequalities are equal. The partners of matched \(J\) exhaust \(N_I(J)\), and the remaining forest is perfectly matched. Append a minimum cover of that forest to \(N_I(J)\) to obtain a minimum cover of \(I\) disjoint from \(J\), and hence from \(P\).

To prove sufficiency, choose a minimum cover \(A\) of \(I\) avoiding \(P\), and a maximum matching \(M\) that saturates \(A\) into its independent complement. Start at all unmatched vertices and all positives in that complement. Reach \(J\) by successively adding every neighbor in \(A\) and its partner in the complement. Put \(A'=N_I(J)\). The set \(A'\) consists exactly of the matched partners of the matched part of \(J\); hence \(|J|-|A'|=|I|-2\nu(I)\). If the root set is empty this part of the construction is unnecessary.

At each \(j\in J\) write \(h_j=\epsilon_j\exp X_j\), with \(\epsilon_j\) the given centre and \(X_j\in\mathbb R u_j\) at negatives, \(X_j\in\mathbb R^3\) at positives. The cap roots can be chosen smoothly to approach the prescribed ones. At \(v\in A'\), which is negative, hold its cap axis fixed and impose the exact equation \(u_v\cdot\operatorname{Im}P_v(X)=0\). The derivative is the rectangular forest map \(W\) with \(v\)-row \(u_v\cdot\sum_{j\sim v} X_j\). The retained matching saturates its rows, so \(W\) has full row rank generically. We need a kernel vector with every block \(X_j\) nonzero and with \(\sum_{j\sim v}X_j\ne0\) at every row, not merely a dimension estimate.

For a negative \(j\), deleting its column still permits a saturating matching of the rows: if it is reached, exchange along an alternating path from a root to expose it. An unmatched root has a spare slot, and a positive root has two slots beyond its ordinary matched slot. The same reasoning covers a root which is itself unmatched. The rank test in Theorem~\ref{geo:tangent} now shows a nonzero kernel projection on this column. Deleting the three columns at a positive site loses at most one row rank, using its ordinary match, so its kernel projection is nonzero as well. At any \(v\in A'\), replacing its scalar row by the full vector sum raises the rank generically by at least one: follow an alternating path from a root and augment at the additional slot of \(v\); if \(v\) is the matched positive root's immediate partner the augmentation is a second copy of the same edge. The rectangular minor criterion applies. Therefore on \(\ker W\) none of the required block or row-vector tests vanish identically. Avoid their finite union of proper linear subspaces. The implicit function theorem now integrates the resulting velocity to an actual curve opening every site of \(J\) and making every new \(A'\) product noncentral.

Old active products that were already noncentral keep a smooth cap equator. Old central active-negative products can be prescribed small noncentral paths in their cap planes: apply the exact word recursion to old \(S\) inputs with the new fibers inserted at their marked places. This does not change equations at \(A'\), which have no neighbors in old \(S\). All remaining inactive centres are negative, since the root set for \(J\) included all positives. The resulting open \(S\cup J\) has Hall deficiency \(\delta\) by the two equalities above; its inactive forest is perfectly matched. Finitely many additional generic nonzero minor conditions at the original inactive axes make it smooth by the first part. This proves the real closure assertion.
\end{proof}

\begin{corollary}[The forced hull and the logarithm count]\label{geo:forced}
Let \(S_*\) be the set of vertices which are exposed by some maximum matching, \(T_*=N_G(S_*)\), and \(I_*=G\setminus(S_*\cup T_*)\). The set \(S_*\) is Hall-maximal and is the intersection of all sets satisfying the matching equality in \eqref{geo:hall}. The forest \(I_*\) is perfectly matched. The exact WRT leading logarithm degree is \(d=|I_*|/2\).
\end{corollary}
\begin{proof}
Use the preceding alternating reach with only unmatched roots. Exchange along an even alternating path to expose each reached vertex. Conversely the symmetric difference of two maximum matchings consists of balanced alternating paths (an unbalanced path would augment one of them); this puts every newly exposed vertex in the reach. The reach is itself Hall-maximal and is contained in every Hall-maximal independent set, which must contain all unmatched roots and the partners of all their reached neighbors. Thus all quantities are independent of the chosen matching and cover. In the minimum-cover contraction of the exact quantum coefficient, the forced cover vertices are the ancestors of the neighbors of the unmatched roots, via a cover vertex's partner and then another cover neighbor. This is exactly \(T_*\) by the same reach. The quantum log degree \(\nu-|T_*|\) is therefore \(|I_*|/2\). This is a combinatorial equality between two independently defined counts, not a reassignment of individual oscillatory weights.
\end{proof}

\subsection{The ambient overlap tube and the radial face forest}
The exact flat character space is not by itself a carrier for literal nonzero edge-product overlap scales. Indeed, independent \(S\) gives
\(\delta_v\delta_w=0\) on every edge of every flat character, where
\(\delta_v=|\operatorname{Im}h_v|\). Any graph closure of these coordinates on the exact flat space, with \(\varepsilon=1/r>0\), still has \(\delta_v\delta_w/\varepsilon\equiv0\). In contrast a marked finite product threshold retains \(n_vn_w/r\) nonzero while both endpoint angular defects may tend to zero. The relevant enlargement here is an ambient approach to flatness, not a putative extra chart on the exact character space.

On the ambient quotient of all unit cap quaternions keep \(h_v=A_v^{-2}\), the same ordered \(P_v\), and \(B_v=A_v^{-1}P_v^{-1}\). Set
\[
 R_{vw}=[h_v,h_w],\qquad R_v=A_vP_vA_v^{-1}P_v,
   \qquad p_v=\operatorname{Im}P_v .
\]
The useful estimates are in fact identities:
\begin{align}
 \|R_{vw}-1\|
   &=2|\operatorname{Im}h_v\times\operatorname{Im}h_w|
     =8|x_vx_w|\,|\mathbf a_v\times\mathbf a_w|,\nonumber\\
 \|R_v-1\|^2
   &=4\bigl(x_v^2|p_v|^2+(\mathbf a_v\cdot p_v)^2\bigr).
 \label{geo:errors}
\end{align}
For an edge multiply the difference by unit factors to get
\(h_vh_w-h_wh_v\); for a vertex multiply by \(P_v^{-1}\) to get the imaginary vector \((\operatorname{Ad}_{A_v}+I)p_v\), with singular values \(2\) axially and \(2|x_v|\) transversely. The second cap-square error has the same norm, by multiplying \(B_v^2-h_v\) on the left by \(A_v^2\) and inverting \(P_v\) in the displayed comparison. The error in commutation of \(B_v\) with \(h_v\) is bounded by twice that error, as \(B_v\) commutes with \(B_v^2\).

All factors in \eqref{geo:errors} are finite marked Wilson data, using \eqref{geo:wilsontraces} and
\[
 |p_v|^2=1-(\operatorname{Tr}P_v/2)^2,\qquad
 \mathbf a_v\cdot p_v=x_v\operatorname{Tr}P_v/2-\operatorname{Tr}(A_vP_v)/2.
\]
The cross norm is a Gram minor. In particular \(\|R_v-1\|=O(\varepsilon)\) retains a true cap-plane lock near a negative centre, but near a positive centre, where \(|x_v|\) is bounded below, it forces the full vector \(p_v=O(\varepsilon)\). This single residual description distinguishes both cases. It is gauge invariant. For fixed marked generators, replacing the finite relator presentation by an equivalent one changes uniform \(O(\varepsilon)\) bounds only by fixed constants, since a relator expressed as a finite product of conjugates of other relators is bounded by their summed errors.

\begin{theorem}[Global radial overlap tube]\label{geo:ambient}
Let \(\epsilon_*=\max(\|R_{vw}-1\|,\|R_v-1\|)\) over the graph. Uniformly on the ambient compact cap quotient,
\begin{equation}
          \max_{vw\in E}\delta_v\delta_w\ \leq\ C_G\epsilon_*.
 \label{geo:globalproduct}
\end{equation}
At the negative endpoint, write \(h_v=-\exp(s_v u_v)\) and take angular deficits \(d_v=|s_v|/\pi\) in a fixed small endpoint collar. The projection of the \(O(\varepsilon)\) relator tube to power rates \(d_v\asymp\varepsilon^{\alpha_v}\), \(0\leq\alpha_v\leq1\), is exactly
\begin{equation}
                         \alpha_v+\alpha_w\geq1\quad(vw\in E).
 \label{geo:rates}
\end{equation}
At any such rate vector, every assignment of finite positive values for \(d_vd_w/\varepsilon\) on a chosen tight subforest can be realized. In the physical log-defect variables \(t_v=1-\alpha_v\) the projected polytope is the exact quantum radial polytope \(t_v+t_w\leq1\). No angular law or signed weight is asserted for this ambient projection.
\end{theorem}

\begin{proof}
Put \(D=\max_{vw}\delta_v\delta_w>0\), \(m=|E|\), and \(B=8(\Delta+1)\), where \(\Delta\) is the maximum degree. The subforests of edges with product at least \(D/B^j\), \(0\leq j\leq m\), have two consecutive equal members, since the first is nonempty and there are at most \(m-1\) strict increases. Take a leaf \(v\) of a nontrivial component of that forest and its sole selected neighbor \(w\). At every other neighbor \(j\), \(\delta_j\leq\delta_w/B\). If \(D>B^{m+1}\epsilon_*\), the selected edge product exceeds \(B\epsilon_*\); its axis angle has sine at most \(1/(2B)\) by \eqref{geo:errors}. In the ordered \(P_v\), replace each other \(h_j\) by its nearest centre with error at most \(\sqrt2\,\delta_j\), and telescope the unit factors. Thus
\[
             p_v=\pm\operatorname{Im}h_w+O(\sqrt2\,\Delta\,\delta_w/B).
\]
Both its magnitude and the absolute component along the \(v\) cap axis are at least \(\delta_w/2\). By \(x_v^2+|\mathbf a_v|^2=1\), the vertex identity forces \(\|R_v-1\|\geq\delta_w>B\epsilon_*\), contradiction. This proves \eqref{geo:globalproduct}, also for parallel limits. A cycle does not have the required leaf: on a quadrilateral, commuting noncentral \(h_1=h_3^{-1}\), \(h_2=h_4^{-1}\) satisfy both kinds of relator with nonzero adjacent products.

Necessity of \eqref{geo:rates} follows since \(\delta_v\) and \(d_v\) are uniformly comparable in the endpoint collar. For sufficiency choose two orthogonal fixed axes for the bipartition of \(G\), put \(h_v=-\exp(s_v u_v)\) on its assigned axis and \(A_v=\pm u_v\exp(-s_v u_v/2)\). Neighbors at each vertex now lie in a common opposite torus, so their ordered product is a central sign times \(\exp(L_v u_w)\), \(L_v=\sum_{w\sim v}s_w\). The two errors are exactly \(2|\sin s_v\sin s_w|\) and \(2|\sin L_v\sin(s_v/2)|\). The latter is bounded by \(|s_v|\sum_{w\sim v}|s_w|\). Taking the indicated powers realizes every rate. This does not say every prescribed angular configuration with those rates is in the tube.

Along each tight subforest, \(\alpha\)'s alternate as \(\alpha,1-\alpha\). Prescribed constants \(w_e>0\) are realized by solving, successively from a root, \(\log c_v+\log c_w=\log w_e\) for the coefficients \(d_v\sim c_v\varepsilon^{\alpha_v}\). These equations are independent with one free scale in each connected component. If \(\alpha=0\) occurs there, it is confined to one bipartition and that free scale makes its fixed defects fit the collar. This proves the finite-constant statement; the \(O(\varepsilon)\) bound can depend on those chosen constants.
\end{proof}

There is a corresponding signed model showing why the positive strata are not additional copies of this negative forest. In the same orthogonal-axis construction let
\(h_v=\epsilon_v\exp(s_vu_v)\), \(\epsilon_v=\pm1\), and use \(A_v=\pm\exp(-s_vu_v/2)\) at positive sites. At a positive site the vertex norm is now
\(2|\sin L_v\cos(s_v/2)|\). For strictly positive rate exponents \(\alpha_w\), prescribe the signs of \(s_w\). In this two-axis model a signed rate assignment obeying \eqref{geo:rates} is realizable exactly when every positive site with minimum neighboring exponent less than one has both signs among neighbors of that minimum rate. Necessity is the leading term of \(L_v\). For sufficiency, solve \(L_v=0\) exactly at these rows. The relevant incidence graph splits into the forests of one row color and the opposite variable color. Root each at a variable; for each row solve for one child of minimum exponent. When the incoming variable has minimum exponent choose an opposite-sign child, taking other minimum coefficients small; when it is higher order choose two opposite minimum children and solve for one of them. At least two minimum variables of opposite signs exist by the test, and all solved variables are new children. This gives the desired signs and rates by a finite recursion, with exact cancellations. No such scalar criterion is asserted for arbitrary axes.

\subsection{A relative marked compactification}
For a finite inventory of marks use the raw invariants in \eqref{geo:wilsontraces}, ordered-product and mixed traces, Gram rank minors, both endpoint distances, and the error factors in \eqref{geo:errors}. Directions at a closing lock or emitted axis, such as \(p_v/|p_v|\) or \(\mathbf a_v/|\mathbf a_v|\) where defined, are retained by their Gram/triple data with the other marked vectors. Stabilizer rank can be recorded both before and after these directions are added. For the exact boundary point its phase is the phase of its base flat character; no off-shell Chern--Simons value is assigned by this construction.

\begin{proposition}[Finite relative overlap models]\label{geo:compactification}
Fix any finite list of comparisons \(M/(M+N)\) of nonnegative monomials of marked magnitudes and \(\varepsilon\). Include, in particular, comparisons of the raw edge products and relator errors with \(\varepsilon\). On the finite support and sign sheets, the marked graph closures give minimal compact extensions over the corresponding raw closures on which these prescribed ratios and directions extend. They are covered by finitely many normal-crossings charts. On each chart the finite visible scaling rates are linear images of a rational polyhedron; actual finite ratio constants belong to the marked boundary, not to the rates alone.
\end{proposition}

\begin{proof}
Embed the simultaneous-conjugation quotient by its quaternion invariant data. Use separate support sheets to distinguish an undefined direction or a comparison with both entries identically zero, from any approached nonzero direction or ratio. The graph closure in the product of the compact raw space, direction data and a ratio cube is compact. A compact extension with these same maps and the given dense sheet maps onto that closure by joint evaluation, uniquely by density. This proves minimality relative to the chosen marks; forgetting marks is compatible with common refinements of the support sheets.

The graph data and all label predicates have finite semialgebraic descriptions; for norms add their nonnegative square-root graph coordinates. Endpoint angular distance is an analytic positive unit times the semialgebraic proxy \(|\operatorname{Im}h|\) in each endpoint neighborhood. The exact ratio after multiplying two magnitudes by units \(u,v>0\) is recovered from proxy ratio \(R\) as \(uR/(uR+v(1-R))\). Thus the proxy does not discard finite comparison data.

Apply real embedded normal crossings to the product of the finite nonzero polynomials defining the graph sheets and the chosen magnitude and rank coordinates, including \(\varepsilon\), as in real rectilinearization \cite{BierstoneMilman1988}. After pullback each factor is a signed unit times a monomial; each Boolean sign predicate is consequently constant on its relevant coordinate orthants, including the lower-dimensional ones. Keep only closures of lifts from the original positive-level sheet. Properness supplies lifts of all its actual graph limits and excludes artificial boundary components; compactness gives a finite cover. The dominating charts are not claimed canonical.

On each such orthant a positive marked magnitude is \(u_\ell z^{m_\ell}\), with \(u_\ell>0\), and the same holds for \(\varepsilon\) with row \(m_\varepsilon\). Thus the normalized logarithmic rates are
\[
 m_\varepsilon\cdot\lambda=1,\qquad
 \lambda_j\geq0,\qquad \operatorname{ord}_\varepsilon f_\ell=m_\ell\cdot\lambda.
\]
A convergent sequence lifts after a subsequence, giving these equations for all finite visible rates; an infinitely fast coordinate invisible to every retained exponent row does not affect them. Conversely choose power paths in the vanishing coordinates, with subpower approaches for a prescribed vanishing coordinate of zero rate. One obtains the indicated rates and retains the actual angular and ratio limits on the graph. A common polyhedral refinement makes the finite rate images compatible. Ties of rates alone do not distinguish finite nonzero ratios from zero or infinite limits at subpower speed; this is why the raw marks and ratio coordinates are retained.
\end{proof}

Theorem~\ref{geo:ambient} therefore identifies the exact radial projection of a concrete ambient approach to flatness, and Proposition~\ref{geo:compactification} retains its intersection and stabilizer data. Neither implies that a signed WRT remainder is supported in the \(O(1/r)\) tube. In particular divergent residual ratios can remain on other approach faces over a flat point, and literal angular insertions need not be diagonal in the unmarked color contraction. The analytic comparisons below require estimates on the actual sewing kernels rather than a replacement by the geometry alone.

\section{All algebraic orders of the signed radial sum}
\label{rad:section}

We first establish existence and a coefficient rule on the two level
parities, before assigning coefficients to Chern--Simons actions.
Beyond paths, the truncation walls at a branching vertex involve
signed sums of several neighbouring colours. Such walls can intersect
multiple endpoint poles. We neither replace them by an ordering of
colours nor discard nonzero aliases.

Let \(G=(V,E)\) be a fixed tree with \(N=|V|\ge2\), and put \(r=k+2\).
Use reflected one-based colours \(a_v=r-n_v\), \(1\le n_v<r\).
We write a radial test as \(F((n_v/r)_v)\). A smooth function of the
original fibre traces is of this form, with \(F\) extending smoothly
by even reflection at both endpoints of each coordinate. This
endpoint condition, more generally, suffices for our argument.
For \(d\ge1\) put \(W_d(x)=(1+\cos\pi x)/\sin^d(\pi x)\). For a fixed
parity vector \(\boldsymbol p\in\{0,1\}^V\) the signed kernel is
\begin{equation}
 \mathcal W_{\boldsymbol p,r}[F]=
 \sum_{\substack{1\le n_v<r\\n_v\equiv p_v\;(\mathrm{mod}\,2)}}
 F((n_v/r)_v)
 \prod_{v\in V}W_{\deg(v)}(n_v/r)
 \prod_{vv'\in E}\sin\!\left(\frac{\pi n_vn_{v'}}r\right).
 \label{rad:sum}
\end{equation}
The actual untwisted paired-cap amplitude is a finite signed linear
combination of these kernels times \((2/r)^{(N+1)/2}\).
The vacuum contribution uses only odd \(n_v\). A cap-product tag
\(q\) uses half the odd weight plus \((-1)^q\) times half the even
weight. Specifying an ordered pair of roots with that product
adds a factor \(1/2\). These are the exact cap identities in the
modular tree contraction \cite{GPPV}, not continuum replacements.
Reflection of an edge contributes \((-1)^{r+n_v+n_{v'}}\), fixed
on each parity choice. The untwisted surgery signature is zero:
the two opposite caps have cancelling Schur contributions, and the
unweighted bipartite backbone adjacency has spectrum symmetric
about zero.

\begin{theorem}[All-order radial expansion]\label{rad:thm}
For each fixed tree, each fixed smooth joint radial test, and each
fixed choice of physical cap projectors, the amplitude admits an
expansion on each parity \(r\equiv\rho\pmod2\). Its terms are
\(r^\alpha P_{\alpha,\rho,F}(\log r)\), where \(\alpha\) lies in a
fixed half-integral translate of \(\mathbb Z\), the power scale is
bounded above, and log degrees are bounded in terms of the tree.
For every \(J>0\), retaining all terms with \(\alpha\ge-J\) leaves
\(O(r^{-J})\); taking extra orders does not change these coefficients.
No other arithmetic subsequence is needed. The coefficients are
continuous linear functionals of finite order in \(F\). Finite
identities between the original signed marked sums hold termwise
in their expansions.
\end{theorem}

We prove this for \eqref{rad:sum}, then restore the prefactor and
signs. Multiplication by the ordinary asymptotic Taylor series of
\(S_{11}^{-1}\) gives the corresponding result for \(\tau_r=Z_r/S_{11}\).
Removing residual parity after the physical action is extracted
is a separate step.

\subsection{Character elimination: one affine wall per branch}
\label{rad:elimination}

Choose a bipartition \(A,B\) of \(G\). Expand \(F\) in its smooth
\(2\)-periodic Fourier series. By individual evenness we can use
the product of cosines at a mode \(\boldsymbol\sigma\in\mathbb Z^V\),
with the original joint Fourier coefficient
\(\widehat F(\boldsymbol\sigma)\). These coefficients decrease
rapidly. For any fixed \(\epsilon>0\), modes
\(\|\boldsymbol\sigma\|>r^\epsilon\) contribute arbitrarily small
powers of \(r^{-1}\), since the original finite sum in absolute
value is bounded by a fixed power of \(r\) for all unit-bounded tests.

At fixed \(A\)-colours and mode, the \(B\)-variables decouple.
At \(b\in B\) with \(d=\deg b\), its sines and denominator combine
into unnormalized characters, with continuous endpoint values:
\[
 \frac{\sin(\pi n_i m/r)}{\sin(\pi m/r)}
  =\sum_{j_i=0}^{n_i-1}
       \exp\!\left(\frac{i\pi(2j_i-n_i+1)m}r\right),
 \qquad i\in N(b).
\]
Put \(h_b(x)=(1+\cos\pi x)\cos(\pi\sigma_b x)\) and let \(h_{b,s}\)
be its coefficients in the basis \(e^{i\pi sx}\). Only
\(s=\pm\sigma_b,\pm\sigma_b\pm1\) can occur. Reflecting
\(m\mapsto 2r-m\) gives equal contributions. The possibly included
\(m=r\) term is zero, whereas the full-period even grid adds
\(m=0\). Exact finite orthogonality
\[
 \sum_{\substack{m\;(\mathrm{mod}\,2r)\\m\equiv p\;(\mathrm{mod}\,2)}}
 e^{i\pi\ell m/r}=r(-1)^{p\ell/r}{\bf1}_{r\mid\ell}
\]
with the right side interpreted as zero off \(r\mid\ell\), gives
the \(b\)-factor
\begin{equation}
 \begin{aligned}
 &\frac r2\sum_s h_{b,s}\sum_{\lambda\in\mathbb Z}
 (-1)^{p_b\lambda}{\bf1}_{L_b\in\mathbb Z}
 \sum_{U\subseteq N(b)}(-1)^{|U|}
 \binom{L_b-\sum_{i\in U}n_i+d-1}{d-1}_{+}\\
 &\quad-{\bf1}_{p_b=0}\frac{h_b(0)}{2}\prod_{i\sim b}n_i,
 \qquad L_b=\tfrac12\left(\lambda r-s+\sum_{i\sim b}(n_i-1)\right).
 \end{aligned}
 \label{rad:alias}
\end{equation}
The subscript \(+\) sets a binomial to zero if its upper argument
is below \(d-1\). Indeed the lattice condition requires
\(\sum_{i\sim b}j_i=L_b\), and excluding each event \(j_i\ge n_i\)
gives the displayed binomial. The final term subtracts the
artificial even-grid origin and must be kept for individual
Fourier modes, even if its aggregate vanishes. The contributing
aliases \(\lambda\) lie in a fixed finite interval if
\(|\sigma_v|\le r\), since each character frequency is bounded
by \(n_i-1\).

Expand the product of \eqref{rad:alias} as a signed sum. An
individual term selects one frequency, alias and subset, or
the origin correction, at each \(b\). Thus it keeps at most
\emph{one affine wall per \(b\)}: no common arrangement of all
competing walls is needed. Set \(r=2R+\rho\) and \(n_i=2j_i+p_i\)
for the retained labels. The integrality test for \(L_b\)
concerns only the fixed parities and the parity of \(s\).
For these allowed choices the binomial is a polynomial times
\begin{equation}
  \lambda_b R+\sum_{i\in N(b)}(1-2{\bf1}_{i\in U_b})j_i
              +c_b(\rho,\boldsymbol p,s)\ \ge0 ,
 \label{rad:wall}
\end{equation}
where the affine constant is integral. The other hard constraints
are \(n_i\ge1,\ r-n_i\ge1\). The remaining trigonometric factors
are smooth numerators times integer powers of those positive
endpoint distances, since \([x(1-x)]^dW_d(x)\) is smooth on
\([0,1]\). Polynomial weights and retained cosines multiply
this macro numerator. This reduction includes high-valence
walls as well as all signed aliases.

\subsection{A lattice decomposition with integral level step}
\label{rad:cone}

Fix the signs of the Fourier shifts and set
\(\upsilon_v=|\sigma_v|\). Homogenize \eqref{rad:wall}, the
endpoint bounds, and \(0\le\upsilon_v\le 2R+\rho\) in
\((R,(j_i)_{i\in A},(\upsilon_v)_{v\in V},H)\),
multiplying constants by \(H\). Add \(R\ge H\ge0\); the physical
slice has \(H=1\), and finitely many lowest levels are immaterial.
Impose the shift congruences on the integer lattice coset.
At \(R=0\) all variables vanish. The rational cone has compact
cross-section at \(R=1\). Pulling-triangulate this cross-section
with its original vertices and take the disjoint relative
interiors of all the resulting simplicial cones, including
lower-dimensional faces.

In one relative interior, a lattice point has a unique expression
\(u+\sum_i k_i e_i\) with independent integral ray vectors \(e_i\),
\(k_i\in\mathbb Z_{\ge0}\), and
\(u=\sum_i\theta_i e_i,\ 0<\theta_i\le1\). Scale rays if needed to
preserve the shift congruences. Only finitely many \(u\) occur.
At fixed \(H=1,\upsilon=|\boldsymbol\sigma|\), any \(k_i\) of a
ray carrying positive \(H\) or shift is \(O_G(Q)\), where
\(Q=1+\|\boldsymbol\sigma\|\), because these components are
nonnegative. Enumeration of those coefficients gives at most
polynomially many base points of size \(O_G(Q)\). Free rays have
\(H=\upsilon=0\).

Every free ray can be taken with \emph{unit \(R\)-component}.
Indeed, on the recession face \(H=\upsilon=0,\ R=1\), the
coefficient matrix in the \(j_i\) has entries \(\pm1\) supported
on a subforest of the original incidence graph, with at most
one row from each \(b\). A square minor cannot have two
contributing perfect matchings, whose symmetric difference
would form an alternating cycle. Consequently all square
determinants are \(0,\pm1\), and adjoining the endpoint unit
rows preserves this property. The right sides at \(R=1\) are
integral, so every recession vertex is integral by the determinant
formula for an independent subsystem of its active constraints.
As \(H=\upsilon=0\) is a face, these vertices are among the
original vertices of the full cross-section. The triangulation
adds no rays, and a free ray needs no congruence scaling since
its shifts vanish. This is where the forest and the single-wall
expansion avoid a new arithmetic period.

The lattice pieces at a fixed \(R\) are therefore simplexes
\begin{equation}
 k_i\in\mathbb Z_{\ge0},\qquad \sum_i k_i=R-R_0 .
 \label{rad:simplex}
\end{equation}
Pieces without a free ray occur only when \(r=O_G(Q)\), outside
the retained Fourier box for large \(r\). Each endpoint distance
is a positive affine form in the free \(k_i\), with base value
at least one and nonnegative slopes, by the recession box bounds.
On each simplex use a smooth partition of unity in barycentric
coordinates, choosing a pivot \(\ell\) with
\(k_\ell/(R-R_0)\ge1/(2m)\) on its support, where \(m\) is the
number of free rays. Eliminate \(k_\ell\) with
\eqref{rad:simplex}. The pivot support is a fixed normalized
distance from the omitted boundary and extends by zero; the
remaining lattice variables now range independently over
nonnegative integers. An endpoint distance seeing the pivot
with positive slope is \(\gtrsim r\) on the support for
\(r\gg_G Q\), so its inverse power gives a known power of \(r\)
and a smooth factor. Other endpoint forms retain their
nonnegative slopes.

Up to explicit powers of \(r\), the resulting sums are
\begin{equation}
 \sum_{k\in\mathbb Z_{\ge0}^d} f(tk,t)
       \prod_{j=1}^{D}\left(b_j+\sum_i a_{ji}k_i\right)^{-p_j}.
 \label{rad:orthant}
\end{equation}
Here \(t=r^{-1}\), \(b_j\ge1\), \(a_{ji}\ge0\), and
\(p_j\in\mathbb Z_{>0}\). The \(a_{ji},p_j\) belong to fixed finite families, and
\(b_j=O_G(Q)\). The profile has a smooth extension to a fixed
compact macro box for \(tQ\ll1\): the eliminated boundary is
avoided, the other orthant walls allow ordinary smooth extension,
and all offsets and retained frequencies are \(O_G(Q)\).
Every fixed derivative order thus costs at most a power of \(Q\).

\subsection{The positive affine orthant lemma}
\label{rad:orthantsec}

Smooth polytope Euler--Maclaurin formulas \cite{BV} and singular
lattice formulas for point-type weights \cite{BK} address related
settings. Here the needed estimate instead follows by resolving
a continuous Laplace kernel after the lattice sum is performed
exactly.

\begin{lemma}\label{rad:lemma}
The sum \eqref{rad:orthant} has an all-order expansion in integer
powers of \(t\) and polynomials in \(\log(1/t)\) of degree at most
\(D\). For each fixed output order both coefficients and remainder
bounds have at most polynomial cost in \(Q\). The power scale is
bounded below uniformly in the base and Fourier mode.
\end{lemma}

\begin{proof}
It suffices initially to take a \(t\)-independent profile \(f(x)\).
The general case follows by Taylor expansion in \(t\) with
polynomial-in-\(Q\) bounds in each specified derivative norm.
A remainder \(O(t^M Q^{C_M})\) costs only \(O(t^{-d})\) when summed
over the original compact support since each affine denominator
is at least one.

For \(g(x)=e^{\sum_i x_i}f(x)\), Fourier-invert its Schwartz
transform, Laplace-transform the denominator powers, and sum
the independent geometric series. The initial exchanges converge
absolutely for \(t>0\). Up to the Fourier convention, the sum is
the integral against \(\widehat g(\xi)\) of
\begin{equation}
 K(t,\xi)=\int_{y_j\ge0} e^{-\sum_j b_jy_j}
       \prod_j\frac{y_j^{p_j-1}}{\Gamma(p_j)}
       \prod_{i=1}^d\frac{1}{1-e^{-z_i}}\,dy .
 \label{rad:laplace}
\end{equation}
Here \(z_i=\sum_j a_{ji}y_j+t(1-i\xi_i)\). The real \(t\)-damping gives the convergent geometric sum and
prevents cancellation of its leading real scale. Since
\(\Re z_i\ge t\) and \(b_j\ge1\), we have
\(|K(t,\xi)|\lesssim (1-e^{-t})^{-d}\) uniformly in \(\xi\).
Insert a smooth cutoff \(\chi(t(1+|\xi|))\), equal to one near
zero and supported below \(c<1/2\). The change after Fourier
integration is \(O(t^M Q^{C_M})\) for any prescribed \(M\),
by the kernel bound and Schwartz decay.

Split the Laplace coordinates at one. In each region order all
small \(y_j\)'s and \(t\) as \(1\ge z'_1\ge\cdots\ge z'_L\),
ignoring null boundaries. Introduce ratios
\(z'_\ell=v_1\cdots v_\ell\), \(0\le v_\ell\le1\). If \(t\)
is in position \(u\), then \(t=\prod_{\ell\le u}v_\ell\).
Denominators seeing a large \(y_j\) with positive coefficient
have bounded smooth inverses with controlled derivatives.
For any other denominator let \(m(i)\) be the position of the
largest relevant variable, including \(t\). The quotient
\(z_i/z'_{m(i)}\) is smooth in the ratios with uniformly
positive real part, since one ratio is one with positive
coefficient. On the cutoff \(|\Im z_i|<c\), so
\(z/(1-e^{-z})\) is smooth, including at zero. It follows
that \(z'_{m(i)}/(1-e^{-z_i})\) extends smoothly to the
ratio cube, with derivative bounds polynomial in \(1+|\xi|\)
at every fixed order. Differentiating the Laplace weights
costs powers of \(Q\), and the large-coordinate integrals
are dominated by \(e^{-\sum_{\rm large}y_j}\) times polynomials.
The smooth cutoff allows extension by zero outside its region.

For the pushforward use \(dt/t\) and \(dy_j/y_j\) on the small
coordinates. The logarithmic change to the ratios has
determinant one. The monomial weight is
\(\prod_\ell v_\ell^{\beta_\ell}\), where
\begin{equation}
 \beta_\ell=\sum_{\substack{j\ {\rm small}\\{\rm pos}(j)\ge\ell}}p_j
    -\#\{i:m(i)\text{ exists and }m(i)\ge\ell\}.
 \label{rad:beta}
\end{equation}
The presence of \(t\) in every denominator gives \(m(i)\le u\).
Hence \(\beta_\ell>0\) for \(\ell>u\), and those ratios and
their required derivative bounds are integrable. Each sector's
Mellin transform in \(t\) is therefore
\begin{equation}
 \int_{[0,1]^u}\prod_{\ell=1}^{u}v_\ell^{w+\beta_\ell}
          A(v,\xi)\prod_{\ell=1}^{u}\frac{dv_\ell}{v_\ell}.
 \label{rad:mellin}
\end{equation}
Here \(A\) is smooth on the closed cube with the stated
seminorm bounds, and vanishes near the all-one corner by the
\(t\)-cutoff.

To control the remainder, apply a sufficiently long Taylor
subtraction at each \(v_\ell=0\) and expand the product of
Taylor and remainder operators. Polynomial terms supply
factors \((w+\beta_\ell+j)^{-1}\); the remaining integrals
are absolutely bounded on the prescribed strip. Poles thus
lie at integers and have order at most \(u\).
For vertical decay, split \(A\) into a finite smooth partition
on each piece of which some \(v_i\) is bounded strictly below
one. Such pieces cover its support. On the initial convergence
half-plane, repeated integration by parts in \(v_i\) with
\(v_i\partial_{v_i}\) has no boundary terms and divides each
time by \(w+\beta_i\). The differentiated amplitudes satisfy
the same Taylor/remainder strip estimates. Meromorphic
continuation gives arbitrary vertical decay on each strip
away from the poles with polynomial bounds in \(Q,1+|\xi|\)
at each fixed order.

Shift Mellin inversion to \(\Re w=-J-\tfrac12\). Residues
give integer powers of \(t\) and log-degree at most
\(u-1\le D\); the shifted integral gives \(O(t^{J+1/2})\)
with the polynomial parameter bounds. Only finitely many
poles are crossed, and \(\beta_\ell\ge-d\) gives the uniform
upper location bound on the poles. Integrating against
\(\widehat g\) preserves the expansion and its estimates.
\end{proof}

\subsection{Coefficient assembly and checks}
\label{rad:checks}

Apply Lemma~\ref{rad:lemma} to \eqref{rad:orthant} on
\(Q\le1+r^\epsilon\) with any fixed \(0<\epsilon<1\),
taking sufficiently many orders before restoring explicit
powers of \(r\). The number of base pieces and all fixed-order
bounds cost only polynomial powers of \(Q\), so the weighted
\(\widehat F(\boldsymbol\sigma)\) sum preserves the remainder.
The coefficient tails decay arbitrarily fast by the
polynomial bounds, and the omitted exact finite-sum tail
does so by the crude absolute bound on \eqref{rad:sum}.
Only finitely many smooth seminorms of \(F\) occur at a fixed
output order. This proves Theorem~\ref{rad:thm}.

A coefficient rule is now explicit: expand \eqref{rad:alias},
triangulate the affine cells, form \eqref{rad:laplace}, extract
the appropriate Mellin residues from \eqref{rad:mellin}, and
sum the Fourier coefficients absolutely. The exact partitions,
even-grid subtraction and signed inclusion--exclusion account
for the overlaps. Uniqueness of ordered power--log expansions
on each parity subsequence shows independence from the
triangulation, pivot partition and growing Fourier box.

\paragraph{Path check.}
For the unmarked four-vertex path, eliminating a leaf and a
degree-two vertex gives the exact expression
\begin{equation}
 Z_r(P_4)=(-1)^r\sqrt{\frac2r}
 \sum_{\substack{m,n<r\\ m,n\ {\rm positive\ odd}}}
 W_2(m/r)W_1(n/r)\min\{m,n,r-m,r-n\}.
 \label{rad:p4}
\end{equation}
The corresponding odd-grid \(W_1\) sine sum is \(r/2\);
the \(W_2\) two-sine sum is \(r/2\) times the displayed folded
minimum. These follow from \eqref{rad:alias}; for the latter
the zero and reflected aliases count the overlap of two
intervals of consecutive weights. The prefactor is
\((2/r)^{5/2}(r/2)^2=\sqrt{2/r}\). This checks the full alias
normalization, not a minimum-kernel hypothesis at higher
valence.

\paragraph{Intersecting branch walls.}
Let the retained set comprise \(a_0,\ldots,a_4\). The eliminated
vertices are two trivalent sites, with neighbours
\(\{a_0,a_1,a_2\}\) and \(\{a_0,a_3,a_4\}\), and an extra
leaf on \(a_0\). Both parts contain a trivalent site. In the
unaliased odd palette write \(n_i=2j_i+1\). If \(j_0\ge j_1+j_2\),
the first zero-weight count is
\((2j_1+1)(2j_2+1)\). In the adjacent regime
\[
 j_0=j_1+j_2-d,\qquad 0\le d\le 2\min(j_1,j_2),
\]
it decreases by \(d(d+1)\),
because each excluded signed tail has count \(1+\cdots+d\).
The second branch has the analogous wall \(j_0=j_3+j_4\).
They give simultaneous additive walls at the endpoint corner.
The shared \(a_0\) column is precisely the forest pattern
of Section~\ref{rad:cone}; outside the zero-weight window
the formula still includes all aliases. Neither part satisfies
a degree-two hypothesis, and no path estimate is needed at
the intersection.

The radial restriction is substantive. Angular Wilson
insertions can carry recoupling not represented by any
fixed smooth function of fibre traces; they need not reduce
to \eqref{rad:alias}. Backbone twists likewise insert
quadratic phases beyond this filter. This theorem also does
not assign different angular components with the same radial
projection their individual amplitudes.

\section{Physical phase separation and the support of radial coefficients}
\label{phase:section}

We use the original colour coordinate \(x_v=a_v/r\) in this section, so that the measured fundamental trace is \(2\cos(\pi x_v)\). In particular the positive central fibre is the endpoint \(x_v=0\). The reflected deficit is \(n_v=r-a_v\). All filters in this section are fixed smooth functions of the trace variables; in angle coordinates they extend smoothly and evenly across each endpoint.

\subsection{The action and the bounded cap-product selectors}
\label{phase:classical}

Write \(N_-,P,S\) for the sets of negative central, positive central and noncentral fibres of a flat character. We use \(E_{--}\) and \(E_{S-}\) for the numbers of edges of the indicated types. At a positive central site the two actual cap roots have product \((-1)^{q_v}I\). The classical action is
\begin{equation}
 \operatorname{CS}=\tfrac12\left(E_{--}+E_{S-}+\sum_{v\in P}q_v\right)
 \quad\bmod\mathbb Z .
 \label{phase:cs}
\end{equation}
Here and below the level factor is \(e^{2\pi i(r-2)\operatorname{CS}}\). At either value in \eqref{phase:cs} this equals \(e^{2\pi i r\operatorname{CS}}\).

By Theorem~\ref{geo:components}, every flat character has a path to an all-central abelian anchor with its locked products preserved. We may compute the action at that anchor because the variation of Chern--Simons vanishes along a flat path.

For the abelian endpoint use any real eigenphase lift \(y\) in units \(2\pi\) for the full plumbing matrix \(Q_{\rm full}\), so that \(Q_{\rm full}y\) is integral. The fundamental abelian Chern--Simons invariant is the torsion linking self-pairing, up to the orientation sign; in these lifts it is \(y^{\mathsf T}Q_{\rm full}y\bmod\mathbb Z\), up to that sign \cite{Auckly}. This formulation also works for singular \(Q_{\rm full}\): the integral vector \(Q_{\rm full}y\) lies in the real image of \(Q_{\rm full}\), and changing a lift by an element of \(\ker Q_{\rm full}\) leaves the pairing unchanged. In the positive Hopf convention let \(b_v=1\) for a negative endpoint fibre and \(0\) otherwise, and \(L_v\) be the sum of neighbouring \(b\)'s. Take the lifts
\[
 y_v=-b_v/2,\qquad
 y_{v,+}=b_v/4+j_v/2,\qquad y_{v,-}=-y_{v,+}+L_v/2.
\]
Here \(j_v\) records the cap-root choice. The spine entries of \(Q_{\rm full}y\) are zero and its two cap entries are \(j_v\) and \(j_v-L_v\). Therefore
\[
 y^{\mathsf T}Q_{\rm full}y\equiv\sum_v(b_vL_v/4-L_v^2/2)
   \equiv \tfrac12(E_{--}+E_{-+})\pmod{\mathbb Z}.
\]
The orientation sign does not matter for a half-integer. An original noncentral vertex incident to a negative vertex contributes half an edge at the anchor regardless of its own central limit sign. At an originally positive vertex the cap product preserves the parity of its negative endpoint neighbours. Constancy of the action along flat paths now gives \eqref{phase:cs}.

It is useful to isolate the action by bounded physical operations. Choose at each vertex a smooth \(\chi_v\) supported away from \(x_v=1\), equal to one near \(x_v=0\), and valued in \([0,1]\). At a site marked by \(\chi_v\), the two cap reciprocity summands can be combined to select the actual cap-product parity \(q_v=0,1\). Explicitly their selected factor, before the powers of \(S_{1a}\) are absorbed, is
\begin{equation}
 A_q(a)=\frac2r\sin^2\!\left(\frac{\pi a}{2r}\right)(-1)^{(r-a+1)q}.
 \label{phase:tag}
\end{equation}
Thus \(A_0+A_1\) is the ordinary cap weight on odd deficits and zero on even deficits.

Here is the Wilson realization of this selection. For an order-\(b\) cap, \(b=\pm2\), write \(C_{b}^{(m)}=(S\Theta^bS)_{a m}\) with cap-core colour \(m\) and with the gluing framing. Its two reciprocal terms \(c^{(m)}_{b,j}\), \(j=0,1\), have, with \(\vartheta=\pi a/(2r)\),
\[
 c^{(1)}_{b,j}=(-1)^{j(r-a+1)}c^{(1)}_{b,0},\qquad
 c^{(2)}_{b,j}=e^{-3\pi i/(2rb)}\,2(-1)^j\cos\vartheta\,c^{(1)}_{b,j}.
\]
The second identity is the sine double-angle formula in the reciprocal sum. The first follows by translating its lift by \(2r\), including the lift's quadratic phase. For the vacuum colour the two opposite order caps have conjugate \(j=0\) terms with squared modulus \(r^{-1}\sin^2\vartheta\). Taking their separately selected \(j\)'s and summing the two choices with \(j_++j_-\equiv q\) gives \eqref{phase:tag}. In particular the single-cap term is
\[
 \tfrac12\left(C_b^{(1)}+
      \frac{(-1)^j C_b^{(2)}}{e^{-3\pi i/(2rb)}\,2\cos\vartheta}\right).
\]
The centralizing roots have traces \(2(-1)^j\cos\vartheta\) on the chart where the fibre is not the negative centre; their product at a positive centre is indeed \((-1)^q I\). The divisor in this expression is uniformly bounded away from zero on the chosen \(\chi_v\) supports. These are signed linear combinations of framed Wilson insertions with a diagonal fibre multiplier, not positive or orthogonal projectors.

Expand the exact tensor partition \(\prod_v(\chi_v+(1-\chi_v))\). Let \(T\) be the sites carrying \(\chi\), \(C=V\setminus T\), and write \(Z_r(T,q;F)\) for the contraction with the chosen \(q_v\)'s at \(T\) and an additional joint smooth trace test \(F\). Set
\begin{equation}
 p(T,q)=|E(G[C])|+\sum_{v\in T}q_v\pmod2,\qquad
 Z_r^{[p]}[F]=\sum_{p(T,q)=p}Z_r(T,q;F).
 \label{phase:split}
\end{equation}
The sum over both \(p\)'s is exact at each level. On the support of a given tag, actual positive centres are in \(T\) and negative centres in \(C\). A tagged noncentral vertex has \(q_v\) equal mod two to the number of its negative neighbours, since all its neighbours are central. Its contribution to \(E_{S-}\) in \eqref{phase:cs} can therefore be included in \(\sum_T q_v\). This proves \(2\operatorname{CS}=p(T,q)\) on every flat consistent with these marks.

The arithmetic phase can also be removed directly from each finite contraction. Reflect the \(C\)-colours to their fixed odd deficits \(n\), and retain the full original grids \(a\) on \(T\). Each \(C\)--\(C\) sine gives \((-1)^r\), and each \(C\)--\(T\) sine gives \((-1)^{a_v+1}\) at its tagged end. Together with \eqref{phase:tag} the extracted factor is
\begin{equation}
 (-1)^{rp(T,q)}\prod_{v\in T}(-1)^{(a_v+1)\lambda_v},
 \qquad
 \lambda_v=q_v+\#\{w\sim v:w\in C\}\pmod2.
 \label{phase:sign}
\end{equation}
Only the exterior phase depends on a moving parity. The remaining grids are fixed; the alternations in the second factor live on the fixed full grids.

\subsection{Projected flat support, including higher valence}
\label{phase:support}

At a positive tagged site of degree \(d\), hold the neighbouring angles \(\theta_i=\pi a_i/r\) fixed. Let \(\nu_\theta\) be probability Haar on the \(SU(2)\) class of angle \(\theta\), and let \(\mu_\theta=\nu_{\theta_1}*\cdots*\nu_{\theta_d}\), using group convolution. Put
\[
 m(u)=\eta(u)u^{d-2}\frac{1-\cos(\pi u)}{\sin^d(\pi u)},\qquad
 K_{r,m}(g)=\sum_{a\ge1} a\,m(a/r)\chi_a^{SU(2)}(g).
\]
Here \(\eta\) is the near-positive trace filter, possibly including a joint radial test evaluated with the other colours fixed. The function \(m\) is smooth, compactly supported away from \(u=1\), and has a smooth even extension at zero. After extraction of the common prefactors and the additional \(1/2\) in the tagged cap factor, the exact local sum is
\begin{align}
 &\sum_{a=1}^{r-1}\eta(a/r)\frac{1-\cos(\pi a/r)}{\sin^d(\pi a/r)}
           (-1)^{(r-a+1)q}\prod_i\sin(a\theta_i) \nonumber\\
 &\hspace{10mm}=
   (-1)^{rq}r^{d-2}\prod_i\sin\theta_i\,
                  (K_{r,m}*\mu_\theta)((-1)^q I).
 \label{phase:convolution}
\end{align}
Indeed the central Fourier coefficient of \(\nu_{\theta_i}\) in the \(a\)-dimensional irreducible is \(\chi_a^{SU(2)}(\theta_i)/a\) by Schur averaging, while \(\chi_a^{SU(2)}((-1)^q I)=a(-1)^{(a-1)q}\). This is the probability normalization of the class-convolution formula in \cite[(29)--(32)]{AMW}, applied to the actual finite cap selector.

Away from identity \(K_{r,m}\) is \(O(r^{-M})\) for every \(M\), uniformly on each fixed compact separated set. To see this, its numerator at angle \(\phi\) is \(\sum_{a\ge1} a m(a/r)\sin(a\phi)\). The odd smooth extension of \(u\,m(u)\) and ordinary Poisson summation give rapid decay away from \(2\pi\mathbb Z\), with the same estimate for any fixed number of \(\phi\)-derivatives. At the antipode the numerator vanishes, so the derivative estimate also controls the division by \(\sin\phi\). Smooth parameter dependence of joint filters is uniform on the compact angle domains.

\begin{proposition}[Physical projected support]
\label{phase:offsupport}
For a given tag set and cap products, the radial projection of the flat locus consists exactly of the following data:
\begin{enumerate}
 \item the noncentral sites are independent, and every tagged noncentral site has its forced parity \(q_v=\#\{w\sim v:h_w=-I\}\bmod2\);
 \item at every positive centre the product of the incoming \(SU(2)\) conjugacy classes contains the target \((-1)^{q_v}I\).
\end{enumerate}
For a fixed smooth marked support disjoint from these data the finite tagged contraction is \(O(r^{-M})\) for every \(M\).
\end{proposition}

\begin{proof}
Necessity follows from the cap equations. For sufficiency, choose at each positive centre a tuple realizing its prescribed class product. The incidence graph between positive and noncentral sites is a forest. Having fixed one incoming noncentral element, simultaneously conjugate the next positive-site tuple so that its already assigned class representative agrees. The rooted tree propagation never inherits two such constraints. At a negative centre a pair of traceless roots can realize any given product \(g\): choose the first imaginary unit perpendicular to \(\operatorname{Im}g\), and the second root determined by the product is again an imaginary unit. At noncentral sites the incident fibres are central, and there are exactly the root choices with the forced product. These give an actual tree representation. The local class-product condition can, in particular, be checked by the \(SU(2)\) inequalities of \cite[Theorem 2.2]{JeffreyMare}.

For the rapid estimate, expand the sine factors into signed exponentials where the noncentral labels are strictly interior. The slope of an untagged noncentral coordinate has to be integral in units of \(\pi r\) for the step-two grid, and the slope of a tagged coordinate before inserting its alternating factor has to be \(q_v\) modulo two for the step-one grid. A leaf of a nontrivial component of noncentral sites sees one strictly interior angle and otherwise integer endpoint angles. Every signed slope is bounded away from the corresponding lattice on a small neighbourhood. The same holds for an incorrectly tagged noncentral with only central neighbours. Discrete summation by parts in that eligible coordinate gives arbitrary inverse powers; other weights have at most fixed polynomial size at the lattice endpoints. If the class product instead fails at a positive centre, compactness of the class product gives a fixed distance from the central target on a neighbourhood. Apply \eqref{phase:convolution} and the rapid group kernel estimate at that site. A finite smooth partition in the trace coordinates on the disjoint support combines the estimates, without any endpoint differentiation of the other factors.
\end{proof}

The class product in Proposition~\ref{phase:offsupport} cannot be replaced by an abelian signed-angle equation. Three classes of angle \(\pi/2\) multiply to either centre by choosing \(i,j,\mp k\), but no signed sum of the three angles is \(0\) or \(\pi\) modulo \(2\pi\). Both targets are regular: the annihilator of the differential of the product map is the intersection of the transported axis lines, which vanishes unless all three axes coincide. The same test also explains why two different CS values can have the same fibre radii on a positive trivalent star.

It now follows that the two phase sums in \eqref{phase:split} are independent of the choice of near-positive cutoffs to every algebraic order. Change one \(\chi_v\). Its difference is compactly supported inside the noncentral interval. Expand the untagged contribution there as \(A_0+A_1\). The two terms with \(v\) tagged and untagged are assigned the same phase exactly when \(q_v=\#\{w\sim v:w\in C\}\bmod2\), where \(C\) excludes the transition vertex. For this parity their changes cancel identically; the other parity is rapidly small by Proposition~\ref{phase:offsupport}. Telescoping over the sites gives cutoff independence. The identity is linear and applies with any \(F\).

\subsection{Uniform resolution of a higher-valence positive wall}
\label{phase:positive-walls}

The local convolution also yields a signed formula useful at a positive-centre boundary. For \(d\ge3\), set \(k=d-2\),
\[
 u_{\epsilon,q}=\sum_i\epsilon_i\theta_i+\pi q,\qquad
 B_k^{\mathrm{per}}(u)=\sum_{a\ne0}a^{-k}e^{iau},\qquad D=-i\partial_u.
\]
After dividing the left side of \eqref{phase:convolution} by \((-1)^{rq}r^{d-2}\), the resulting transfer is
\begin{equation}
 \frac{(-1)^q}{2(2i)^d}
 \sum_{\epsilon_i=\pm1}\left(\prod_i\epsilon_i\right)
       [m(D/r)B_k^{\mathrm{per}}](u_{\epsilon,q}).
 \label{phase:bernoulli}
\end{equation}
This follows by extending the even summand to all nonzero integers and expanding every sine. Off \(2\pi\mathbb Z\) the periodic distribution is the Bernoulli polynomial of degree \(k\), with convention
\(B_k^{\mathrm{per}}(2\pi t)=-(2\pi i)^k B_k(t)/k!\) for \(0<t<1\).
Near a lifted wall \(u=2\pi j+w\), it is exactly \(P_-(w)+c_k w_+^{k-1}\), where \(P_-\) is the left polynomial continued across the wall and \(c_k=2\pi i^k/(k-1)!\). Use the step distribution when \(k=1\). With \(\check m(y)=(2\pi)^{-1}\int e^{isy}m(s)\,ds\),
\begin{align}
 [m(D/r)B_k^{\mathrm{per}}](2\pi j+w)
 &=\sum_{\ell=0}^{k}\frac{m^{(\ell)}(0)}{\ell!\,r^\ell}D^\ell P_-(w)
   +c_k r^{1-k}\int_{\mathbb R}\check m(y)(rw-y)_+^{k-1}\,dy
   +O(r^{-\infty}).
 \label{phase:wall}
\end{align}
The estimate is uniform on a sufficiently small fixed wall neighbourhood, including any fixed derivatives. In fact Poisson gives the periodized kernel \(r\sum_j\check m(r(u+2\pi j))\); only the nearby polynomial and truncated power survive modulo rapid error, with their displayed exact convolutions. On either tail the scaled truncated-power convolution agrees with its respective polynomial branch to rapid error. Simultaneous walls at the same vertex occur linearly in \eqref{phase:bernoulli}. Intersections of walls from different vertices are treated by the signed cell expansion of the full radial sum, rather than by adding independent wall contributions.

\subsection{Complete phase-labelled expansions}
\label{phase:all-orders}

\begin{theorem}[Two-action radial expansion]
\label{phase:main}
For each fixed paired-cap tree and every fixed smooth joint trace test \(F\), the physical decomposition \eqref{phase:split} has
\[
 Z_r^{[p]}[F]\sim (-1)^{rp}
     \sum_{\alpha,j}r^\alpha(\log r)^j\mathcal C_{\alpha,j}^{[p]}(F),
       \qquad p=0,1,
\]
to arbitrary algebraic order. The power scale and log degrees are those of the all-tree radial expansion. There is no level-residue dependence after the displayed action phase. At every order \(\mathcal C_{\alpha,j}^{[p]}\) is a finite-order distribution in the trace coordinates, independent of the cutoff choices, supported on the radial image of the flat locus of action \(p/2\). The finer tagged coefficients have the support of Proposition~\ref{phase:offsupport}.
\end{theorem}

\begin{proof}
The all-tree lattice expansion in Theorem~\ref{rad:thm} supplies existence to every algebraic order on the even and odd subsequences. The cap selection \eqref{phase:tag} is half the standard odd-deficit magnitude on each of the two grids, with the even grid carrying \((-1)^q\). Thus no new summation walls are introduced by the phase labels: one uses the exact parity-dependent local character transform, including its even-grid omitted zero term, before the affine orthant and Mellin steps. In particular the remainder and Fourier-mode coefficient bounds there apply to these same jointly marked physical contractions. It remains to remove the subsequence dependence after \eqref{phase:sign}.

For a fixed tagged cell strip the factor \((-1)^{rp}\) and prolong the level continuously to \(R>1\) on the fixed grids of \eqref{phase:sign}, keeping all indices integral. All folded profiles have smooth support separated from their opposite endpoint \(x=1\), and have finite Laurent order at \(x=0\); no grid point at zero is used. The prolongation \(f(R)\) is smooth by locally finite sums. We establish, for a fixed small \(\varepsilon>0\) and some \(B\) independent of derivative order,
\begin{equation}
 |\partial_R^j f(R)|\le C_j R^{B-\varepsilon j}.
 \label{phase:symbol}
\end{equation}
To prove this, expand the edge sines and split every folded coordinate smoothly into small and large at \(R^{-\varepsilon}\). Hold the small labels \(n\) on their fixed grids, so their scaled coordinates \(y=n/R\) are \(O(R^{-\varepsilon})\), and Poisson-sum all the large coordinates \(x\) after smooth extension by zero near their endpoint. Alternation factors on held labels are constant in \(R\). Each remaining integral has phase
\[
 \pi R \Phi(x,y),\qquad
 \Phi=\sum_{vw\in E}\epsilon_{vw}\widetilde x_v\widetilde x_w
                         +\sum_{\mathrm{large}\ v}c_v x_v ,
 \qquad \widetilde x=(x,y),
\]
where \(c_v\in\mathbb Z\) are the remaining alternations and Poisson aliases. Shifted-grid factors are independent of \(R\). Unbounded aliases have gradient comparable to their size and give an absolutely summable rapidly small contribution by integration by parts, including after a fixed number of \(R\)-derivatives.

For the finitely many bounded aliases, use a smooth cutoff to \(\lvert\nabla_x\Phi\rvert\lesssim R^{-\varepsilon}\). Its complement is similarly rapid: each integration by parts gains \(R^{-1}\), while differentiating the reciprocal gradient and scale cutoffs costs at most \(R^{3\varepsilon}\) per step, up to a fixed initial pole factor. The large-coordinate profile derivatives have this bound because their singular endpoint is at distance at least a constant multiple of \(R^{-\varepsilon}\). Take, for instance, \(\varepsilon<1/16\). These gains dominate any fixed finite number of differentiated factors and the fixed polynomial cardinality bound on the held-label sums.

On the retained region set \(y=0\). The gradient then differs by \(O(R^{-\varepsilon})\) from zero and is affine in \(x\), with a fixed linear part for each alias. If its zero system is inconsistent the retained region is empty for large \(R\). Otherwise project to the affine zero set, moving \(x\) by \(O(R^{-\varepsilon})\). The projected coordinates have modulus strictly less than one, since all supports were separated from one. Their nonzero support is a forest-independent set: an induced leaf of a nontrivial component would give one nonzero neighbour of modulus less than one in a slope which must be integral. All edge products therefore vanish at the projected point, and the linear coefficient \(c_v\) is zero at each nonzero coordinate. The residual phase there is zero. Hence
\(\Phi(x,y)=O(R^{-\varepsilon})\) throughout the retained region.

At fixed \(x\) and fixed held \(n\), the first derivative of \(R\Phi(x,n/R)\) is
\(\Phi-\sum y_v\partial_{y_v}\Phi=O(R^{-\varepsilon})\).
Its \(j\)-th derivative for \(j\ge2\) is \(O(R^{1-j})\), since the phase is quadratic-affine. Held Laurent profiles differentiate by \(-y_v\partial_{y_v}/R\), which retains their fixed pole order. Scale cutoffs, Poisson prefactors and the near-gradient cutoff gain \(R^{-1}\) per \(R\)-derivative on their transition: the \(y\)-part of that scaled gradient is \(O(R^{-\varepsilon})\). Thus every exponential derivative gains at least \(R^{-\varepsilon}\), while the undifferentiated polynomial bound on amplitudes and held counts is independent of \(j\). This proves \eqref{phase:symbol}; if all labels are held small the initial phase products give the same bound directly.

An additional parity-dependent coefficient in an all-orders power--log expansion of \(f(r)\) would contradict \eqref{phase:symbol}. Choose its largest exponent, and then its largest nonconstant logarithmic coefficient. A sufficiently high forward difference lowers all larger constant-coefficient terms in power, but acts invertibly on the mean-zero part of the two-periodic coefficient. The resulting term at the selected order survives on a subsequence. On the other hand arbitrarily high forward differences are bounded by integrating \eqref{phase:symbol} over unit increments. This contradiction proves residue independence after the actual exterior action.

Finally, the fixed-order estimates of the lattice expansion have polynomial costs in Fourier-mode size, so only finitely many trace seminorms are required for each coefficient functional. This gives distributional continuity. Proposition~\ref{phase:offsupport} makes every such coefficient vanish on test functions with support disjoint from the corresponding projected flat set. The finite partition and the cutoff argument above give the stated support and cutoff independence.
\end{proof}

A fibre-only partition is insufficient even for the two actions at some radii. For example, on the path \(P_4\) a positive centre at site 2 and negative end at site 4 admit noncentral sites 1 and 3 with either \(h_3=h_1^{-1}\) or \(h_3=-h_1^{-1}\). At trace-zero \(h_1,h_3\) these have identical fibre radii but opposite values of \(q_2\) and of the action. The bounded cap-product test separates them. Even these tags do not separate all same-action angular configurations; all-negative fibres have trivial radial and cap-trace variation regardless of their relative axes.

\subsection{Absolute adjoint eta and the sewing phase}
\label{phase:eta}

Use the positive mixed odd signature operator
\(\mathcal D_A(\phi,a)=(d_A^*a,d_A\phi+*d_Aa)\).
For every flat character the adjoint relative invariant in this convention is
\begin{equation}
 \rho_{\rm ad}(M_G,A)=
       \eta(\mathcal D_A)-3\eta(\mathcal D_{\rm triv})=0 ,
 \label{phase:rho}
\end{equation}
where zero eigenvalues are omitted. This is a real equality, including reducible and singular strata. The following computation also fixes the sign of the seam correction independently of the finite colour phase.

Start with rational noncentral eigenphases and choose positive primitive filling slopes \(a_e z+b_e h\) at one end of each torus, with coefficients exchanged at the other and both fundamental meridian holonomies trivial. For two negative centres use \((1,1)\); for a positive--negative edge take \((2,1)\) at the positive end; for two positive centres use any positive primitive pair. If the noncentral eigenphase is \(m/(2n)\) in units \(2\pi\), coprime \(m,n\), take \((a,b)=(1,n)\) for odd \(m\) and \((2,n)\) for even \(m\) against a negative centre. Against a positive centre take \((1,2n)\). These choices ensure trivial holonomy at both swapped ends. Let \(K_v=\sum_{e\ni v} b_e/a_e>0\) in its own end coordinates. The single-site case has no seams and gives \(S^2\times S^1\) directly.

Let \(I_v\) be the invariant adjoint subspace of the entire vertex representation and \(W_e\) the peripheral invariant subspace, of dimension three on central--central seams and one otherwise. Write \(i_v=\dim I_v\), \(w_e=\dim W_e\). In separately transported frames let \(J_{v,e}:I_v\to W_e\) be the inclusion and \(T_e\) the orthogonal transport between the two edge frames. In section/fibre coordinates the kernel of the boundary-to-block homology map is exactly
\begin{equation}
 x_{v,e}=J_{v,e}c_v,\qquad
 \sum_{e\ni v}J_{v,e}^* y_{v,e}=0 .
 \label{phase:boundary-kernel}
\end{equation}
On the invariant summand this is the ordinary section sum and common fibre, since the two cap Euler corrections cancel. If the fibre is noncentral its transverse system is acyclic. If the fibre is central, on the complement of \(I_v\) the product fibre classes vanish by \(H_0=0\), and the boundary section classes of the punctured base jointly inject by relative duality \(H_2(\Sigma,\partial\Sigma)=0\) with those coefficients. The two genuine cap fillings kill exactly their own section classes; the cap degree-zero maps are isomorphisms on their fixed subspaces. This gives \eqref{phase:boundary-kernel} without assumptions on neighbouring cap axes.

To evaluate the triple Maslov indices in the filling comparison, put \(\beta=b/a\) at the first end of an edge and
\(z_{v,e}=y_{v,e}-\beta x_{v,e}\), \(q_e=z_{v,e}/\sqrt\beta\). The seam kernel is the anti-graph of the marked swap, not its graph. With a simultaneous swap sign \(\epsilon_e\), the other end has \(z_{w,e}=\epsilon_e\beta^{-1}T_e z_{v,e}\). Define the incidence map \(A:\bigoplus W_e\to\bigoplus I_v\) by its two endpoint columns \(\sqrt\beta J_{v,e}^*\) and \(\epsilon_e \beta^{-1/2}J_{w,e}^*T_e\). Put \(K_I=\bigoplus K_v I_{I_v}\) and \(Q_I=K_I-AA^*\); the ordinary version \(Q\) is the negative signed backbone adjacency. Decomposing a seam-kernel vector into a filling-meridian part and \eqref{phase:boundary-kernel} gives \(K_I c=-Aq\). With the outward intersection convention \(\omega(z,h)=1\), their skew pairing is
\(\langle Aq,K_I^{-1}Aq\rangle-\langle q,q\rangle\).
Indeed the endpoint terms in \(\sum x z\) on each seam add to \(-\lVert q_e\rVert^2\), and the remaining part comes from \(c\). Two Schur complements of the block matrix with diagonal entries \(K_I,I\) and off-diagonal \(A,A^*\) give
\[
 \tau_{\rm ad}=\sum_v i_v-\sum_e w_e-\operatorname{sig}Q_I,\qquad
 \tau_{\rm ord}=|V|-|E|-\operatorname{sig}Q .
\]
No nonsingularity of the last two forms is required.

Apply the signature gluing formula \cite[Theorem 3.9]{Toffoli} to the blocks and seam cylinders with these meridian fillings. The auxiliary seam closures have equal meridians on both ends and are \(S^2\times S^1\) with flat extension over \(D^3\times S^1\); they contribute zero. At a central site compare the tilted fillings with fibre fillings \(h\). The fibre-filled comparison manifold is a connected sum of the two order-two cap lens pieces and copies of \(S^2\times S^1\): filling a boundary fibre splits the genus-zero base into the individual cap balls and ordinary handles. The cap relative invariants vanish, since the adjoint is trivial at a positive fibre and the order-two cotangent sum is zero at a negative fibre. The comparison Maslov form on the section vector \(c\) in \eqref{phase:boundary-kernel} is \(-K_v\lVert c\rVert^2\): decompose \((Jc,\beta Jc)\) into its fibre-filling part and section part. Thus the filled central contribution is \(3-i_v\) plus the comparison lenses.

At a noncentral site the axis is trivial and the two transverse unitary lines are conjugate. Compare against section fillings \(z\) for these adjoint lines: the original sections are central in the fundamental system and hence trivial in the adjoint. The resulting bare two-cap Euler-zero piece is \(S^2\times S^1\); its transverse line extends over \(D^3\times S^1\). The twisted boundary tori of each line are acyclic, whereas the ordinary comparison form is
\(\sum_e a_e b_e q_e^2\) subject to \(\sum_e b_e q_e=0\). This follows by writing the filling vector as \((a_e q_e,b_e q_e)\) and removing its common section part. Its signature is \(d_v-1\). The two-line contribution is consequently \(-2(d_v-1)\) plus its comparison lenses. A comparison lens for a mixed seam is orientation-reversed across the swap: the central pair of meridians \((b z+a h,h)\) matches the noncentral pair \((a z+b h,z)\). The surviving line character is transported or inverted, which gives the same relative invariant on the two conjugate lines. Thus all such lenses cancel, and comparison lenses on central--central seams have trivial adjoint coefficients.

The remaining correction is \(3\operatorname{sig}Q-\operatorname{sig}Q_I\): substitute the local factors and \(\sum_e w_e=3|E|-2\sum_{v\in S}d_v\) into the two Maslov terms. Both signatures vanish by bipartite sign symmetry, also in the presence of null directions. This proves \eqref{phase:rho} at rational eigenphases.

To pass to general eigenphases one can rationalize within constant local cohomology ranks. At a central tuple with a common invariant line the fixed axis lifts give affine rational angle equations; rational points are dense in their solution space. At a locally noncommuting central tuple the product map of the fixed input classes is a submersion: the annihilator of its differential is the intersection of the successive transported invariant axes, equal by induction over the factors to their common axis intersection. Thus the implicit-function theorem adjusts the axes when the input angles vary to rational values. Other local ranks and signed axis transports may be held fixed, since the gluing graph is a tree. The vertex \(H_1\) dimensions are \(3d_v+2i_v-c_{v,+}-c_{v,-}\) at central sites, where \(c_{v,\pm}\) are the fixed dimensions of the two genuine cap holonomies, and \(d_v\) on a noncentral site. These follow from the base Euler characteristic and the same cap section eliminations as \eqref{phase:boundary-kernel}. The kernel of the degree-one seam map is \(\ker Q_I\) by that equation and the anti-swap gluing; the degree-zero seam rank is that of the fixed transported inclusion-difference map. Hence the total \(h^0+h^1\), which is the odd-signature nullity, is constant on each local rationalizing path. Eta is continuous there, since the small eigencluster does not gain or lose nonzero modes. This proves \eqref{phase:rho} everywhere. The argument uses unreduced eta and makes no identification of the absolute determinants of different strata.

\subsection{Regular positive covers and finite termination}
\label{phase:regular-cover}

There is a regular higher-valence regime in which one can compute the entire phase series without a singular lattice expansion. Let \(G=P\sqcup S\) be a bipartition and write \(p_0=|P|\), \(s=|S|\), \(N=p_0+s\), \(E=N-1\). Keep the \(S\)-radii in compact noncentral windows using any smooth joint \(F(x_S)\); mark each \(P\)-radius near \(+I\), with its cutoff equal to one near that centre, and fix its product \(q_v\). Assume throughout the compact angle support that the \(d_v\)-fold product of incoming classes at every \(v\in P\) has a nonempty regular fibre over \(z_v=(-1)^{q_v}I\). Its Haar probability density \(\rho_v\) is then smooth in a common neighbourhood of the target. This is a genuine higher-valence region; a nonempty regular product at the centre requires at least three noncentral inputs. Set
\[
 \mathcal J_q(x_S)=
   \prod_{j\in S}(1-\cos\theta_j)\prod_{v\in P}\rho_v(z_v),
 \qquad D_0=3s-3p_0-3.
\]
The exact leading coefficient in \(Z_r\) is
\begin{equation}
 (-1)^{r\sum_Pq_v}
   2^{(1-N)/2-p_0}\pi^{\,2p_0-E}r^{D_0/2}
       \int F(x_S)\mathcal J_q(x_S)\,d^s x .
 \label{phase:regular-leading}
\end{equation}
After the displayed phase and power have been extracted, the full expansion is a finite polynomial in \(r^{-2}\), modulo \(O(r^{-\infty})\), of degree at most
\(\sum_{v\in P}\lfloor(d_v-2)/2\rfloor\). In particular in the trivalent positive-cover regime the quadratic coefficient is the only algebraic term under these marks.

To prove the quantum assertions note that \(P\) is independent, so its conditional colour sums \eqref{phase:convolution} factor with the \(S\)-colours held fixed. At a regular target the local multiplier expansion is
\[
 (K_{r,m_v}*\mu_v)(z_v)\sim
       \sum_{\ell\ge0}\frac{m_v^{(2\ell)}(0)}{(2\ell)!\,r^{2\ell}}
                                   (1-\Delta)^\ell\rho_v(z_v).
\]
Indeed the rest of the class measure is a fixed distance away and is rapidly small. A smooth extension of the local density has rapidly decreasing character coefficients, uniformly in the tested angles, and
\(-\Delta\chi_a^{SU(2)}=(a^2-1)\chi_a^{SU(2)}\) in the unit quaternion metric. All signed-angle walls at the central target are excluded by regularity, since a hit would realize a collinear tuple. Hence \eqref{phase:bernoulli} acts on polynomials of degree at most \(d_v-2\), and the even multiplier terminates at the stated order, with a rapid remainder. In the remaining compact smooth \(S\)-sum the ordinary step-two Poisson density is \(r/2\) per site with no nonprincipal aliases to any algebraic order. The \(\sin\theta_j\) factors from \eqref{phase:convolution} cancel the sine denominators at \(S\). The common coefficient is consequently the product of
\[
 (2/r)^{(N+1)/2},\quad 2^{-p_0},\quad (r/2)^s,\quad
       \prod_{v\in P}\bigl(r^{d_v-2}\pi^{2-d_v}/2\bigr),
\]
where \(m_v(0)=\pi^{2-d_v}/2\) and \(\sum_{v\in P}d_v=E\). This is \eqref{phase:regular-leading}.

The absolute quadratic interpretation can be checked independently. In primitive section/fibre periods the noncentral block at \(j\in S\) has adjoint half-torsion
\((1-\cos\theta_j)/(2\sin\theta_j)^{d_j}\): the perpendicular fibre determinant is \(4\sin^2\theta_j\), each cap core contributes \(4\sin^2(\theta_j/2)\), and the two primitive cap slopes give longitudinal half-torsion \(1/2\). At a positive regular block there are no common adjoint invariants. Its two true cap cores are adjoint-trivial. In base trace coordinates with the paired fibre-area copies the cap-free torsion cancels degree by degree; each cap has degree-one determinant \(2^3\) and unit remaining pivots. The local half-torsion is therefore \(1/8\). Each seam retains just its axial line. Differences about a noncentral star, one-sided lifts and the complementary sum-zero section periods have unit relative determinant. The oriented torus cup pairs the common radial period with its excluded error, and the fixed-class surface tangent with its fibre-area copy by the Goldman form. Thus there is a radial factor \(\pi\,dx_j\) per \(S\), and \(E-s\) bend circles of period \(2\pi\). The coarse fixed-class volume at \(v\in P\) is
\[
 \frac{2}{(2\pi^2)^2}\prod_{j\sim v}(2\pi\sin\theta_j)\rho_v(z_v);
\]
this is the regular class-convolution reduction formula of \cite[Section 5.1]{AMW}, in the unit quaternion metric. The identity of the seam determinants can also be read by taking relative periods \(t_i-t_{d_j}\); a one-sided lift complements each basis \(e_i-e_{d_j}\) with determinant of modulus one.

There are \(2^{s+p_0}\) genuine cap-root sheets. The independent positive-block centres divide the bend volume by \(2^{p_0-1}\), because signs acting trivially through each noncentral star must coincide on the whole tree. Use a further \(1/2\) for the single global centre in the stack convention. The resulting stack half-torsion pushed to the radii is
\[
            2^{E-4p_0}\pi^{2E-4p_0}\mathcal J_q(x_S)\,d^s x.
\]
Here the dimensions are \(s\) radial, \(E-s\) bend, and \(\sum_{v\in P}(2d_v-6)\) fixed-class surface directions, adding to \(D_0\). The invariant quadratic factor \((r/(2\pi^2))^{D_0/2}\) in these metric units \cite{Adams}, with the action \eqref{phase:cs} and zero adjoint relative eta \eqref{phase:rho}, gives precisely the absolute coefficient and phase in \eqref{phase:regular-leading}. Finite termination is proved for the indicated marked quantum responses; it is not a statement about each angular or higher-loop coefficient of an unmarked path integral.

\section{The actual classical Hall action and its normalization}
\label{classic:sec:action}

We use the trace pairing \((a,b)=-\operatorname{Tr}(ab)/2\) on
\(\mathfrak{su}(2)\), so that unit imaginary quaternions have norm one,
\(\operatorname{Vol}SU(2)=2\pi^2\), and the metric effective group has
\(\operatorname{Vol}SO(3)=\pi^2\).
The exponent convention is \(e^{2\pi i r\,\mathrm{CS}}\), where \(r=k+2\).
On the flat actions \(0,\frac12\), replacing \(k\) by \(r\) leaves the
phase unchanged; in the leading normal integrals the level shift changes
neither the scale nor the leading coefficient. All classical integrations
below are finite residual integrals: the nondegenerate quadratic complement
is integrated invariantly, the finite zero modes are retained, and no
identification with a nonperturbative functional integral is assumed.

Let \(S\) be the genuinely noncentral fibres in an all-negative-centre
Hall chart, \(T=N(S)\), and let \(I\) be the induced forest on the other
vertices. Here the Hall hypothesis on the chart is that \(I\) is balanced
with a perfect matching (and \(S\) has the prescribed noncentral data).
The case \(S=T=\varnothing\) with a perfectly matched tree is included.
Write \(\alpha\) for the active parameters. For
\(v\in I\), let \(u_v\) be its cap axis on the central plane,
\(h_v=-\exp(\pi y_v u_v)\) to first order, and
\(t_{vw}=\epsilon_{vw}u_v\cdot u_w\) be a signed physical dot in
the marked seam frames. Put
\[
 B_I(u)_{vw}=
 \begin{cases}t_{vw}&vw\in E(I),\\0&\text{otherwise}.\end{cases}
\]
Signs can equivalently be moved to the axes since the graph is a forest.
All axis dots below refer to the common framed configuration, not to
independently chosen frames for different tests.

\begin{theorem}[All-angle cyclic Hall normal form]
\label{classic:thm:hall-action}
At every globally irreducible flat point in such a chart, the actual
finite harmonic Chern--Simons reduction is stably real-analytically
right equivalent to
\begin{equation}
 \mathrm{CS}_0+\frac{c}{2}y^{\mathsf T}B_I(u)y,
 \qquad |c|=\tfrac12 .
 \label{classic:eq:normal}
\end{equation}
Stability adds or removes only the finite gapped quadratic radial
pivots of this graph form. The equivalence fixes the entire central
section \(y=0\) pointwise in its literal angular and old active
parameters. It is jointly analytic through central active product
locks, including locally collinear locks, and through parallel as well
as perpendicular dormant angles. A choice of orientation fixes the
common sign convention for \(c\). The theorem does not identify all
off-shell extensions of the period marks under the normal change.
\end{theorem}

\subsection{Analytic representation scheme in peripheral marks}
\label{classic:sec:scheme}
We first describe the true flat ideal, before drawing any conclusion about
a potential. Use the graph-of-groups relations
\[
 [h_v,h_w]=1,\qquad A_v^2=h_v^{-1},\qquad
 {\rm Ad}_{A_v}P_v=P_v^{-1},
\]
with \(P_v\) the actual ordered product of transported neighboring fibres
(with their seam signs). The other cap root is
\(A_v^{-1}P_v^{-1}\). Local logarithms are taken after dividing off a
fixed central sign. All changes below use analytic units or invertible
analytic matrices and therefore hold in the analytic relation ideal,
not just on real points.

Write \(h_t=-\exp(p_tu_t)\) at an active centre \(t\). If an old active
neighbor of \(t\) is not parallel to \(u_t\), the corresponding
commutator kills \(p_t\) with a unit coefficient. At a remaining
\(t\), every incoming active axis is parallel to \(u_t\), and the
negative-cap equation at the base point implies \(P_t\) is central.
Thus each remaining column \(t\) touches at least two active
neighbors. An active \(s\) has \({\rm Ad}_{A_s}+1\) invertible; its
reversal equation kills \(\log(\pm P_s)\). Its longitudinal projection
has derivative a nonzero real entry on each incident surviving column
\(t\). This rectangular forest matrix has independent columns:
the graph induced by any nonempty supporting columns, with all their
neighboring rows, is a forest with every column degree at least two
and therefore has a row leaf. That row rules out its sole supported
column. Select a full-rank minor, including the earlier commutator
pivots. Every selected equation vanishes identically when \(p_T=0\).
Analytic division by that coordinate ideal and the invertible minor
therefore give the exact ideal \((p_T)\).

Once \(p_T=0\), all constraints at \(S\) are automatic. If \(P_t\)
is noncentral the scalar \(u_t\cdot\operatorname{Im}P_t=0\) solves
one angular coordinate. If \(P_t\) is central at the reference, use
the \emph{whole ordered product} as a local group parameter. To do
this simultaneously, root the product-incidence forest at old active
vertices. At each locked centre reserve one child active vertex;
its group element is solved from \(L h^\epsilon R=Q_t\) by actual
multiplication, including the dormant factors in their original
positions. Each reserved variable is used once. Hence all new
\(Q_t\) are independent analytic parameters. The locked constraint
is \(u_t\cdot\log(\pm Q_t)=0\); its two transverse product coordinates
are free. The lock is their joint zero, not an additional critical
factor. The result is an analytic product with the dormant critical
scheme and smooth passive \(\alpha\), even when several locks coincide.
A global effective rotation slice may be fixed at irreducibles.

For the all-negative dormant forest the nonparallel edge commutators
give \(q_e=p_vp_w\) in the ideal. The axial equation in the full
angular domain is
\[
 F_v=\sum_{w\sim v}b_{vw}(p,u)p_w,\qquad
 b_{vw}(0,u)=t_{vw}.
\]
Indeed, after removal of the fixed signs the product is an analytic
product of \(\exp(\pm p_wu_w)\). At a negative cap the longitudinal
projection of the reversal equation is identically a unit times its
axial logarithm. Multiply \(F_v\) by \(p_v\), remove all already known
\(q_e\), and peel the forest of edges parallel at the reference:
each leaf row has a single remaining coefficient with value \(\pm1\).
Thus all \(q_e\) belong to the ideal. Conversely the transverse
reversal equations vanish modulo \((q_e,F_v)\); there \(p_v\log P_v=0\),
and the transverse part is the pure half-turn equation. It follows
that \((q_e,F_v)\) is exactly the remaining ideal.

For use through parallel angles, write the row coefficients as
\(b_{vw}=t_{vw}+c_{vw}\), with \(c_{vw}\) in the ideal of \(v\)'s
neighboring periods, by radial integration of the ordered product
expansion. In \(\sum_w c_{wv}p_w\) each monomial is divisible by an
edge product \(p_wp_j\). Therefore
\(\widetilde t_{vw}=t_{vw}+c_{vw}+c_{wv}\) is a symmetric coefficient
whose radial rows agree with \(F_v\) modulo \((q_e)\). Differentiating
\(\sum_e\widetilde t_e q_e\) shows that its Jacobian is contained in
\((q_e,F_v)\). To reverse the inclusion use independent child-subtree
rotations on nonparallel reference edges, affecting only that edge's
literal dot, and add the rows \(p_v\partial_{p_v}\) on the parallel
forest. These rows generate all \(q_e\) by an invertible coefficient
minor: at zero periods the angular block is diagonal with nonzero
entries on the nonparallel edges, and the parallel subblock is the
same leaf-eliminable incidence matrix. The ordinary radial rows then
recover \(F_v\). Finally on parallel edges solve recursively
\(d_vd_w=\widetilde t_{vw}/t_{vw}\) with positive analytic units
\(d_v=1\) on \(p=0\). Set \(p'_v=d_vp_v\). On each nonparallel edge
use its independent relative subtree rotation to realize
\(\widetilde t_e/(d_vd_w)\) as the corrected literal dot. The
construction is invertible and fixes the central angular plane
pointwise. Thus the true dormant ideal in these ambient coordinates
is the Jacobian of
\[
 \frac12 p^{\mathsf T}B_I(u)p=\sum_e t_e p_vp_w .
\]
This is a graph-scheme statement only.

For clarity, the identification with the analytic Coulomb critical
scheme used here is the general holonomy correspondence on finite
local Artin rings \cite[\S\S5--6]{GoldmanMillson}, not a quadratic
formality theorem for this three-manifold. Fix an irreducible
effective slice on both sides. On a small Coulomb slice the
off-harmonic gradient equation has an invertible elliptic
linearization and is eliminated by the analytic implicit function
theorem (use Sobolev degree high enough for multiplication). The
remaining gradient represents the flat functor. The holonomy
comparison and the graph pivots identify these local quotient
analytic rings. In the ambient comparison, first split the radial
nondegenerate block of the graph Hessian at the reference. The
minimal rings then have the same tangent dimensions; lifts of
their analytic coordinate functions give inverse-function
isomorphisms of the ambient germs. They may be chosen to retain
the literal coordinate trace on the smooth common central section,
by correcting lifts by their trace there. Restore the finite
nondegenerate quadratic block on the Coulomb side. This gives a
stable coordinate comparison with the same Jacobian ideal and the
named central section, without claiming that equality of ideals
alone determines the Chern--Simons value away from it.

\subsection{Cyclic comparison of the actual radial Hessian}
\label{classic:sec:cyclic}
The extra input is a \emph{family} of symmetric radial differential
comparisons; comparing torsion numbers or signatures pointwise would
not suffice. First freeze the literal old active tuple \(\alpha_0\)
and let the dormant axes vary. Active pieces see only fixed central
dormant boundary holonomies on this plane, while each dormant piece
has a constant flat local system up to analytic rotation of its
cap axis. Seam changes of frame are their actual rotations. If the
active tuple itself is collinear, use a dormant nonparallel axis
for the residual azimuth; global irreducibility guarantees such an
axis. Thus freezing is a restriction on actual holonomies, not a
coordinate freeze that changes active holonomies when the dormant
axes move.

Use the twisted de Rham mapping cone of thickened pieces and collars.
There are no triple overlaps. The ordinary restriction map to this
cone is a chain map and induces the closed cohomology isomorphism
by the partition-of-unity proof of Mayer--Vietoris. Locally
analytic flat frames on a fixed cover can be patched with a smooth
partition of unity to identify nearby coefficient bundles; the
patched linear maps remain invertible near the base parameter and
polar decomposition preserves the inner product. Contract the
fixed-rank piece and seam complexes using analytic Hodge (or cellular)
contractions within this frozen family. If \(\delta\) is the
vertex-to-edge restriction difference and \(h\) is the block-diagonal
contraction, \(\delta h\delta=0\) because \(\delta\) kills the edge
summand. Thus the transferred differential is exactly the
cohomological restriction cone, with no higher-transfer terms.

Here is the precise radial part of its primitive elimination.
Use the common fibre class at each piece in its invariant space,
section periods \(c_{ve}\) with their primitive constraint
\(\sum_e J_{ve}^{*}c_{ve}=0\), and complementary one-sided area
lifts, where \(J_{ve}\) is the inclusion of the local invariant
space into the seam invariant space in the seam frame. Noncentral
torus complements are acyclic. At a central piece the
non-invariant part has independent restrictions on its central
exterior section planes, by the relative cohomology exact
sequence and the absence of invariant vectors there. Hence those
planes have unit elimination pivots. The swap identifies each
half-edge section with its neighboring fibre period. Substitution
in the section sum leaves the compressed signed adjacency.
In degree zero the inclusion difference on invariant vectors
and in degree two its adjoint are eliminated as injective and
surjective maps respectively, since the closed reference is
irreducible. They do not change the radial block. The connecting
cup formula is explicit in a collar: take a normal bump of
integral one and cup the associated section error with the
original fibre form. A radial error consequently pairs with its
own unit fibre period; a tangent with those dormant periods
identically zero pairs to zero. The ordinary determinant version
is also used in Theorem~\ref{thm:hall:torsion}.

If every \(T\) piece has no invariants, the radial block is
\(B_I\) plus the passive old active directions. At a locally
collinear lock its invariant line must instead be retained. Let
\(T_1\) consist of those lines in the frozen family. In the
\(S,T_1,I\) order the radial block is exactly
\begin{equation}
 \begin{pmatrix}0&U^{\mathsf T}&0\\
                 U&C&V\\
                 0&V^{\mathsf T}&B_I
 \end{pmatrix}.
 \label{classic:eq:activeblock}
\end{equation}
The entries of \(U\) on the incidence forest are signed unit
axis pairings; every \(T_1\) row has at least two old active
neighbors. The leaf proof above again shows \(U^{\mathsf T}\)
is injective. Split \(S=S'\oplus\ker U\) so the restriction
\(U'\) to \(S'\) is square invertible. The inverse of the
\(S'\oplus T_1\) block is
\[
 \begin{pmatrix}
 -(U')^{-1}C((U')^{\mathsf T})^{-1}&(U')^{-1}\\
 ((U')^{\mathsf T})^{-1}&0
 \end{pmatrix}.
\]
In particular it has \emph{zero hub--hub block}. The dormant
coupling reaches only that block, so the Schur correction is
identically zero and the retained graph is \(B_I\) itself.
A dormant tangent may acquire an \(S'\) correction, but its
cup pairing with a dormant connecting error is still the unit
dormant period. This is why local collinear locks do not induce
a new dormant spring or rescale its radial units.

On the true closed Hodge side, an analytic sufficiently small
positive gap for the twisted Laplacian around the irreducible
reference isolates degrees one and two, with star identifying
the two low spaces. The analytic Hodge projection away from
exact forms uses the nonzero scalar gap. The tangent to the
central family gives analytic harmonic modes in the kernel of
the corresponding coexact curl; split these angular modes,
using the star-dual split in degree two. Similarly split the
already nondegenerate radial block of \(B_I\) at the reference
by ordinary analytic symmetric Gaussian elimination. The closed
low complex and the model cone now give analytic chain maps
\(T_1',Q_2'\) on their remaining radial blocks:
\[
       B\,T_1'=Q_2'D ,
\]
where \(B\) denotes the graph Schur block and \(D\) the low
symmetric differential in physical units after the fixed
conversion described below. The maps on radial blocks are
invertible at the reference and hence nearby. At the base
point the cone map is a cup isomorphism; in the splits just
described the degree-one map sends angular tangents to angular
tangents, which pair to zero with radial degree-two connecting
errors. The induced quotient pairing of the radial tangents
and radial connecting errors is the unit one by the primitive
torus computation. Thus
\((T_1')^{\mathsf T}Q_2'=1\) at the reference. No pointwise
minimal basis is chosen as a function of the dormant rank.

To compare with the \emph{nonlinear} harmonic reduction's Hessian,
project from the static to the moving Coulomb slice along exact
directions; on a flat connection the Hessian kills those
directions. The static transverse high complement is a small
analytic graph \(J\) over the moving high spectral complement.
Writing the moving forms as \(D_L,D_H\), its static Schur
Hessian is
\[
 D_L-D_LJ(D_H+J^{\mathsf T}D_LJ)^{-1}J^{\mathsf T}D_L.
\]
The gapped inverse is analytic in the elliptic Sobolev spaces
and the correction acts through a finite-dimensional space.
This form is \(A D_L\) with \(A\) analytic, equal to the
identity at the reference and symmetric after multiplication
by \(D_L\). Accordingly the analytic square root gives an
actual congruence. The same square-root argument applies to
the cone comparison: \(P=(T_1')^{\mathsf T}Q_2'\) is near
one and \(PD=(T_1')^{\mathsf T}BT_1'\) is symmetric, so its
analytic principal root obeys
\(P^{1/2}D=D(P^{1/2})^{\mathsf T}\), even at rank changes.
The Schur complement comparison extends back to the
stabilized radial block by using the explicit nondegenerate
pivot. Finally the primitive fibre period in \(y\) is
\(\pi\,dy\), its degree-two partner is cup-dual, and the CS
Hessian is \(\pm (2\pi^2)^{-1}\) times the coexact curl
pairing, obtained by differentiating
\((8\pi^2)^{-1}\int\operatorname{Tr}(A\,dA+2A^3/3)\).
The symmetric Hessian in signed Wilson units is thus
\(\pm\tfrac12 B_I\), independent of any WRT coefficient.

It remains to extend the frozen comparison analytically in
\(\alpha\), as the cut-piece cohomology need not remain
constant. Write \(F(y,u,\alpha)\) for the actual stabilized
germ in the common Jacobian chart, with its constant flat value
subtracted and \(F_2=\tfrac12 y^{\mathsf T}D(\alpha,u)y\).
Let \(G=\tfrac{c}{2}y^{\mathsf T}B_I(u)y\). The shared
Jacobian ideal is generated in radial degrees one and two
by the linear rows \(B_Iy\) and all edge products. At
\(\alpha_0\) use the analytic radial congruence just obtained.
It preserves the linear ideal by the degree-one equality.
It also preserves the degree-two edge module modulo the
linear ideal: \(\operatorname{Jac}F_2\subseteq
\operatorname{Jac}F\) by radial homogeneity in the common
graded ideal. The induced invertible symmetric-square map
on quadratic polynomials therefore sends that module into
itself. Its inverse is a polynomial in the map, by
Cayley--Hamilton with unit constant term, so the inclusion
reverses. This argument does not require the edge module to
be locally free. Apply this one \(u\)-dependent coordinate
change to all active parameters. Now
\(D(\alpha,u)=A(\alpha,u)cB_I(u)\) for analytic \(A\), by
the linear ideal equality, and \(D(\alpha_0,u)=cB_I(u)\)
for every \(u\) in the chart. A forest with a perfect matching
has determinant \((-1)^d\prod_{e\in M}t_e^2\) (only the
matched two-cycles contribute). It is invertible at a dense
set of literal angle configurations, including in a chart
through a perpendicular or parallel reference. Hence
\(A(\alpha_0,u)=1\) identically and its principal square root
gives the analytic congruence jointly in \((\alpha,u)\).
This avoids passing a cut-piece contraction through a jump.

Lastly the radial-quadratic truncation has the same full
Jacobian by the preceding degree-two argument. Use the
straight homotopy \(F_s=G+s(F-G)\) after the congruence.
The gradient differences express in the homogeneous graph
generators with coefficient matrix whose value at \(y=0\)
has only the radial-row/angular-column block possibly
nonzero. Indeed the radial differences start at degree two,
angular differences at degree three, and each edge product
is a combination of angular rows and \(y\) times radial
rows, as in the nonparallel/parallel pivots above. Thus
the coefficient matrix for the gradients of \(F_s\) is
uniformly invertible near \(y=0\), since its constant
correction is strictly block triangular. Radial Euler
integration of \(F-G\) supplies coefficients of order
\(O(y^2)\) on radial gradient rows and \(O(y)\) on angular
gradient rows. Applying the inverse preserves these orders.
The analytic Moser vector field solving
\(dF_s(X_s)=-(F-G)\) therefore has \(X_y=O(y^2)\) and
\(X_{u,\alpha}=O(y)\). Its local flow fixes the entire
central section; this proves Theorem~\ref{classic:thm:hall-action}.
The preliminary radial congruence need not be the identity
on all off-shell period directions.

The ideal elimination at old active centres also works
when some of them are positive. Use the three-vector
period at a positive vertex; each genuinely noncentral
neighbor kills the two transverse components, and any
remaining axial column has at least two old active
neighbors (its product must be central). The longitudinal
forest minor again eliminates them all, and the locked
positive product target is fixed in all three components.
This is only a separation of the framed critical scheme.
A positive \emph{dormant} component has vector periods and
is not represented by \eqref{classic:eq:normal}.

\subsection{Absolute density and what the leading angular law measures}
\label{classic:sec:leading}
Put \(a=|T|\), \(b=|S|\), \(2d=|I|\), \(N=a+b+2d\)
and \(\nu=a+d\) on the intrinsic Hall family. Let
\[
 p_G=\frac{3N-2\nu-3}{2},\qquad
 C_G=2^{(3N+1)/2-2b-3d}\pi^{1-N}.
\]
Write \(dM_*\) for the forced radial/ordered-product law of
\eqref{eq:hall:weight}, and \(dP\) for its accompanying common
normalized angular law. In particular every dormant seam dot
has conditional marginal \(dt/2\) before test weights are inserted:
a successive relative rotation along a tree produces fresh spherical
probability independently of the previous seam dots. The other
angles are \emph{conditioned jointly}, not recreated separately
for each insertion. Let \(dP_\perp\) denote this law with the
matched edge dots set to zero by conditional probability.

Contract the unique dormant matching and write an unmatched
edge \(a_i b_j\) as \(i\to j\), for a fixed bipartition within
each component. The order polytope \(\Omega\) has one
coordinate per matched pair, inequalities \(u_i\le u_j\)
along these arrows, and \(0\le u_i\le1\).
In the classical Fourier proof one can use
\(u_i=1+\log|y_{a_i}|/\log r\); the use of
\(-\log|y_{a_i}|/\log r\) instead reverses the inequalities
and has identical volume.

\begin{theorem}[The finite canonically normalized Hall coefficient]
\label{classic:thm:hall-stationary}
On the irreducible all-negative Hall charts of
Theorem~\ref{classic:thm:hall-action}, the leading angular
density of the finite CS residual model, after removing the
phase \((-1)^{r(N-1)}\) and \(r^{p_G}(\log r)^d\), is
\begin{equation}
 C_G\operatorname{Vol}(\Omega)\,dM_*\,dP_\perp .
 \label{classic:eq:hallleading}
\end{equation}
This is valid for fixed smooth tests on compact chart
localizations; it patches on compact Hall interiors. If
\(d>0\), a fixed smooth test whose trace on the indicated
matched angular support vanishes is \(O(r^{p_G}(\log r)^{d-1})\)
on a fixed compact localization. Exhausting the resolved
leading measure in the forced variables gives the entire
vacuum coefficient in Theorem~\ref{thm:hall:main}.
Equality of those totals does not assign every other classical
stratum a zero contribution.
\end{theorem}

\begin{proof}
The invariant Gaussian factor for \(h\) retained irreducible
zero modes, in cohomology frames of the stated metric, is
\(\tfrac12(r/(2\pi^2))^{h/2}\sqrt{T_{\rm adj}}\); the factor
\(1/2\) divides by the ineffective central stabilizer, not
by a continuous rotation volume a second time.
This is the quadratic integration and framing normalization
of \cite[eqs.~(3.3), (3.6)--(3.7), (3.14), (4.5)]{Adams}.
Use the analytic spectral gap before separating small eigenvalues:
its complement determinant is smooth as flat parameters vary.
If \(L\) is the Hessian of the finite reduced action on the
normal quotient, the corresponding small coexact curl product
has square root
\((2\pi^2)^{\dim L/2}\sqrt{|\det L|}/\sqrt{\det G_N}\),
where \(G_N\) is the normal quotient metric. One way to see
the precise relative statement near degeneracy is the
triangular Schur graph of the high block, followed by the
invertible low spectral projection; it intertwines the
bilinear forms up to factors tending to identity, rather
than by an additive estimate of individual eigenvalues.
The coordinate volume splits off exactly
\(\sqrt{\det G_N}\) by the Schur identity. Thus the normal
metric cancels out of the determinant-line restriction.
This gives the usual stationary determinant law here
\emph{also} as one approaches the rank-changing points
through the regular chart.

The independent regular coarea computation in
Theorem~\ref{thm:hall:torsion} gives the coefficient with
\(r^{p_G}\) removed,
\[
 D_G\,\frac{dM_*\,dP}{\prod_{e\in M}|t_e|},
 \qquad
 D_G=(2\pi^2)^{-p_G}\,2^{a+N-1}\pi^{2b+2d-2}.
\]
The stabilized graph's radial Gaussian at generic angles has
factor \(2^d/(r^d\prod_M |t_e|)\) because \(|c|=1/2\),
\(|\det B_I|^{1/2}=\prod_M|t_e|\), and its signature is zero.
Hence the smooth stabilized amplitude at zero radii is
\(2^{-d}D_G\,dM_*\,dP\), with common \(r^{p_G+d}\).
It agrees first on an open dense angular set by the
determinant comparison and extends therefrom by smoothness.
For all negative fibres/centres the CS action is
\((N-1)/2\bmod1\) and the relative adjoint eta phase is zero
by the phase atlas. This also removes any putative extra
signature from the comparison Gaussian.

For the leading singular integration write
\(x_i=y_{a_i}\), \(z_j=y_{b_j}\).
First Fourier integrate the \(z_j\)'s of a smooth compact
amplitude at the frequencies
\(r c\sum_i t_{ij}x_i\). Repeated integrations by parts
give arbitrary product decay in these frequencies. Integrate
one dominant dot per \(b_j\); these dots are distinct and
their product marginal is bounded. The resulting bound is
\(\prod_j C/(1+r\max_{i\sim j}|x_i|)\).
Each violating order ratio \(|x_i/x_j|>1\) then costs
the integrable reciprocal ratio, thus at least one logarithm
relative to \(d\) independent \(d|x_i|/|x_i|\) scales.
At the short radial end \(|x_j|\lesssim r^{-1}\) the
same bound gives one fewer log, and the finite upper end
has the same property. On the ordered chamber set
\(w_j=rx_j z_j\). Extending the compact radial domain in
these Fourier coordinates and replacing the radial amplitude
by its central trace have one-log error: the bounds
are integrable with relative factors \(|x_i|\) and
\((r|x_j|)^{-1}\), or the rapid Fourier moments of
\(|w_j|/(r|x_j|)\). Writing the matched-frequency relation
as
\[
 t_{jj}+\sum_{i\ne j}(x_i/x_j)t_{ij}=0
\]
shows that its matched-dot centres shift by
\(O(\sum |x_i/x_j|)\), again integrable with one logarithm
lost. One may first discard a bounded-width log strip near a
tie to keep these centres in the interior of every matched
dot interval. Fourier inversion there gives \(|c|^{-1}\)
per matched dot and the product density \(dt/2\) supplies
\(1/2\) per dot. This uses ordinary smooth conditional
densities near the \emph{matched} dot zero only. Nonmatched
angle poles can be handled by using spherical frames and
bounded smooth tests throughout, without a coordinate
singularity at a pole. The signed \(x_i\) ranges supply
\(2^d\); the ordered absolute log region has volume
\((\log r)^d\operatorname{Vol}(\Omega)+O((\log r)^{d-1})\).
Multiplying these factors leaves \(D_G=C_G\) and the
joint law \eqref{classic:eq:hallleading}, with the stated
loss if its trace vanishes. For \(d=0\) this step is
just regular evaluation. Relative exhaustion of the
finite forced mass and its equality with the exact
vacuum coefficient are the independent quantum estimates
in the Hall theorem; the two limits are not interchanged.
\end{proof}

The smooth product resolution makes a negative active lock
a real two-vector zero. Its mass on a compact fixed-rank
resolved chart is \(O(\epsilon^2)\) when that vector lies
in a ball of radius \(\epsilon\); a projection to active
radii can instead have a cap-convolution singularity.
One should use its independently resolved Haar estimate
rather than infer bounded projected density from the smooth
chart. When \(d>0\), a globally collinear all-negative
Hall reference has \(|t_e|=1\) on every dormant edge;
the graph block is invertible there and no dormant
logarithm arises in its fixed compact residual gauge model.
For \(d=0\) a collinear orbit has a circle stabilizer rather
than a smooth ordinary quotient. In the fixed slice use
Haar division before quotient: a generator of orbit speed
\(v\) and its cup dual contribute a torsion square-root
factor \(|v|\) on the nonfixed locus, the orbit-volume
Jacobian of this division. The effective orbit period is
half the circle's full period. This local pushforward
adds no atom on fixed orbits. These observations are not
uniform estimates near a change in the noncentral set.

\subsection{Tied scales carry a different angular support}
\label{classic:sec:overlap}
There is a direct subleading test of the same actual
normal phase. Choose any subset \(E'\) of unmatched arrows
and insert compact smooth cutoffs on the signed ratios
\(\rho_e=x_i/x_j\in\mathbb R^*\) for \(e=i\to j\) in \(E'\).
The constrained pairs form a forest. Let \(p=d-|E'|\) be
the number of its components, retaining isolated matched
pairs, and let \(\Omega_{E'}\) be the order polytope of the
roots under the unselected arrows. In any chart with smooth
density in the matched dots on the selected support (in
particular compact interior angular traces), the same
finite model has order \(r^{p_G}(\log r)^p\); its
coefficient is obtained by the replacements
\[
 \begin{gathered}
  2^d\operatorname{Vol}(\Omega)
      \longmapsto 2^p\operatorname{Vol}(\Omega_{E'}),\\
  \delta(t_{jj})\longmapsto
  \delta\!\left(t_{jj}+
       \sum_{\substack{i\to j\\(i,j)\in E'}}\rho_{ij}t_{ij}\right),
  \qquad
  \prod_{e\in E'}\chi_e(\rho_e)\frac{d\rho_e}{|\rho_e|}.
 \end{gathered}
\]
Here the densities and Fourier modulus per matched dot
are the same as before, a delta outside the angular
cutoff contributes zero, and the error has one fewer
root logarithm on fixed compact supports. Indeed, change
the absolute log coordinates to one root per selected
component and one edge difference per selected edge.
The incidence minor has determinant of modulus one;
the selected ratios fix descendant signs so only \(2^p\)
root signs remain. Unselected ratios have the same
dominant-bound and Fourier-shift estimates as above.
The selected coherent frequencies cannot be put equal
to zero in the integrand. Fourier inversion gives the
displayed deltas with unit matched-angular derivative.

If all arrows in one dormant connected component are
selected, the delta relations exhibit a genuine flat
curve: set all \(z\)'s to zero, let the \(x\)'s be a
common scalar times the vector of the selected signed
ratios, and keep the indicated \(u\) fixed. Every edge
period product vanishes and its only nontrivial radial
row is the displayed balance. The inverse analytic
chart therefore gives actual flat holonomies with this
physical first-order direction. A single selected arrow
in \(P_4\) can already have
\(t_{11}=0,\ t_{22}=-\rho t_{12}\ne0\), away from the
principal perpendicular-pair support. This is not a
positive splitting of the full unmarked next
coefficient; the normal off-shell ratios are not
automatically the same tests as the literal finite
quantum ratio marks.

\subsection{The absolute two-vertex and four-vertex tests}
\label{classic:sec:benchmarks}
These computations compare independent quadratic normalization and
exact finite level, so in particular check the volume of the effective
rotation group and the finite centre factor. The all-negative
\(P_4\) test also checks why an apparently singular intermediate
angle does not create an extra leading logarithm.

\paragraph{Two sites: the flat \(H\) model.}
The two-site \(H\) is the didicosm. Its degree-four rectangular torus
cover has lengths \(L_1,L_2,L_3\); the three nontrivial deck motions
have linear half-turns about the coordinate axes and half-period
translation on their fixed axis. These data, and the orthogonal
quaternionic flat lift, follow from the two glide presentations
in \cite[\S2]{P49}. Put \(V=L_1L_2L_3/4\).
The twist is the three-dimensional linear representation of the
deck quotient \(V_4\). Its harmonic one-forms are
\(X\mathbf i\,dx_1+Y\mathbf j\,dx_2+Z\mathbf k\,dx_3\) and
there is no invariant vector. These parallel forms close under the
quadratic curvature followed by Hodge star, so the complement
equation is solved identically by zero. The residual CS magnitude
is \(V|XYZ|/\pi^2\), directly by the normalized trace of the cubic.

The torsion can be computed without a seam convention. In the
weighted heat trace
\[
 \sum_q(-1)^q q\operatorname{Tr}(e^{-s\Delta_q}),
\]
the identity deck insertion has zero weighted form character.
A half-turn has form character \(-4\), and
adjoint character \(-1\); deck averaging gives \(+1\).
Its restricted Fourier sum is
\(\sum_{n\in\mathbb Z}(-1)^n e^{-s(2\pi n/L_i)^2}\).
Omitting its zero term gives in zeta variables
\[
 \zeta_i(z)=2(L_i/(2\pi))^{2z}
       (2^{1-2z}-1)\zeta_{\rm R}(2z),\qquad
 \zeta_i'(0)=4\log2-2\log L_i .
\]
Therefore \(T_{\rm adj}=\exp(\frac12\sum_i\zeta_i'(0))=16/V\).
The coexact curl has equal opposite helicity multiplicities in
both twisted and untwisted deck traces: on any nonidentity fixed
wavevector the half-turn acts with the same sign on both
transverse helicity lines. The scalar/exact parts also pair,
so the relative eta is zero. The harmonic coordinate metric
has volume \(V^{3/2}\,dX\,dY\,dZ\). The normalized zero-mode
amplitude is thus \(r^{3/2}V/(\sqrt2\pi^3)\,dX\,dY\,dZ\).
The cubic with \(\mathrm{CS}_0=\frac12\) gives
\((-1)^r\sqrt{2r}\log r/\pi\). Here and below a split
cubic \(cXYZ\), with a smooth cutoff around zero, has
coefficient \(2/(r|c|)\log r\) times the amplitude at zero:
Fourier transform one pair, integrate the remaining
signed \(dX/|X|\) between \(r^{-1}\) and a fixed radius,
and bound the two endpoint errors by the same rapid
Fourier estimate as in the Hall proof. This agrees
with the independent finite \(H\) computation in \cite{P48,P49}.

\begin{theorem}[Two normalized singular \(P_4\) models]
\label{classic:thm:p4}
For the all-negative central \(P_4\), let \(s_1,s_2\) be its
two outer cap-axis dots and \(t\) its middle dot. At
\(s_1=s_2=0\) and \(0<|t|<1\), the physical reduced
quartic magnitude in outer signed Wilson units is
\[
 \frac{|s_1s_2 y_1y_2|}{2|t|},\qquad
 \sqrt{T_{\rm sing}}=
 \frac{2\pi^4}{|t|}\,ds_1\,ds_2\,dy_1\,dy_2\,dt\,d\gamma .
\]
Here \(d\gamma\) denotes two passive azimuth probabilities
and \(\sqrt{T_{\rm sing}}\) is the coarse half-torsion.
After canonical gauge/framing normalization the leading
central model and the independent exact finite middle-dot
orbital law each give
\[
 \frac{(-1)^r\sqrt2}{4\pi^3}\,
      r^{5/2}(\log r)^2\,dt
\]
after passive angular integration. Hence the closed coefficient
is \((-1)^r\sqrt2/(2\pi^3)\) at \(r^{5/2}(\log r)^2\)
in \(Z_r\). In the literal disjoint trace box with signs
\((-,+,+,-)\), there is a phase-zero singular cubic at
perpendicular negative endpoint caps. Its total normalized
coefficient is \(+\sqrt2/(16\pi)\) at \(r^{1/2}\log r\),
matching the exact finite trace-box contraction.
\end{theorem}

\begin{proof}
For the all-negative corner, eliminate the two central
radial variables in the exact graph action:
the quadratic graph expression before its common coefficient
\(c\) is \(tX_1X_2+s_1X_1y_1+s_2X_2y_2\).
Its stationary value is \(-s_1s_2y_1y_2/t\), with retained
signed outer period linear units on the singular stratum.
For an independent absolute check, an all-central negative
Seifert piece contributes adjoint torsion 4 in primitive
peripheral bases. To see this, remove its genuine caps.
The base--fibre product cochain complex cancels against its
one-degree shift in the paired section/fibre bases. On the
plane perpendicular to the cap axis both cap cores act by
\(-1\), giving determinant 4 on each cap solid torus; their
tori are acyclic on that plane. On the invariant line the
two degree-one cap restrictions have determinants 2 and
2 with unit section pivots. The degree-zero and degree-two
pivots have modulus one; the opposite slopes cancel in
the surviving section relation. Thus the local product is
\(4^2/2^2=4\). The trivial positive central block, used
below, has instead three line copies and torsion \(4^{-3}\).

The same relative cone computation as in the Hall proof
then leaves two angular maps \(U,U^*\) and the dot
adjacency in degree \(1\to2\); the four piece contributions
have product half-factor \(2^4\). The radial pivot scalar
is \(|t|\) in half-torsion, and the retained two period
units give \(\pi^2\). For the angular determinant, lift
axis variations to \(\delta u_v=2a_v\times u_v\). Edge
differences plus the root generator are a triangular
unimodular change on the tree. Removing axial columns
and the global rotation leaves
\[
  \frac{\prod_{v=1}^4 dA(u_v)}{4^4\pi^2}
       =\frac{\pi^2}{8}\,ds_1\,ds_2\,dt\,d\gamma
\]
as a local quotient density, the right side using three
uniform edge dots. This yields \(2\pi^4/|t|\) as stated.
There is no second quotient by \(SO(3)\).
The clean outer-active branch is an independent
determinant check: per signed quadrant its coarse
density is
\(\pi^2\cot(\pi |y_1|/2)\cot(\pi |y_2|/2)
\,d|y_1|\,d|y_2|\,d\operatorname{Vol}_{SO(3)}\),
where \(d\operatorname{Vol}_{SO(3)}=(\pi^2/2)dt\,d\gamma\).
Its residue is \(2\pi^2/|y_1y_2|\). Dividing the singular
coarse density by the small curl product
\(2\pi^2|\kappa y_1y_2|\), with \(\kappa\) the physical
quartic coefficient, gives this residue exactly for
\(|\kappa|=1/(2|t|)\). The small eigenvalues have opposite
signs; the phase is \((-1)^r\) under
Section~\ref{phase:classical}.

For the finite quartic \(s_1s_2y_1y_2\), Fourier integrate
\(s_1,s_2\) at fixed signed \(y_1,y_2\). The four signed
quadrants each have logarithmic area \(1/2\) after the
gate \(r|y_1y_2|\gtrsim1\), yielding coefficient
\(2/(r|\kappa|)\,(\log r)^2\). The omitted gate strip and
the lower radial ends are controlled by the same Fourier
moments and lose one log. Combining the \(h=7\) Gaussian
normalization with the independent singular density leaves
the given constant multiple of \(dt\).

There is a global exact angular check without uniform
classical estimates at \(t=0,\pm1\). Use a smooth even
middle orbital test \(\eta(t)\) and its kernel
\(K_\eta(\xi)=\frac12\int_{-1}^1\eta(t)e^{i\xi t}dt\).
In the finite reflected deficit sum eliminate both
outer leaves by the odd sine identity
\[
 \sum_{\ell<r,\ \ell\ {\rm odd}}
     \cot\!\bigl(\tfrac{\pi\ell}{2r}\bigr)
                    \sin\!\bigl(\tfrac{\pi n\ell}{r}\bigr)
     =\tfrac r2,\qquad 0<n<r,\quad n\ {\rm odd}.
\]
For instance the identity follows by subtracting adjacent
odd frequencies using sine/cosine addition on the odd
grid and evaluating the first frequency. Using the
prefactor and vertex factors of \eqref{rad:sum}, the
remaining exact expression is
\[
 (-1)^r\frac{4\sqrt2}{\pi^3}r^{5/2}
 \sum_{\substack{n,m<r\\ n,m\ {\rm odd}}}
 \frac{F_2(n/r)F_2(m/r)}{nm}\,
 K_\eta(\pi nm/r),\qquad
 F_2(x)=\left(\frac{\pi x/2}{\sin(\pi x/2)}\right)^2 .
\]
Indeed the leaf identity in sinc units is \(r/(4n)\);
the middle vertex has precisely the bounded endpoint
factor \(F_2\). The normalized odd harmonic measures
converge to \(\tfrac12\,du\) with \(u=\log n/\log r\).
Away from the endpoint strips and \(u+v=1\), the
bounded kernel tends uniformly to \(K_\eta(0)\) on
one triangle and 0 on the other. The strips have small
normalized mass and \(F_2\) tends to one on interior
logarithmic scales. The triangle's weight is \(1/8\).
This gives the exact claimed density for even orbital
tests, including exhaustions of the open generic
angular set. The check does not substitute an orbital
integral variable for a genuine finite Wilson operator.

For the mixed sign box, fix one root sign at each positive
site. Use signed radian endpoint periods \(p,q\), axes
\(u,v\), and \(t=u\cdot v\). The positive centres' vector
periods are eliminated by reversing their ordered
products: their caps have \({\rm Ad}+1\) invertible,
forcing the two middle periods to be \(-qv,-pu\).
The edge commutators then give \(pq=0\), with a unit
from a component of \(u\times v\) at perpendicular
axes. The two negative caps' axial equations give
\(pt=qt=0\); the transverse equations hold modulo these
generators. Thus the exact critical scheme is
\(\operatorname{Jac}(pqt)\). The genuine CS germ has no
quadratic terms in these three retained coordinates.
Only \(pqt\) can occur in its cubic with these gradient
generators, and its coefficient cannot vanish.
The analytic straight-homotopy argument for this
graded ideal gives the exact cubic with a
tangent-identity change (equivalently the normal-form
argument of \cite{P49B}); there is no change in the
physical period units. \(\mathrm{CS}_0=0\) since the
two positive product bits are both one, and the
relative rho vanishes by \eqref{phase:rho}.

The singular half-torsion at \(p=q=t=0\) per root pair
is \(dp\,dq\,dt/32\). In fact the four blocks now give
product half-factor \(2^2(1/8)^2=1/16\).
The degree-one compressed adjacency has half-block
\((x,b)\mapsto(u\cdot x,x+bv)\); its nonzero singular
values are \(\sqrt2,\sqrt2,1\).
The two retained primitive radial tangent vectors
have norm \(\sqrt2\) and are orthogonal, cancelling
the determinant divisor 2. The two axis areas
divided by \(4^2\pi^2\) have total weight one,
with normalized dot density \(dt/2\). This verifies
the stated singular factor. On the adjacent reflected
smooth arm \(p\ne0,\ q=t=0\), the coarse half-torsion
per root pair tends to \(dp/(32|p|)\).
The compact-regular formula is
\((\pi/64)\cot(\pi x/2)\csc^2(\pi(1-x)/2)\,dx\)
with \(x=|p|/\pi\). Its primitive block computation
has two ordinary noncentral factors
\(\frac12(4\cos^2(\pi x_j/2))
 (4\sin^2(\pi x_j))^{-d_j/2}\),
the positive centre factor \(1/8\) and the negative
end factor 2; the invariant radial path block has
the common unit period in radians and all
transverse section/area pivots have determinant
one in coarea bases. Substituting \(d_j=1,2\) and
\(x_j=x,1-x\) gives the displayed formula.
The small-curl determinant comparison now yields
\(|\kappa|=1/(2\pi^2)\) in \(pqt\) radian units.
With the \(h=3\) canonical factor and the split
cubic residue, the coefficient is
\(1/(32\sqrt2\pi)\) per positive-root pair. The four
such pairs are distinct. Ordered perpendicular
negative axes form a single effective rotation
orbit; the two signed period variables already
include both adjacent Weyl directions. The total
is \(+\sqrt2/(16\pi)\).

For comparison the exact finite trace box can be
reduced by the same two leaf sums with cutoffs.
After reflecting the two small positive-centre
labels \(a,b\), the middle sine changes by
\((-1)^r\), cancelling the external parity.
The remaining form is
\[
 \frac{4\sqrt2}{\pi^4}r^{-1/2}
 \sum_{a,b\equiv r-1\, (2)} C_2(a/r)C_3(b/r)
   \sin(\pi ab/r)\,Q_1^{(r)}(a)Q_4^{(r)}(b),
\]
where the positive-centre profiles are smooth with
\(C_j(0)=\pi^2/4\), and cut leaves give
\(Q_i^{(r)}(a)=1+O(1/a)\) on the sufficiently small
positive arguments. The latter estimate follows by
summation by parts in the smooth difference from
the uncut cotangent identity. For \(a\) in this range,
two geometric summations by parts in \(b\) give
\[
 \sum_{b\equiv r-1\, (2)} C(b/r)\sin(\pi ab/r)
    =\frac{r\,C(0)}{2\pi a}+O_C(r/a^2+1):
\]
the first endpoint term is the imaginary part of
\(C(s/r)e^{is\pi a/r}/(1-e^{2i\pi a/r})\),
\(s=1\) or 2, and the second difference has total
variation \(O(r^{-1})\). The single leaf errors cost
\(O(r\sum a^{-2})\), the double error \(O((\log r)^2)\).
Now the half-grid harmonic sum gives \(\frac12\log r+O(1)\)
on either parity. The result is exactly
\(+\sqrt2\,r^{1/2}\log r/(16\pi)+O(r^{1/2})\).
It is a literal trace box, not a finite projector onto
a single character, and two independently integrated
smooth endpoint arms must not each be credited with
the whole logarithm.
\end{proof}

\section{Actual Wilson insertions and their joint angular limit}
\label{sec:ang:main}

In this section we keep the inserted Wilson ribbons in the surgery diagram. In particular, a function of a relative cap angle is not inserted as an assumed diagonal multiplier in a backbone-color sum. The result is an ordinary, undamped level limit for every fixed finite collection of integer-spin Wilson words on a perfectly matched backbone. The construction and the uniform signed estimates also apply to the mixed Hall sector in the next section.

\subsection{The marked diagram and its exact tree transfer}

We use one-based colors \(V_m\), \(1\leq m<r\), and put
\[
 \theta=\pi/r,\qquad q=e^{i\theta},\qquad
 d_m=\frac{\sin(\theta m)}{\sin\theta},\qquad
 {\cal D}=\frac{\sqrt{r/2}}{\sin\theta}.
\]
The actual surgery backbone color at \(v\) is \(a_v=r-m_v\). The unpierced \(-2\) cap, at each vertex, enforces \(m_v\) odd. Draw the tree in disjoint planar vertex boxes and edge strips, with only the designated two-crossing positive Hopf clasp in each strip. At any positive cap disk mark finitely many passages parallel to the backbone; a fixed-color auxiliary ribbon then follows these passages in its specified word order. Between boxes it can use the outside side of the unique tree corridor, with distinct heights for different strands. An inverse letter may be implemented with the dual passage. Other twists and crossings in the auxiliary part are fixed, and no additional heavy-heavy crossing is introduced. These conventions define the framed representative to which all finite-level results below apply. No invariance under changing an auxiliary route while preserving only its free homotopy word is needed.

A collection \(W\) can contain several ribbons, with arbitrary shared cap visits. Auxiliary spins are integers, so their representations are insensitive to the central sign in lifting an adjoint holonomy. On the all-negative flat slice, all backbone meridians are central and may be filled in the adjoint complement. The separated cap circles then have disjoint spanning disks, and reading a routed ribbon's crossings of those disks gives exactly its ordered cap word, up to a common conjugation and central lifts. Thus its classical evaluation is
\begin{equation}
 \label{eq:ang:word}
 P_W(u)=\prod_{\ell\in W}
   \chi_{j_\ell}\!\left(\prod_{v\text{ along }\ell}^{\rm ordered}
                                   i u_v\cdot\boldsymbol{\sigma}\right),
          \qquad u_v\in S^2 ,
\end{equation}
where \(\chi_{j_\ell}\) is the ordinary spin-\(j_\ell\) character and the \(\boldsymbol{\sigma}\) are the Pauli matrices. For integer \(j_\ell\), the inverse of a half turn has the same representation matrix. In particular \(P_W\) is a polynomial in the axes, invariant under their independent reversals and under simultaneous rotations.

To see how the marking changes the finite sum, let
\[
 f_\pm(a)=\frac{(S\Theta^{\pm2}S)_{a1}}{S_{1a}},
\]
where the twists \(\Theta\) have the same convention as in the unmarked surgery sum. A positive cap through which a fixed auxiliary packet \(C\) passes acts on \(V_a\otimes C\) by \(f_+(x)\) on the total simple channel \(V_x\), including every multiplicity of that channel. Divide by the unpierced \(f_+(a)\); the surviving parity of \(a\) makes the denominator nonzero. The two-term quadratic Gauss sum for the \(+2\) cap gives the exact ratio
\begin{equation}
 \label{eq:ang:cap}
 \frac{f_+(x)}{f_+(a)}
  =\frac{\sin(\theta m/2)}{\sin(\theta(r-x)/2)}
          \exp\!\left(-\frac{i\theta}{4}(x^2-a^2)\right),
                     \qquad m=r-a .
\end{equation}
For example, on the nonzero parity, \((S\Theta^2S)_{a1}\) is a common nonzero scalar times \(\sin(\theta a/2)\exp(-i\theta a^2/4)\); the two Gauss translates of \(a\) differ by \(2r\). Dividing by \(S_{1a}\) and taking the ratio gives \eqref{eq:ang:cap}. Since \(|x-a|\) is bounded by the packet size, these ratios are uniformly bounded on all admissible charges.

Now remove the invertible object \(\mathsf J=V_{r-1}\) from every \(V_a=\mathsf J\otimes V_m\). Its monodromy with every integer-spin object is the identity, so its crossings with both the stripped heavy wires and the bounded auxiliary packets can be switched as exact morphisms. The cap projector is carried from the \(x\)-channel to the \(r-x\)-channel by tensoring with \(\mathsf J\). The separated invertible-object link contributes one factor \(M_{\mathsf J,\mathsf J}=(-1)^r\) per heavy Hopf clasp, since \(d_{\mathsf J}=1\) and the heavy framings are zero. This reduction takes place in the unitary ribbon category before any magnetic-basis limit.

Let \({\cal K}_{r,W}(\mathbf m)\) be the remaining diagram value with all positive caps normalized by \eqref{eq:ang:cap} and one normalized categorical trace on each heavy backbone. The physical marked invariant is
\begin{equation}
 \label{eq:ang:normalized-sum}
 Z_{r,W}=(-1)^{r(N-1)}
      (2/r)^{(N+1)/2}\sin(\theta)^{-(N-1)}
      \sum_{\mathbf m\ {\rm odd}}
        \prod_v \bigl(1+\cos(\theta m_v)\bigr)\,
                            {\cal K}_{r,W}(\mathbf m).
\end{equation}
Any diagonal radial mark multiplies the summand in this formula. To check the scalar normalization, the normalized unpierced Hopf clasp at \(m,n\) is
\[
 \lambda_{r,0}(m,n)=
 \frac{\sin\theta\,\sin(\theta mn)}{\sin(\theta m)\sin(\theta n)}
   =D_r(m,n)\frac{\sin(\pi mn/r)}{\pi mn/r},
 \qquad
 D_r(m,n)=\frac{\pi mn}{r\,\sin\theta\, d_m d_n}.
\]
Thus factoring the original unmarked vertex weights
\((1+\cos(\theta m))/\sin^{\deg(v)}(\theta m)\) against one sine per edge produces \eqref{eq:ang:normalized-sum}. Equivalently factor each unchanged scalar \(f_-(a)f_+(a)\) and the backbone surgery weight before evaluating the decorated diagram. In neither calculation is there division by a possibly vanishing Hopf entry.

\begin{lemma}[Finite-rank tree reduction]\label{lem:ang:network}
For the declared routes there is a fixed finite index set, independent of \(r\) and of the heavy charges, such that
\begin{equation}
 \label{eq:ang:network}
 {\cal K}_{r,W}(\mathbf m)
  =\sum_\alpha
       \prod_v B_{v,\alpha}(m_v;r)\,
       \prod_{e=vv'}\lambda_{r,L_e(\alpha)}(m_v,m_{v'}).
\end{equation}
Here each \(L_e(\alpha)\) is bounded in terms of \(W\), each \(B_{v,\alpha}\) is bounded, and \(\lambda_{r,L}\) is the one-port rank-\(L\) transfer through the true heavy Hopf clasp.
\end{lemma}

\begin{proof}
At a tree cut, move its finitely many auxiliary strands to the outside of the clasp as in the declared route. Bend their outgoing duals; the resulting fixed packet \(C_e\) has finitely many simple components and fixed multiplicity spaces. Performing the subtree diagram up to that cut leaves the parent-facing heavy through-line and a morphism \(V_m\otimes C_e\to V_m\). This description does not bend either heavy leg through its neighbor. After selecting an auxiliary simple \(V_{2L+1}\), the space \(\operatorname{Hom}(V_m\otimes V_{2L+1},V_m)\) is zero or one-dimensional. In the latter case use a normalized coisometry \(T_{m,L}\), with \(T_{m,L}T_{m,L}^*=1_{V_m}\).

The intervening clasp and normalized traced child heavy line take \(T_{n,L}\) to a scalar multiple of \(T_{m,L}\); the packet multiplicity is only an exposed finite index. Closing the other line against \(T_{m,L}^*\) identifies the scalar with the Hopf graph with the two normalized trivalent vertices and the rank-\(L\) bridge. This defines \(\lambda_{r,L}\), extended by zero when a channel is absent. Bending and routing the fixed packet uses only its braids with the appropriate heavy through-line, not another heavy monodromy. Recursing through the tree puts all couplings among packet indices at their respective vertices. The final normalized trace at a vertex includes only that vertex's own heavy \(m_v\), its cap and its finite auxiliary couplings. This gives the product and sum in \eqref{eq:ang:network}, even when several Wilson words use the same cut. The fixed packet dimensions and unitary categorical norm bound all factors. In a forest with shared auxiliary words the same local construction gives a finite tensor contraction coupling the components, rather than independent observables on them.
\end{proof}

We will need the following elementary strengthening of the vertex bound:
\begin{equation}
 \label{eq:ang:vertex-bv}
 \sup_m|B_{v,\alpha}(m;r)|+
      \operatorname{Var}_{m\ {\rm on\ its\ fixed\ parity}}B_{v,\alpha}(m;r)
          \leq C_W .
\end{equation}
It holds on every physical interval or folded branch. One can see the finite complexity explicitly by evaluating each vertex through a fixed list of elementary recouplings. Only one heavy through-line occurs: its neighboring channel labels are all \(m+O_W(1)\) and the auxiliary labels are bounded. In the finite quantum Racah formula (e.g. \cite[Sec.~2, Eqs.~(6)--(7)]{TaylorWoodward2006}) for a one-small term \(\{A\,B\,E;\,J,E+\xi,B+\eta\}\), set \(T=A+B+E\), \(U=A+B-E\), \(V=A+E-B\), \(W_0=B+E-A\). The Racah summation parameter is \(T+\kappa\) with bounded \(\kappa\); its long factors have arguments \(U+\kappa-J-\xi,V+\kappa-J-\eta,W_0+\xi+\eta-\kappa,T+\kappa+1\). They cancel the corresponding long triangle factorials up to bounded shifts. Admissible denominators are nonzero below level; terms whose numerator reaches the zero quantum integer are omitted before cancellation. The remaining expressions have only fixed-length \(q\)-integer products and their square roots, on a bounded number of linear admissibility pieces. The same holds for the normalized trivalent gauge; its relative norm is given below in \eqref{eq:ang:macdonald}. Fixed-large braid phases and \eqref{eq:ang:cap} have fixed frequencies in \(m/r\) at each bounded shift; no heavy self-twist survives zero framing. In the normalized real Racah gauge there is no hidden sign alternating on the step-two progression. An elementary move with one heavy input has one heavy total spin, and the normalized Racah prefactor has sign exponent equal to the sum of the three external input spins and that total spin, so its slope in the heavy spin is two. The large triangle used for \(T+\kappa\) likewise has exactly two heavy labels and slope two. In a fixed-large braid the eigen-sign exponent contains the heavy input minus the heavy output, each with slope one. These signs, as well as the duality sign \((-1)^{2j}\), therefore remain constant on a physical parity progression. The trivalent gauge of \eqref{eq:ang:macdonald} takes the highest allowed coefficient positive up to the fixed duality and braid convention; only fixed-shift norm and fixed-frequency phase factors remain. After \(y=\tan(\theta m/8)\), the resulting real and imaginary parts are semialgebraic families of uniformly bounded complexity. Uniform monotonicity gives a bounded number of intervals, and on the sampled points of each such interval variation is controlled by the bounded sampled values. This yields \eqref{eq:ang:vertex-bv} for the gauge used below.

\subsection{Joint compact-product limit}

For \(e=vv'\) put \(w_e=m_vm_{v'}/r\), and let \(d\Pi=d\Omega/(4\pi)\) be spherical probability. The common limiting kernel is
\begin{equation}
 \label{eq:ang:angular-kernel}
 {\cal K}_{W}(\mathbf w)=
  \int_{(S^2)^V}P_W(u)
    \exp\!\left(i\pi\sum_{e=vv'} w_e\,u_v\cdot u_{v'}\right)
                                               \prod_v d\Pi(u_v).
\end{equation}

\begin{theorem}[Actual joint symbol]\label{thm:ang:compact}
If \(\min_v m_v\to\infty\), \(\max_v m_v/r\to0\), and all \(w_e\) remain bounded, then
\[
       {\cal K}_{r,W}(\mathbf m)-{\cal K}_W(\mathbf w)\longrightarrow0
\]
uniformly on these scales. There is no comparable-spin requirement, and the theorem applies to any fixed number of simultaneous ribbons.
\end{theorem}

\begin{proof}
After the exact \(\mathsf J\)-removal, every fixed cut of the reduced diagram contains a total highest weight \(o(r)\). The tensor products are inside the alcove, so we may use ordinary type-one \(U_q(\mathfrak{sl}_2)\) weight matrices without negligible summands. Elementary matrix norms in this basis, not transported categorical unitary norms, are used for the local norm limits.

Let \(j=(m-1)/2\) and \(U=\mathbf J/j\), with \(\mathbf J\) the usual spin matrices. The \(q\)-raising and lowering entries differ from the ordinary normalized entries by \(O((j/r)^2)\). The triangular universal \(R\) from \cite[Sec.~7]{RT1990} has Cartan factor \(q^{2J_z\otimes J'_z}\), and its \(b\)-th nilpotent coefficient is \((q-q^{-1})^b/[b]!\) times a unit phase and the harmless \(q^{\pm J_z}\)-dressings of the symmetric coproduct. Up to the stopping index \([b]\geq c b\); raising and lowering norms give a dominating sequence \((Cw)^b/b!\) for \(w=mn/r\). For each fixed \(b\), the terms approach their ordinary exponential terms. Consequently the unflipped \(R\) is, in norm up to \(o(1)\), \(e^{2i\theta J_zJ'_z}e^{2i\theta J_+J'_-}\). The corresponding opposite-order factor and \([U_i,U_j]=O(1/j)\) then give
\[
                  R_{21}R
                         =e^{\,i\pi w\,U\cdot U'}+o(1).
\]
The coefficient agrees with the physical balancing spectrum
\(\exp(i\pi(x^2-m^2-n^2+1)/(2r))\) on the \(x\)-channel. Bounded auxiliary crossings become ordinary flips. The factorial bound gives uniformity for all bounded \(w\), even with widely separated heights of the two heavy labels.

For a cap packet with fixed total ordinary spin matrices \(S\), use \(C_q=FE+[J_z+1/2]^2\) with its symmetric coproduct. The packet channel \(r-x=m+2s\) has eigenvalue \([j+s+1/2]^2\) for the coupled Casimir. On the entire tensor product,
\[
 \frac{\Delta C_q-[j+1/2]^2}{2j}
                             = U\cdot S+O(j/r+j^{-1}).
\]
Indeed the change of the pure heavy raising-lowering term from multiplication by \(q^{2S_z}\) is \(O(j^2/r)\) before dividing, the mixed terms reduce to \(J_-S_++J_+S_-\), and the divided difference of the sine-squared Cartan term gives \((J_z/j)S_z\). Subtracting the exact single-spin Casimir cancels the large pure term. This is why no hypothesis on \(j^2/r\) is imposed. There are only finitely many bounded integer shifts \(s\), and their normalized eigenvalues approach these distinct integers. The ratio \eqref{eq:ang:cap} on each of them is \((-1)^s+o(1)\). Polynomial interpolation on this separated finite set gives its operator limit with symbol \(\exp(i\pi u\cdot S)\). On a full packet that is the tensor product of the same half turn in each constituent spin representation; hence the simultaneous words, not just their separate means, enter the calculation.

Finally the spin-coherent resolution of the identity gives the normalized heavy trace of each fixed ordered polynomial of the \(U_i\)'s as its spherical integral up to \(o(1)\). Indeed rotation-averaging the highest-state projector gives \(1/(2j+1)\) times the identity by irreducibility, and its rotated coherent states have mean \(u\) and component variances \(O(j^{-1})\); apply Cauchy--Schwarz and telescope each ordered monomial. Bounded exponential factors can be approximated by fixed polynomial products uniformly at bounded \(w\). The magnetic pivotal trace factor differs from the ordinary one by \(O(j/r)\), and \(d_m/m\to1\). The finite auxiliary contractions now give exactly the prescribed characters \eqref{eq:ang:word}. This proves \eqref{eq:ang:angular-kernel} and the stated uniform limit.
\end{proof}

For later use, in a rooted tree the dot \(t_e=u_v\cdot u_{v'}\) of a fresh child with its parent has probability \(dt_e/2\) on \([-1,1]\), and the child's transverse azimuth is an independent circle variable. These variables are independent over all edges in the vacuum, but the inserted \(P_W\) is a joint function of the shared axes. Its conditional expectation \(p_W(\mathbf t)\) is a polynomial: successively average leaves using the transverse plane projector \(1-uu^{\mathsf T}\); odd powers of the transverse unit circle vanish and even powers give powers of \(1-t_e^2\) times projector tensors. This also makes the conditional formula regular across any seam \(t_e=0\).

\subsection{Uniform signed tails from the punctured S-matrix}

An absolute \(1/w\) estimate on one clasp cannot remove its radial tail, since the associated harmonic integral diverges. We now prove a signed prefix estimate in the whole physical alcove, including large and reflected colors.

\paragraph{Two exact carriers.}
The bulk-independent fixed-rank input is the once-punctured modular \(S\)-matrix. With the same bridge contraction as in Lemma~\ref{lem:ang:network}, bend the two through-heavy ports into the handle \(V_m\otimes V_m^*\). Applying the Hopf clasp and closing the second line is the punctured \(S\) morphism of \cite[Eqs.~(1.18), (2.1), (5.7)--(5.9)]{Kirillov1996}. In that formula the \(d_n/{\cal D}\) factor replaces the normalized partial categorical trace. Closing against the other normalized trivalent intertwiner adds the second dimension denominator. The highest-zero-weight basis in that paper has relative squared norms
\(h_m=\prod_{\ell=1}^L [m+\ell]/[m-\ell]\), and bending preserves those relative norms with a common packet factor. Therefore in the normalized coisometry convention
\begin{equation}
 \label{eq:ang:macdonald}
 \lambda_{r,L}(m,n)
     =\frac{{\cal D}}{d_m d_n}\sqrt{\frac{h_m}{h_n}}\ S^{\rm raw}_{mn}
     = c_{r,L}
       \frac{\prod_{\ell=0}^L\sin(\theta(m-\ell))}
            {\sin(\theta m)\sin(\theta n)}
       \sqrt{\frac{h_m}{h_n}}\ P_{n-L-1}(q^{-m}).
\end{equation}
Here \(P_l\) is the monic \(A_1\) Macdonald polynomial at parameters \(q^2,q^{2(L+1)}\). In the raw convention its overall scalar is
\[
          c_{r,L}=\tfrac i2(-2i)^{L+1}q^{-L(L+1)/2}\sin\theta .
\]
Reversing the clasp conjugates the expression; a fixed puncture twist contributes a constant phase. The \(L=0\) case is exactly the normalized ordinary Hopf sine above, independently checking the overall dimension factor. A term is zero if one of \(m,n\) lies outside \(L+1,\ldots,r-L-1\).

The reflection is exact: \(h_{r-m}=h_m^{-1}\), the sine product changes by \(h_m\), and monic polynomial parity and Weyl symmetry give
\begin{equation}
 \label{eq:ang:reflection}
 \lambda_{r,L}(r-m,n)=(-1)^{n-L-1}\lambda_{r,L}(m,n).
\end{equation}
The normalized transfer is symmetric under interchanging the two boundary ports, equivalently by the norm identity for the raw matrices in \cite[Eq.~(5.12)]{Kirillov1996}, so the other reflection has the same form. Thus on each independent folding branch the signs are constant on the relevant parities. Every original interval mark becomes its corresponding folded interval. No endpoint is counted twice.

Fold now \(m,n\leq r/2\). For fixed \(L\) the Rogers generating function \cite{IsmailStanton1988} takes a rational form. Writing \(l=n-L-1\),
\begin{equation}
 \label{eq:ang:rogers}
 P_l(z)= B_L(l)^{-1}[t^l]\prod_{j=0}^{L}
                  (1-q^{2j}zt)^{-1}(1-q^{2j}z^{-1}t)^{-1},
 \quad B_L(l)=\prod_{j=1}^{L}
                   \frac{1-q^{2(l+j)}}{1-q^{2j}} .
\end{equation}
The products telescope at integer parameter \(\beta=(q^2)^{L+1}\), and \(B_L(l)\) is exactly the highest Laurent coefficient by the finite binomial identity. It has magnitude comparable to \(n^L\) when the channel is present. Let \(z=q^{-m}\), and at \(m>4(L+1)\) group the simple poles into \(q^{2j}z\) and \(q^{2j}z^{-1}\) (reciprocal positions). The partial fraction coefficients of each group in \eqref{eq:ang:rogers} are \(z^{\pm l}\) times the order-\(L\) divided difference at \(u_j=q^{2j}\) of
\[
            \frac{u^{\,l+2L+1}}
                      {\prod_{j=0}^{L}(u-q^{2j}z^{\mp2})}.
\]
This grouping is essential: estimating the individual poles first loses powers at the edge of the alcove. The opposite-group separations on the convex hull of the nodes are bounded below by \(c_Lm/r\). By the simplex formula for divided differences, differentiation \(p\) times in the numerator costs \(O_L(n^p)\) and the remaining derivatives cost \(O_L((r/m)^{2L+1-p})\). For \(w=mn/r\geq1\) this is \(O_L(n^L(r/m)^{L+1})\). Combining with \eqref{eq:ang:macdonald}, \eqref{eq:ang:rogers}, and the folded sine comparabilities proves the exact representation
\begin{equation}
 \label{eq:ang:carriers}
 \lambda_{r,L}(m,n)=E_+(m,n;r)e^{i\theta mn}
                         +E_-(m,n;r)e^{-i\theta mn},
                 \qquad |E_+|+|E_-|\leq C_L/w\quad(w\geq1).
\end{equation}
After factoring those carriers, the remaining \(n\)-dependent exponentials in the pole expression are \(q^{2jl}\) with fixed \(j\leq L\). Thus the envelopes are fixed-complexity trigonometric rational families, with bounded-complexity square roots and feasibility intervals. Their sampled sup and total variation on \(w>K\geq1\) are both \(O_L(K^{-1})\) by uniform monotonicity. On \(w\leq1\) keep the carriers unsplit: their argument is in a fixed compact interval, so they are definable in the same restricted analytic family and have uniformly finitely many monotonicity intervals \cite{vdDMiller1996}. Their sampled norm is bounded by the unitary category, giving uniformly bounded variation. For \(m\leq4(L+1)\) one may use symmetry of \(\lambda_{r,L}\) in \eqref{eq:ang:macdonald}, which makes its polynomial degree uniformly bounded and therefore gives the same bounded variation on the whole folded interval without coincident-pole denominators.

\begin{lemma}[Signed whole-alcove prefix estimate]\label{lem:ang:prefix}
Fix a vertex of degree \(D\) with neighboring folded labels \(m_a\) and bounded ranks \(L_a\), and let \(M=\max_a m_a\). Let \(A(n)\) be any folded one-vertex weight of bounded sup and discrete total variation. On either folded parity, with any hard interval in \(n\leq r/2\), one has
\begin{equation}
 \label{eq:ang:prefix}
 \left|\sum_{Mn/r>K} A(n)\prod_{a}\lambda_{r,L_a}(m_a,n)\right|
          \leq C\,\frac{r}{MK}\quad(K\geq1),\qquad
 \left|\sum_n A(n)\prod_a\lambda_{r,L_a}(m_a,n)\right|
          \leq C\,\frac{r}{M}.
\end{equation}
The constant is linear in the sup-plus-variation bound on \(A\), and otherwise depends only on the fixed tree and ranks.
\end{lemma}

\begin{proof}
If a second neighbor is at least \(M/(2D)\), two transfers give absolute \(O_D((1+w)^{-2})\) at \(w=Mn/r\), giving the tail estimate directly. In the remaining case the sum of the smaller neighbors is less than \(M/2\). Partition at their finitely many transitions \(m_a n/r=1\), and expand the carriers of the high-product factors. The maximizing edge has envelope sup and variation \(O(K^{-1})\); other high envelopes and low factors have bounded sup and variation. Any combined frequency satisfies \(M/2\leq|\nu|\leq3M/2\leq3r/4\); hence geometric summation of \(e^{i\theta\nu n}\) on the actual step-two lattice costs at most \(Cr/M\). Multiplying the envelope bound by this geometric bound and summing by parts proves the first inequality. Bounded exceptional neighbors can be kept unsplit; if the maximizing neighbor is bounded, the large tail is empty and the finitely bounded \(K\)'s are covered by absolute counting. Finally \(w\leq1\) costs \(O(r/M)\) in absolute counting. The same proof gives the bound with an interval prefix or suffix, and the reflection signs \eqref{eq:ang:reflection} are constant on each progression.
\end{proof}

This estimate is for an actual signed conditional sum, not the total variation of the full original multiple-color sum. In particular it remains valid after multiplying by finitely many one-variable BV profiles. For smooth joint radial tests we use the anchored mixed-variation representation by rectangular prefixes: iteratively apply the fundamental theorem of calculus in each variable, with its mixed derivatives or boundary values as coefficients. Uniform prefix estimates then extend with the corresponding finite mixed variation factor. The representation also applies to smooth functions of log radii, since coordinatewise monotone reparametrization preserves the bound.

\subsection{Ungated angular law for a perfectly matched backbone}

Let \(G\) have its unique perfect matching \(M\), \(N=2d\). Choose a bipartition and name the matched edges \(a_i b_i\), \(i=1,\ldots,d\). On their true log variables \(u_i=\log m_{a_i}/\log r\), put
\[
      \Omega_M=\{u\in[0,1]^d:\ u_i\leq u_j
                               \text{ if }a_i b_j\notin M\}.
\]
Write \(\iota(u)\) for the full log-radius array with matched pairs \((u_i,1-u_i)\). Let \(\mathbb E_{\rm fr}\) denote independent Haar orthogonal frames on the matched pairs, with quotient by simultaneous rotation and individual axis reversals when evaluating even Wilson words. A matched frame means the two ordered perpendicular axes and their right-handed cross product.

\begin{theorem}[Ungated joint Wilson law, balanced case]\label{thm:ang:perfect}
For the declared finite collection \(W\), arbitrary fixed smooth joint marks \(F\) of all original log radii, and \(G_0\) of their ordinary radii \(m_v/r\), one has
\begin{equation}
 \label{eq:ang:perfect}
 Z_{r,W}[F,G_0]
  =(-1)^{r(N-1)}
       \frac{2^{(N+1)/2-d}}{\pi^{N-1}}\,
             r^{\,N-3/2}(\log r)^d
       \left(G_0(0)\int_{\Omega_M}\! F(\iota(u))\,d^du\
                        \mathbb E_{\rm fr}P_W +o(1)\right).
\end{equation}
No radial, product or Abel cutoff is present in the invariant on the left.
\end{theorem}

\begin{proof}
First keep all original colors on the lower branch in a power box \(r^\delta\leq m_v\leq r^{1-\delta}\). Fix all \(a_i\)-labels and use \eqref{eq:ang:network}. The \(b\)-sums separate termwise. As a result of \eqref{eq:ang:prefix}, if an unmatched competitor \(a_i\) is larger than the matched \(a_j\) at \(b_j\) by \(r^\epsilon\), the conditional bound is suppressed by \(m_{a_j}/M_{b_j}\leq r^{-\epsilon}\) against the usual matched harmonic bound. The fixed finite union of order walls of width \(\epsilon\) has at most \(O(\epsilon)(\log r)^d+O((\log r)^{d-1})\) harmonic mass. Thus only strict allowed order chambers matter after sending \(\epsilon\) to zero. In those chambers the matched label maximizes the frequency at its partner. Lemma~\ref{lem:ang:prefix} discards \(w_i=m_{a_i}m_{b_i}/r>K\) at normalized cost \(O(1/K)\); the opposite \(w_i<\kappa\) has cost \(O(\kappa)+o(1)\) by absolute counting. Other \(b\)-sites are always kept as signed conditional sums. Prefix integration permits the same statements for the joint tests.

Inside the fixed window \(\kappa\leq w_i\leq K\), the nonmatching products tend to zero uniformly in a strict chamber, and Theorem~\ref{thm:ang:compact} applies. At fixed \(a_i\), the odd \(b_i\)-mesh in \(w_i\) has spacing \(2m_{a_i}/r\); summing odd \(a_i\) against \(1/m_{a_i}\) gives \((\log r)\,du_i/2\). Thus each matched lattice pair contributes \(r\log r\,du_i\,dw_i/4\) before the cap profiles. Each of its two profiles tends to \(2\). Moreover \(\log m_{b_i}/\log r=1-u_i+\log w_i/\log r\), so the retained log mark approaches \(F(\iota(u))\) and the ordinary mark approaches \(G_0(0)\).

With only the matched products nonzero in \eqref{eq:ang:angular-kernel}, independent sign flips turn the exponential in each \(w_i\) into a cosine. The undamped improper integrals have the same value as their Abel limits by the uniform signed tail just proved (or by the smooth conditional Fourier representation). For a single coordinate,
\[
       \int_0^\infty e^{-\varepsilon w}\cos(\pi wt)\,dw
             =\frac{\varepsilon}{\varepsilon^2+\pi^2t^2}
                          \longrightarrow\delta_0(t),
\]
with mass one on the real line. The dot density is \(dt/2\); each matched edge therefore supplies \(1/2\) and conditions its two axes to be perpendicular. These conditions are on disjoint pairs, leaving independent Haar frames and the one shared \(P_W\). The resulting raw lattice factor is \((r\log r/2)^d\). The common prefactor in \eqref{eq:ang:normalized-sum} gives \eqref{eq:ang:perfect} on the interior chambers.

It remains to remove the initial radial power box. In any fold branch, sum one bipartition first and use \(\prod_b(r/M_b)\leq r^d\prod_a m_a^{-1}\), with folded labels. For any restricted endpoint vertex, choose the bipartite direction that leaves it unsummed; its folded harmonic mass outside \([r^\delta,r^{1-\delta}]\) is \(O(\delta\log r+1)\). Repeat on the finitely many intersections of endpoint restrictions. The same estimates apply to all radial prefixes, and to their mixed-variation combinations. A branch with an originally upper label but all folded labels in the power box has the original profile \(1-\cos(\theta m)\leq C r^{-2\delta}\) at that vertex. Keep it on the unsummed side; the corresponding branch is negligible. Sending \(\delta\) to zero removes every artificial radial restriction. The uncolored case gives the independently normalized signed top coefficient and no additional dihedral or radial sheet factor.
\end{proof}

\paragraph{Completeness of fixed Wilson tests.}
The joint law is physically testable, not a completion of a collection of unrelated one-angle means. Set \(R_v=\operatorname{Ad}(i u_v\cdot\boldsymbol{\sigma})\) and \(Q_v=(1+R_v)/2=u_vu_v^{\mathsf T}\). For a cyclic list of distinct vertices \(v_1,\ldots,v_l\),
\[
 \prod_{j=1}^l(u_{v_j}\cdot u_{v_{j+1}})
   =\operatorname{tr}(Q_{v_1}\cdots Q_{v_l})
   =2^{-l}\sum_{A\subseteq\{1,\ldots,l\}}
                      \operatorname{tr}\!\left(\prod_{j\in A}^{\rm ordered}R_{v_j}\right).
\]
Every term on the right is a literal spin-one ribbon word, or the empty trace \(3\). Using length two gives the squared pair dots; the sign of each simple nonzero Gram cycle fixes the remaining signs after choosing a spanning-forest sign gauge. The complete Gram matrix then fixes the axes up to an orthogonal transformation. An improper transformation may be made proper by reversing all the axes, already in the deck group; rank decline creates no ambiguity in this statement. Thus a finite set of actual spin-one linear combinations separates the compact deck quotient, including zero-dot seams. Their joint polynomial moments determine its limiting law. This does not require a growing-color approximation to a finite-level angle projector.

For example, on the four-site path with two matched frames let \(R\) be their relative rotation. The simultaneous spin-one tests on \(R_{11},R_{12}\) have Haar limits
\[
       \mathbb E(4R_{11}^2-1)=\mathbb E(4R_{12}^2-1)=\tfrac13,\qquad
       \mathbb E\bigl[(4R_{11}^2-1)(4R_{12}^2-1)\bigr]=-\tfrac35 .
\]
The second assertion uses the uniform first row of \(R\) on \(S^2\), with \(\mathbb E R_{11}^2R_{12}^2=1/15\). It rules out using independent angle marginals, which would instead give \(1/9\). This is a consequence of the true simultaneous marked invariant in \eqref{eq:ang:perfect}, not merely a calculation on a proposed classical model.

\section{Hall reduction: absolute densities and the mixed angular law}
\label{sec:hall}

A deficient matching leaves macroscopic colors at sites that cannot be
treated by the low-deficit angular expansion. We first identify the positive
measure carried by these sites through an absolute determinant calculation.
We then prove that the finite, signed angular transfer joins it to the
dormant matching law. Throughout this section \(r=k+2\),
\(Z_r=S_{11}\tau_r\), and the Lie algebra is in unit-quaternion coordinates,
with inner product \(-\operatorname{Tr}(XY)/2\). In particular,
\(\operatorname{Vol}(\mathrm{SU}(2))=2\pi^2\) and the unit axial circle has
period \(2\pi\). All seam signs and fibre periods refer to the stated
untwisted paired-cap surgery.

\subsection{Intrinsic Hall data and the active cap measure}
\label{sec:hall:data}

Let \(N\) be the number of backbone vertices, \(\nu\) its matching number,
and let \(S_*\) consist of the vertices exposable by some maximum matching.
Set
\[
 T_*=N_G(S_*),\qquad I=G\setminus(S_*\cup T_*),\qquad
 a=|T_*|,\quad b=|S_*|,\quad d=|I|/2,\qquad \nu=a+d .
\]
The inactive forest \(I\) has a unique perfect matching \(M_I\).
Equivalently, begin with a maximum matching saturating a minimum cover
and follow alternating paths from the unmatched vertices in its
complement. The reached complement vertices are precisely \(S_*\), and
their reached cover neighbors are \(T_*\). Each \(t\in T_*\) has at least
two distinct neighbors in \(S_*\): its matching partner and its
predecessor on a path from an exposed vertex. Also, no edge joins \(S_*\)
to \(I\). These properties give the minimum cover \(T_*\cup A_I\), where
\(A_I,B_I\) are the two bipartitions of the dormant forest. The
bipartitions may be chosen independently in each component. If \(a_jb_j\)
is the \(j\)-th matched edge, an unmatched \(a_i b_j\) gives the relation
\(i\to j\). Write
\[
 \Omega_I=\{u\in[0,1]^d:u_i\le u_j\text{ whenever }i\to j\},\qquad
 \iota_I(u)=(u_j,1-u_j)_{j=1}^d .
\]
The coordinates in \(\iota_I\) retain their physical vertex labels.
Reversing a component's bipartition reflects its root coordinates and
does not change any resulting labelled measure.

For a noncentral fibre \(h_s\) put \(x_s=1-\arg(h_s)/\pi\).
At \(t\in T_*\), denote by \(P_t\) the ordered product of the incoming
noncentral conjugacy-class variables, with the prescribed central
factors included. Set
\begin{equation}
 dM_*=\prod_{s\in S_*}2\cos^2(\pi x_s/2)\,dx_s\prod_{t\in T_*} L_t(x),
 \qquad
 L_t(x)=\mathbb E\!\left[\frac{\pi}{4|\operatorname{Im}P_t|}\right].
 \label{eq:hall:weight}
\end{equation}
The expectation uses independent probability conjugacy-class variables at
each cut active block. Successive tree gluing fixes the relative
orientations; it does not replace \(P_t\) by the product of its traces.

Here is the accompanying angular coarea rule. For the two trace-zero
cap quaternions \(A=u,B=v\) with probability sphere law \(d\Pi\), the
real part of \(Q=AB\) is \(-u\cdot v\), uniform on \([-1,1]\).
The group-Haar density of \(Q\) is therefore
\(\pi/(4|\operatorname{Im}Q|)\). At \(Q\ne\pm1\), rotation about its
axis gives the conditional uniform law of \(u\) on its perpendicular
equator. In a common local frame the resulting cap kernel at \(P_t\)
can be written
\begin{equation}
              \frac{\pi}{2}\,
          \delta\bigl(u_t\cdot\operatorname{Im}P_t\bigr)\,d\Pi(u_t).
 \label{eq:hall:cap-coarea}
\end{equation}
In fact the integral of \(\delta(u\cdot q)\) against \(d\Pi(u)\) is
\(1/(2|q|)\) for \(q\ne0\). At a central product this formula is used as
coarea in the variables of the ordered group relation, not as a
normalized pointwise equator. The conditional circle generally depends
on the direction from which the product approaches the center.

The same \(L_t\) is forced by the finite quantum characters. With
\(\chi_n(h)=\sin(n\arg h)/\sin\arg h\), the coefficient of \(\chi_n\) in
\(\pi/(4\sin\arg h)\) under probability group Haar is \(1/n\) for odd \(n\)
and zero for even \(n\). Averaging a product character over independent
conjugacy-class representatives gives
\(n\prod_s(\chi_n(h_s)/n)\), where the product contains one factor
\(\chi_n(h_s)/n\) per incoming active vertex. Thus on compact interior
macro radii
\begin{equation}
 L_t(x)=\sum_{\substack{n\ge1\\ n\ {\rm odd}}}
             \prod_{s\in S_*\cap N(t)}\frac{\chi_n(h_s)}n .
 \label{eq:hall:characters}
\end{equation}
The convergence is absolute there because \(t\) has at least two incoming
active neighbors. Formula \eqref{eq:hall:cap-coarea}, rather than the
pointwise endpoint values of \eqref{eq:hall:characters}, defines the
positive measure used in endpoint exhaustion.

\subsection{The absolute adjoint determinant line}
\label{sec:hall:torsion}

We give the normalization on a slightly more general chart, in order to
state the mixed-sign overlap without specializing the number of sites.
Let \(S\) be an independent noncentral set, \(T=N_G(S)\) and
\(I_0=V\setminus(S\cup T)\), of sizes \(b_0,a_0,n_0\).
Assume each \(t\in T\) is negative central and its actual cap tuple is
locally irreducible; this holds in particular for noncentral \(P_t\).
Work first where its fixed-periphery volume is regular, and on globally
irreducible characters. Let \(n_+\) be the number of positive centers in
\(I_0\), fix their cap-root choices, and sum the two active roots at
each \(s\). At each inactive site its invariant coefficient space \(V_i\)
is the line of its negative cap axis or the full quaternion algebra at a
positive center. Let \(\mathcal B_{I_0}\) be the symmetric projected
adjacency on \(\bigoplus_{i\in I_0}V_i\) given by the marked seam maps,
and set \(K=\ker\mathcal B_{I_0}\). All these spaces carry the indicated
unit metric. Use \(dM_{S,T}\) for the analogue of \eqref{eq:hall:weight}
and \(dP\) for the conditional normalized active-surface and bending
Haar variables. The latter are genuine cut variables on the effective
group quotient. In cohomology, use the dual \(H^2\) basis to the actual
deformation basis of \(H^1\).

\begin{theorem}[Absolute Hall determinant factorization]
\label{thm:hall:torsion}
The stack half-torsion in these coordinates is
\begin{equation}
 d\mu_{\rm stack}
 =2^{a_0+N-1-4n_+}\pi^{\,2b_0+n_0-2-n_+}\,
    dM_{S,T}\,dP\,
    |\det{}'\mathcal B_{I_0}|^{-1/2}\,d\operatorname{vol}_{K}.
 \label{eq:hall:torsion}
\end{equation}
Here \(\det{}'\) is the product of nonzero eigenvalues, with empty
product one, and \(d\operatorname{vol}_{K}\) uses the actual unit
quaternion fibre periods. The formula gives the density in all retained
harmonic variables on a fixed-rank chart. It does not claim that each
harmonic vector at a singular reference is tangent to a flat family.
If the positive roots are also summed, multiply by \(2^{n_+}\).
\end{theorem}

\begin{proof}
Use ordinary unit coefficient cellular cochains, in which a differential
from degree \(j\) to \(j+1\) contributes the absolute determinant to
power \((-1)^j\). A torus with a noncentral peripheral holonomy has only
its invariant line in cohomology, in primitive section and fibre units;
the perpendicular complex is acyclic with torsion one. A central seam
has the same unit assertion in all three directions.

At a noncentral vertex of valence \(k\), the cap-free piece is a
punctured sphere times the fibre circle. The perpendicular fibre
differential has determinant \(4\sin^2(\pi x)\), and its Euler exponent
from the base is \(-k\). Each cap solid torus contributes a
perpendicular determinant \(4\cos^2(\pi x/2)\). On the axis the product
cochains have canceling degrees. At each filling the section difference
has unit pivot and the primitive fibre evaluates by \(\pm2\) on the
core. The degree-one cap mismatch therefore gives \(1/2\) in torsion
per cap; the degree-zero and degree-two pivots are one. Thus the piece
half-torsion is
\[
                  \frac{2\cos^2(\pi x/2)}{(2\sin(\pi x))^k}.
\]
At a negative central piece the two cap cores act by perpendicular
half-turns, each with determinant \(4\). The same axial primitive
subdeterminants give \(4/2\) per cap in torsion, hence the piece
half-torsion is \(2\). In the inactive positive case the coefficient
system is trivial and the axial cap calculation occurs in three
directions, giving half-torsion \(1/8\).
For a locally irreducible negative piece, choose the cap-free base
cohomology in paired degrees using all peripheral traces, possible by
relative duality since the invariants vanish. The cap determinants
then remove exactly their own two trace lines. The same formula \(2\)
holds, with fixed-periphery cohomology retained.

The remaining Mayer--Vietoris calculation can be made with primitive
half-edge section coordinates. In degree one augment each vertex
restriction by its one section sum in the local invariant space.
Eliminate the section mismatches with identity pivots. A noncentral
vertex keeps its one common fibre period; its incident section values
use a sum-zero basis such as \(e_i-e_k\), with unit complementary lift
\(e_k\). A locally irreducible active center has no invariant fibre
period. Since \(S\) is independent and no \(S\)--\(I_0\) edge exists, the
remaining fibre adjacency is zero on the old \(S\) periods and is
exactly \(\mathcal B_{I_0}\) on the inactive periods. In degree zero the
map is the inclusion-difference \(U\) of vertex invariant lines or
algebras in the transported seam spaces. The degree-two half-edge
augmentation gives \(U^*\): one-sided lifts of the seam area sums have
determinant one, so no norm of a diagonal or sum-zero lattice is
inserted. The connecting collar cup product pairs each retained unit
fibre period with its own unit quotient dual, and the fixed-periphery
pairs use the Goldman form. Cross terms from alternate lifts are
triangular, since the connecting collars pair through the indicated
boundary periods. In the unit fibre metric the residual determinant is
therefore \(|\det{}'\mathcal B_{I_0}|^{-1/2}\), with its induced kernel
volume. This specifies the volume on a slanted kernel rather than
replacing it by the volume of any selected coordinate coefficients.

For completeness, put \(p=\sum_{s\in S}\deg s\) and \(q=N-1-p\). The
\(p\) noncentral seam circles and \(q\) central seam groups contribute
\((2\pi)^p(2\pi^2)^q\). Their effective quotient follows from the paired
\(U,U^*\) determinants, or by inserting the vertex invariant generators
as columns and taking the determinant complement. Divide by the
\(b_0\) noncentral vertex circles,
the \(n_0-n_+\) inactive circles and \(n_+\) inactive groups, and the
\(a_0\) locally irreducible centers.
The resulting stack angular factor is
\[
 \frac{2^{-a_0}(2\pi)^{p-b_0}(2\pi^2)^q}
      {(2\pi)^{n_0-n_+}(2\pi^2)^{n_+}}.
\]
Indeed the kernel of the combined local conjugation is precisely the
single global center when the closed tuple is irreducible. Dividing by
the product of the full local stabilizers includes that final stack
factor; dividing by it a second time would be incorrect.

The active unit period changes to \(\pi\,dx_s\). If \(k_t\) is the
number of active neighbors of a negative center, its absolute reduced
Goldman volume is
\[
   2^{k_t+1}\pi^{k_t-2}
        \left(\prod_{s\sim t,\ s\in S}\sin(\pi x_s)\right)L_t .
\]
For this last normalization use the group-valued volume formula
\cite{AMW}: a noncentral conjugacy class has volume
\(2\pi\sin\theta\); the two traceless classes each have volume \(2\pi\);
and the coarse fixed-product reduction uses
\(2/\operatorname{Vol}(\mathrm{SU}(2))^2\) times the convolution at
identity in these probability-Haar conventions. Formula
\eqref{eq:hall:cap-coarea} supplies its convolution factor.
Equivalently this is the relative cup pairing on the punctured
sphere followed by one primitive fibre product. Multiplying these
volumes and all block factors, cancel the active incidence sines,
use \(\sum_t k_t=p\) and \(p+q=N-1\), and include the \(2^{b_0}\)
true noncentral root sheets. This gives \eqref{eq:hall:torsion}.
\end{proof}

In particular let \(S=S_*,T=T_*,I_0=I\), with all dormant signs
negative. For each dormant edge let \(t_e\) be the actual signed cap
axis dot. The determinant of the weighted bipartite adjacency of a
perfectly matched forest satisfies
\[
          |\det\mathcal B_I|^{1/2}=\prod_{e\in M_I}|t_e| .
\]
Every nonzero determinant monomial gives a perfect matching and two
different matchings on a forest would have an alternating cycle.
On its resolved regular charts this gives
\begin{equation}
 d\mu_{\rm reg}
      =C_{\rm abs}\,dM_*\,\frac{dP}{\prod_{e\in M_I}|t_e|},
 \qquad C_{\rm abs}=2^{a+N-1}\pi^{\,2b+2d-2}.
 \label{eq:hall:resolved-density}
\end{equation}
When a dormant matched component is orthogonal and both radial
periods at every matched edge are retained, the block determinant is
instead the corresponding literal \(\pi\,dy_i\) per signed fibre
defect, using exactly the same angular density.

Rooting the central seams makes the conditional joint dormant axes
transparent. The incoming central relative rotation of a new dormant
axis has independent group Haar law and so produces a uniform sphere
axis conditional on earlier axes, local active cut shapes and active
radii. The old ordered products see their dormant neighbors' central
fibre signs, not these cap axes. Thus the pre-Fourier dormant axis
measure is the product of spherical probabilities, modulo the common
rotation and actual cap sign actions. Its dots on tree edges are
independent uniform variables. Conditioning on \(t_e=0\) for each
\(e\in M_I\) gives an ordered orthogonal frame on that edge and the
joint law of independent Haar frames. A joint Wilson word is evaluated
in these shared axes, not in separate angle marginals. This description
does not impose independence conditional on arbitrarily prescribed
global relative orientations of old active blocks.

For the actual local action and its calibrated stable reduction
through arbitrary dormant angles and interior active locks, see
Theorems~\ref{classic:thm:hall-action} and
\ref{classic:thm:hall-stationary}. In particular the \(2\pi^2\)
conversion between the quaternion-unit curl and the action Hessian
compensates the radial coarea factor \(\pi^{2d}\). Applying the
Gaussian normalization to \eqref{eq:hall:resolved-density} gives the
absolute factor that will occur in the joint quantum law:
\begin{equation}
 p_G=(3N-2\nu-3)/2,\qquad
 c_G=C_{\rm abs}(2\pi^2)^{-p_G}
    =2^{(1-N)/2+2a+d}\pi^{1-N}.
 \label{eq:hall:normalization}
\end{equation}
The signed phase follows independently from the capped flat action;
on this Hall stratum it is \((-1)^{r(N-1)}\).

\subsection{Active endpoints and bounded-product overlap}
\label{sec:hall:overlap}

The cap measure has a useful uniform bound at genuine macro
endpoints. Put \(\alpha_s=\min(x_s,1-x_s)\). On the open cube,
\[
       L_t(x)\asymp_G
        \frac{1}{\max_{s\in S_*\cap N(t)}\alpha_s}.
\]
For the upper bound average last the class with largest reduced
angle \(\alpha\), using radians for the following calculation. For a
fixed other quaternion of angle \(\beta\), the angle of their product
has density \(\sin\theta/(2\sin\alpha\sin\beta)\) on an interval of
length \(2\min(\bar\alpha,\bar\beta)\), where a bar reduces angles to
\([0,\pi/2]\) by central signs. This follows directly from the uniform
relative axis dot, and extends to a central fixed input by
continuity. The integral of \(1/\sin\theta\) is bounded by a constant
times \(1/\sin\alpha\). The lower bound uses the triangle inequality
on the distance of the product to the center.

For \(\alpha_s\asymp2^{-k_s}\), let \(A_j=\{s:k_s\ge j\}\) and
\(B(A_j)=\{t:N_{S_*}(t)\subset A_j\}\). The mass of the dyadic box in
\eqref{eq:hall:weight} is bounded by a constant times
\[
 2^{-\sum_s k_s+\sum_t\min_{s\sim t}k_s}
       =2^{-\sum_{j\ge1}(|A_j|-|B(A_j)|)} .
\]
For nonempty \(A_j\), the forest on \(A_j\) and \(B(A_j)\) has at least
\(2|B(A_j)|\) incidences and at most \(|A_j|+|B(A_j)|-1\).
Thus the exponent is at least \(\max k_s\). Summing by that maximum
gives, for \(b>0\),
\begin{equation}
 \int_{\{\min_s\alpha_s<\epsilon\}}dM_*
 \le C_G\epsilon(1+|\log\epsilon|)^{b-1}.
 \label{eq:hall:classical-end}
\end{equation}
The same absolute bound holds with any bounded joint classical angular
mark. There is also no endpoint mass hidden at a lock of \(P_t\).
On \(\alpha_s\ge\epsilon\) the product of two incoming classes has
group-Haar density bounded by \(C\epsilon^{-2}/\sin\theta\).
The remaining input products translate this bound. The integral of
the translated inverse-distance factor times \(1/\sin\theta\) over a
ball of radius \(\delta\) at either center is \(O(\epsilon^{-2}\delta)\),
uniformly over the translate. If the two poles are close to the ball,
apply Cauchy--Schwarz and the \(O(\delta)\) inverse-square integrals in
dimension three; otherwise integrate directly. Bounding the other
\(a-1\) cap factors by \(C\epsilon^{-(a-1)}\), and using
\eqref{eq:hall:classical-end} on the complement, gives a uniform
vanishing mass, for example
\(O(\delta^{1/(a+2)}(1+|\log\delta|)^{b-1})\), for the union of these
locks when \(a>0\). The statement does not require a pointwise choice
of conditional angular law at a collapsed equator.

We next make the normalization of all bounded physical transition
regions explicit. Let \(W\) be a fixed collection of integer-spin
Wilson words piercing only positive dormant cap disks. The words may
join dormant components. Use the exterior-compatible tree-corridor
routing: auxiliary tracks run on one side of each clasp corridor,
above the old forced sublink, and descend only at their dormant cap
ports. In the diagram with the external \(T_*\)--\(I\) Hopf clasps
removed, the ordinary forced sublink can be placed below a plane
separating it from the dormant and auxiliary ribbons. This is a
specified class of actual framed links, not an identification of all
embeddings of an abstract word. It satisfies the cut hypotheses of
Lemma~\ref{lem:ang:network}. The cap-word evaluation on the central
dormant axes is \(P_W\) in \eqref{eq:ang:angular-kernel}. In particular,
there are no old active cap passages or deliberately retained
linkings with a macroscopic ribbon in the separated diagram.

There is also an exact stronger cut identity for words routed
entirely inside individual connected dormant components. Contract
those components to vertices of the incidence tree and prune leaves.
Cut the one simple backbone ribbon at the exterior clasp, keeping
the clasp outside the cut. A closed decorated leaf then gives an
endomorphism of a simple object, hence its normalized standalone
scalar. Reattaching contributes just the ordinary normalized Hopf
factor. Repeating the cut gives exactly the product of the standalone
decorated dormant transfers and the unchanged outside scalar Hopf
factors, at every level. This argument never divides by a scalar
Hopf entry, which can vanish. A connecting word can place a nonscalar
auxiliary packet on that cut, which is why its proof uses the
conditional forced exhaustion below.

Let \(E'\) be any subset of the unmatched dormant arrows. Keep a fixed
compact positive product annulus with smooth profile \(\psi_j(w_j)\)
at each physical matched scale
\(w_j=m_{a_j}m_{b_j}/r\), and a fixed compact positive profile
\(\chi_e(v_e)\) at each \(v_e=m_{a_i}m_{b_j}/r\), \(e=(i\to j)\in E'\).
Contract the \(E'\) arrows on the matching tree. Let \(p=d-|E'|\),
let \(\Omega_{E'}\) be the order polytope on the \(p\) components, and
let \(\iota_{E'}\) insert the corresponding exponent pair
\((u_C,1-u_C)\) at every physical matched pair in its component $C$.
At the other unmatched edges one may use fixed bounded upper-product
gates equal to one near zero; alternatively support the smooth
logarithmic mark \(F\) in the strict interior chamber for the
remaining order comparisons. Use fixed negative-centre profiles
equal to one near zero on the forced and dormant sites and a fixed
smooth macro mark \(\eta\) with compact interior support.

\begin{theorem}[Absolute mixed overlap at bounded physical products]
\label{thm:hall:overlap}
For the compatible literal words and these marks,
\begin{align}
 &\lim_{r\to\infty}
 \frac{(-1)^{r(N-1)}Z_{r,W}[\eta,F,\psi,\chi]}
      {2^d c_G\,r^{p_G}(\log r)^p}
 =\left(\int\eta\,dM_*\right)
  \left(\int_{\Omega_{E'}}F(\iota_{E'}(u))\,d^pu\right)
 \notag\\[-0.2em]
 &\quad{}\times
  \int_{\mathbb R_+^{d+|E'|}}
    \mathcal K_{W,I}(\mathbf w_M,\mathbf v_{E'},0_{\rm rest})
    \prod_j\psi_j(w_j)\,dw_j
    \prod_{e\in E'}\chi_e(v_e)\,\frac{dv_e}{v_e}.
 \label{eq:hall:overlap}
\end{align}
Here \(\mathcal K_{W,I}\) is the actual joint kernel of
Theorem~\ref{thm:ang:compact}, on the dormant sphere variables with the
original shared words; it is not a product of single-edge observables.
\end{theorem}

The exhaustion proved below justifies first fixing a finite forced tuple,
also for connecting words. To see the normalization, take
\(A=T_*\cup A_I,\ B=S_*\cup B_I,\ \alpha=|B|=N-\nu\).
If \(\mathcal S_{r,W}\) denotes the sum with vertex factors
\(1+\cos(\pi m_v/r)\) and the normalized categorical kernel of
\eqref{eq:ang:normalized-sum}, put
\(X_{r,W}=(2/r)^\alpha 2^{-\nu}\mathcal S_{r,W}\).
The exact external factor in this normalization is asymptotic to
\[
 (-1)^{r(N-1)}2^{(3N+1)/2-2\alpha}\pi^{1-N}r^{p_G}
   =(-1)^{r(N-1)}2^d c_G r^{p_G}.
\]
At fixed forced labels the macro summations have ordinary odd-grid
Riemann limits with conditional density
\(2\cos^2(\pi x_s/2)\,dx_s\) and one normalized character for every old
active incidence. Summing the forced characters gives
\eqref{eq:hall:characters}. At every dormant matched partner, the
odd grid of \(w_j\) has spacing \(2m_{a_j}/r\), and the normalized
conditional mass is \(2\,dw_j/m_{a_j}\). The annuli on \(E'\) bound
all ratios between anchors in each selected component. Keep its
one root anchor and its remaining edge log ratios. Each nonroot
anchor contributes \(dR_e/(2R_e)\), \(R_e=m_{a_i}/m_{a_j}\);
the log-ratio incidence Jacobian is one, by successive leaf
elimination on a tree. Each root yields \((\log r)\,du_C/2\).
The replacement \(v_e=R_e w_j\) then gives the physical
\(dv_e/v_e\), including when several arrows use the same matched
partner. Since \(p+|E'|=d\), the powers of two cancel exactly.

On separated compact log cells all these meshes tend to zero and
all dormant products in the kernel are bounded. The true packet
limit in Theorem~\ref{thm:ang:compact} is uniform there. On the
allowed order chamber the other products tend to zero; on violated
cells they are removed by the stipulated gates. The boundary strips
have arbitrarily small normalized harmonic mass, with uniformly
bounded compact-annulus conditional transfers. This gives
\eqref{eq:hall:overlap}, first for separated smooth tests and then
for smooth joint masks by the usual tensor expansion or mixed
variation bounds. A single measured intersection of several
positive annuli is integrated in its complete \(E'\) chart. Its
contribution must not be replaced by the sum of independently
integrated single-seam residues.

\subsection{Ungated active-Hall angular gluing}
\label{sec:hall:leading}

\begin{theorem}[Ungated literal joint moments for the intrinsic Hall sector]
\label{thm:hall:main}
Let \(\eta\) be any fixed smooth joint test of the literal ordinary
radii of \(S_*\), including both endpoints, and let \(F\) be any
fixed smooth joint test of the \(2d\) literal logarithmic radii of
\(I\). For the exterior-compatible literal words above,
\begin{align}
 Z_{r,W}[\eta,F]
 &=(-1)^{r(N-1)}c_G\,r^{p_G}(\log r)^d
 \left[
       \left(\int\eta\,dM_*\right)
       \left(\int_{\Omega_I}F(\iota_I(u))\,du\right)
       \mathbb E_{\rm frame} P_W+o(1)\right].
 \label{eq:hall:main}
\end{align}
There are no original forced-color prefixes, radial endpoint,
product-scale, or Abel gates. The expectation is the genuine
joint law of independent ordered orthogonal Haar frames on
\(M_I\), modulo common rotation and the cap sign actions.
Ordinary smooth radius arguments on \(T_*\cup I\), if included in a
joint test, may be evaluated at their all-negative centres in the
limiting coefficient. The theorem concerns fixed tests of the
signed invariant, not convergence in unsigned total variation.
\end{theorem}

\begin{proof}
Use the finite-rank network of Lemma~\ref{lem:ang:network} and the
signed full-alcove prefix estimate of Lemma~\ref{lem:ang:prefix}.
Expand in its finitely many rank labels. Fold each color independently
into \([1,r/2]\); the reflected parity is fixed on each branch.
The vertex factors have bounded sup and sampled variation, including
any interval prefix, and the unaltered scalar cap profile is
\(1\pm\cos(\pi n/r)\). Sum \(B\) conditionally first. In \(X_{r,W}\)
normalization a \(B\)-site \(z\) contributes at most \(C/M_z\), where
\(M_z\) is its largest folded neighbor. For its matched anchor
\(n_j\) extract \(1/n_j\), leaving at most \(C n_j/M_z\le C\).
All \(A\)--\(A\) coefficients are bounded. This is an estimate after
each signed \(B\) transfer, not a claim about the absolute variation
of the original color array.

For every \(t\in T_*\) choose the alternating path towards an exposed
vertex. Along that path a matched partner contributes at most a
constant times \(n_{\rm next}/n_{\rm previous}\), and the final
unmatched partner gives \(1/n_{\rm last}\). Their product is bounded
by \(C/n_t\). The factors along each chosen path are distinct, and all
of them lie in the old forced part. Take the geometric mean over
\(T_*\) and restore the exterior harmonic factors. For \(a>0\) the
resulting bound on the remaining \(A\)-summand is
\begin{equation}
 C\prod_{t\in T_*} n_t^{-1-1/a}\,
            \prod_{j\in A_I}n_j^{-1}.
 \label{eq:hall:uniform-forcing}
\end{equation}
If \(a=0\) only the dormant harmonic part is present. This holds on
every folded parity and for all rectangular prefixes. In
particular a forced folded prefix \(n_t\le H\) may be inserted at
uniform normalized cost \(O(H^{-1/a})\). An originally reflected
forced vertex within a bounded folded prefix carries its untouched
\(1-\cos(\pi n_t/r)=O_H(r^{-2})\) factor and is negligible. Thus the
relevant fixed tuple can be taken in the original low-deficit
branch. Mixed variation expansion against prefixes gives the same
statements for the fixed smooth joint tests in the theorem.

This estimate also controls old macro endpoints without a compact
interior assumption. With the other forced labels at most \(H\),
the original endpoint condition at any \(s\in S_*\) restricts its
independent \(B\)-sum to at most \(O(\epsilon r+1)\) entries. Its
normalized transfer is then bounded by
\(C(\epsilon+r^{-1})\le CH(\epsilon+r^{-1})/M_s\).
Keep this improving factor in
\eqref{eq:hall:uniform-forcing} and add the high-\(H\) tail.
Optimizing gives, for \(a>0\), a bound on the normalized magnitude of
the \emph{signed} endpoint contribution of
\begin{equation}
                       C(\epsilon+r^{-1})^{1/(a+1)} .
 \label{eq:hall:quantum-end}
\end{equation}
Here the normalization is \(r^{p_G}(1+\log r)^d\). The bound
includes a finite union of endpoint bands, possibly shrinking with
the level. It is distinct from the positive classical exhaustion
\eqref{eq:hall:classical-end}.

At fixed forced \(H\), choose a dormant bipartition so that a given
dormant endpoint is unsummed. Its folded harmonic strips
\(n<r^\zeta\) and \(n>r^{1-\zeta}\) cost normalized
\(O(\zeta+1/\log r)\). For intersections choose one such restriction
as unsummed and use the signed prefix estimates for the other
ones. If an originally reflected dormant site lies in the remaining
mesoscopic rectangle, its untouched \(1-\cos\) supplies
\(O(r^{-2\zeta})\) in the same orientation. Thus only the original
negative mesoscopic rectangle remains. In a strict violation of one
dormant Hall inequality by logarithmic gap \(\gamma>0\), the factor
\(n_j/M_{b_j}\) loses \(O(r^{-\gamma})\); thin comparison strips cost
their log width. On a strict allowed chamber each dormant matched
anchor dominates its incident labels by a power, including the
bounded forced ones. The signed tail estimate removes matched
products \(w_j>K\) at normalized \(O(1/K)\) cost. The volume of
\(w_j<1/K\) is \(O(1/K)+o(1)\) by the odd mesh on separated log
cells. All of these cutoffs have now been justified in the
original exact signed diagram, prior to any approximation of an
attaching matrix.

In a retained compact product cell each \(T_*\)--\(I\) attaching
monodromy tends to the identity in the unitary ribbon category.
For fixed forced color \(n\le H\) and dormant deficit \(m=o(r)\),
the balancing phase on a fused spin \(j+s\), \(|s|\le(n-1)/2\),
differs from one by \(O((Hm+H^2)/r)\). This bounds the operator
norm because fusion channels are orthogonal. Replacing these
finitely many monodromies therefore changes the normalized compact
cell by \(o(1)\): the normalized quantum traces and the bounded
fixed packets do not amplify the error, and the absolute
compact-annulus radial count was already bounded. There is no
absolute summation of this perturbation over unbounded products.
With the chosen exterior routing the resulting old sublink is
split and undecorated. The old macro bracket remains exactly the
ordinary normalized Hopf bracket at its original values. The
dormant calculation is now on the disconnected heavy forest with
all original fixed auxiliary contractions, so the actual
coherent-state kernel of Theorem~\ref{thm:ang:compact} retains their
joint \(P_W\). This proves the fixed-prefix assertion in
Theorem~\ref{thm:hall:overlap}, including connecting words.

For each fixed old odd color \(n\) the normalized bracket is
\[
 \frac{\sin(\pi/r)}{\sin(\pi n/r)}
                      \frac{\sin(\pi n x)}{\sin(\pi x)} .
\]
It is uniformly bounded for the fixed \(n\) on the full macro
interval, with removable values at both endpoints. Thus its Riemann
limit with the actual cap factor is the ordinary normalized
character integral. The limit as \(H\to\infty\) is
\eqref{eq:hall:weight} by \eqref{eq:hall:characters} on compact
macro windows. On the full box the same identification follows
from \eqref{eq:hall:uniform-forcing},
\eqref{eq:hall:quantum-end}, and the positive exhaustion
\eqref{eq:hall:classical-end}. This removes the old macro and
fixed hub gates without referring to a limiting normalized active
angle at a degenerate cap relator.

Finally use the previously controlled signed dormant tails to
let the compact products exhaust the positive half-lines in the
strict Hall chamber. In the bounded-product joint kernel,
independent cap sign flips realize any assignment of signs to
the edge dots of a forest and leave \(P_W\) invariant. Thus the
matched Fourier integrations may be taken with real cosines.
For one matched edge the Abel kernel
\[
       \int_0^\infty e^{-\epsilon w}\cos(\pi wt)\,dw
          =\frac{\epsilon}{\epsilon^2+\pi^2t^2}
\]
tends to \(\delta_0(t)\) with mass one on \(\mathbb R\). The signed
tail bound identifies this auxiliary Abel value with the
ungated physical leading integral. Each bare dot probability
\(dt/2\) contributes \(1/2\), canceling the \(2^d\) in
\eqref{eq:hall:overlap}. Conditional Haar at the matched zero
dots gives precisely the joint frames specified in the theorem.
Exhausting the compact chambers and the old forced tuples yields
\eqref{eq:hall:main}. This is also the classical absolute
normalization in \eqref{eq:hall:normalization}, obtained by a
different computation.
\end{proof}

The genuine joint spin-one words of Section~\ref{sec:ang:main}
therefore determine the normalized dormant law on its
unprojected frame quotient, even across components in the
compatible routing. For example, a single edge-band spin-one
character \(4t^2-1\) has normalized leading mean \(-1\) on a
matched edge and \(1/3\) on an unmatched dormant edge; products
of distinct non-edge observables must still be taken in the shared
frame law. In contrast, positive physical product gates
\(\psi\) supported in \(1<w<2\) on a matched edge can give a
negative vacuum overlap coefficient in
\eqref{eq:hall:overlap}, from the negative sinc on that interval.
The positive frame measure after signed integration is therefore
not a positive joint measure of unintegrated product and angle
data. The theorem does not include arbitrary active angular
Wilson insertions, a route that truly links the separated forced
sublink, or growing packet colors. Nor does it assign lower-order
coefficients of a globally ungated integral to all other
character strata.

\subsection{A deficient five-vertex validation}
\label{sec:hall:broom}

The first mixed diagram with a nontrivial dormant matching is
the five-vertex tree with edges
\(s_1c,s_2c,cd,de\). Its intrinsic sets are
\(S_*=\{s_1,s_2\}, T_*=\{c\}, I=\{d,e\}\).
Here \(p_G=4\) and \(c_G=2/\pi^4\). At fixed compact old macro radii
the measured dormant regular negative branch has half-torsion
\[
    32\pi^4\,dM_*\,
       \frac{d\Pi(u_d)d\Pi(u_e)}{|u_d\cdot u_e|}.
\]
At the perpendicular overlap the replacement by \(\pi^2\,dy_d\,dy_e\)
and the calibrated cubic integration give
\((2/\pi^4)r^4\log r\,dM_*\).
This is also a direct finite check without the packet theorem.
The two compact old macro transfers decay rapidly in the odd hub
index by ordinary sine summation. At fixed hub index \(n\), the
remaining terminal transfer is
\((1+O(1/m))/m\) at its adjacent dormant label \(m\) in negative
trace windows equal to one at their endpoints. Its middle sinc
changes the harmonic sum by \(O(1+\log(1+n))\), so the coefficient
is the odd harmonic half-log times the old forced character
sum. The latter is exactly \(\int dM_*\) with the chosen macro
tests. This recovers the absolute coefficient and a remainder
\(O(r^4)\).

As a normalization check across a true positive end, let
\(h_d=-1,h_e=+1\). For each fixed positive \(e\) root the dormant
adjacency on the line and full Lie algebra has nonzero eigenvalues
\(\pm1\). Its kernel is the actual plane
\(v=\log(h_e)/\pi\perp u_d\); by
\eqref{eq:hall:torsion} its boundary density is
\(2\pi^5dM_*\,d\Pi(u_d)\,d^2v\).
There are two positive root choices. Writing \(y=1-|v|\), the
continued activated cover has exact radial stack density
\(16\pi^5\cot(\pi y/2)\,dy\,dM_*\), whose endpoint agrees with
the plane density after polar integration. With a smooth literal
upper-end mark \(\eta_e(y)\) and the other dormant and hub
negative windows as above, the exact quantum leading term is
\[
      \frac{r^4}{\pi^3}
           \left(\int\eta\,dM_*\right)
           \left(\int_0^1\eta_e(y)\cot(\pi y/2)\,dy\right)
      +O(r^3).
\]
Indeed the \(e\)-transfer at upper \(y\) is the sine transform of the
odd-reflected smooth cap profile and is rapidly summable in
the remaining lower index. The attaching sinc can be replaced by
one at \(O(r^{-1})\) cost in the normalized sum. The odd sine
series gives the e integral and the same forced character
factor as before. There is no new logarithmic residue.

If the other dormant site is positive instead, its opening changes
the hub product to the actual three-input ordered product. For the
activated cover write
\[
 dM'=\left(\prod_{s=s_1,s_2,d}2\cos^2(\pi x_s/2)\,dx_s\right)
          L_c(x_{s_1},x_{s_2},x_d)\,L_e(x_d),
 \qquad L_e(y)=\frac{\pi}{4\sin(\pi y)} .
\]
For a literal \(d\) upper-end mark \(\eta_d\) and \(e,c\) negative
windows the corresponding coefficient is
\[
                4\pi^{-4}r^4\int\eta\,\eta_d\,dM'+O(r^3).
\]
One can verify the exact coefficient directly: fold the upper \(d\)
label \(m\) to \(r-m\) in the terminal transfer, whose error is
\(O((r-m)^{-1})\). For odd hub \(n\) the middle sine is
\(O(n(r-m)/r)\), so this error is summable at \(O(r^{-1})\) in the
normalized sum. Its odd-mesh integral supplies the third normalized
character at \(c\) with the factor \(2\cos^2(\pi y/2)L_e(y)\).
The two old compact macro profiles make the resulting triple
character series summable all the way to the positive \(d\) endpoint.

If both dormant sites are positive they form a separate
four-dimensional smooth chart: their full projected adjacency
is invertible and there is no nearby noncentral dormant site.
Equation~\eqref{eq:hall:torsion} gives total stack
\((\pi^2/2)dM_*\) after the four positive roots. Its action is
\(1/2\) modulo one. With both positive windows equal to one at
their endpoints the exact coefficient is
\((-1)^r r^2\int\eta\,dM_*/(8\pi^2)+O(r)\).
For the finite check fold both dormant labels and write
\(m=r-i\) on \(d\). The terminal sine transfer has the double
reflection factor \((-1)^r\) and the folded transform of
\[
                 H(z)=\frac{\pi}{2}\eta_e(1-z)\tan(\pi z/2).
\]
The middle sinc supplies \(i/r+O(r^{-2})\) at fixed odd hub \(n\)
and fixed \(i\), and the remaining dormant profile supplies
\(\pi^2/(4r)+O(r^{-2})\). The fixed-parity sine derivative identity
\[
 \sum_{\substack{i>0\\ i\equiv r-1\ (2)}} i\,\frac2\pi
      \int_0^1 H(z)\sin(\pi i z)\,dz
            =\frac{H'(0)}{2\pi^2}=\frac18
\]
follows by reflecting the sine series and using \(H'(1)=0\).
Smooth odd continuation at the reflected end makes the weighted
transforms summable uniformly, so the same estimates give the
indicated remainder. This separates the two-positive chart from
the one-positive smooth endings of leading regular covers; it
is not an additional Hall logarithmic residue.

\section{Synthesis and scope}\label{intro:conclusion}
The common structure is a matching decomposition of the tree, together with the exact phase split. The exposed fibres carry the macroscopic class-convolution data. The remaining matched pairs carry a radial product and a conjugate relative axis, and their cutoffs are coupled. In the classical reduced action this coupling appears as a graph quadratic in the fibre periods with angular coefficients. In the exact finite-level transfer it appears as a signed Fourier kernel. The independent determinant normalization and the quantum transfer give the same leading measure and absolute coefficient.

The all-orders expansion has a separate role: after projection to the actual radial and cap data, it fixes every coefficient in each of the two actions. It gives sharper global information than a vacuum fit or a single saddle computation. The joint Wilson laws and transition formulas retain additional angular information and give nontrivial tests of the classical weights. Higher-order angular coefficients associated to all individual strata, in particular unrestricted insertions probing the active skeleton, are left open.

\paragraph{Checks on the absolute convention.}
For the two-site tree \(H\), the one dormant matched pair gives a unit interval of scales, and the vacuum leading term in \(Z_r\) is
\[
 Z_r(H)\sim (-1)^r\frac{\sqrt2}{\pi}\,r^{1/2}\log r.
\]
For the four-site path, the matched scale region is a triangle of volume \(1/2\), so
\[
 Z_r(P_4)\sim (-1)^r\frac{1}{\sqrt2\,\pi^3}\,r^{5/2}(\log r)^2.
\]
Each formula is in the same \(Z_r(S^2\times S^1)=1\) normalization. In a five-site broom with two leaves directly attached to the hub and a two-site tail, the exposed pair of leaves stays in the forced measure. The leading density on a compact old-macro chart is \((2/\pi^4)r^4\log r\) times that measure. Imposing a negative-central collar on every charge would discard the exposed radii and miss this mechanism.

\begin{thebibliography}{30}
\bibitem{RT1990} N.~Reshetikhin and V.~G.~Turaev, Ribbon graphs and their invariants derived from quantum groups, \emph{Commun. Math. Phys.} \textbf{127} (1990), 1--26, \href{https://doi.org/10.1007/BF02096491}{doi:10.1007/BF02096491}.
\bibitem{GPPV} S.~Gukov, D.~Pei, P.~Putrov, and C.~Vafa, BPS spectra and 3-manifold invariants, \emph{J. Knot Theory Ramifications} \textbf{29} (2020), 2040003, \href{https://arxiv.org/abs/1701.06567}{arXiv:1701.06567}.
\bibitem{BV} N.~Berline and M.~Vergne, Local asymptotic Euler--Maclaurin expansion for Riemann sums over a semi-rational polyhedron, \href{https://arxiv.org/abs/1502.01671}{arXiv:1502.01671}.
\bibitem{BK} A.~A.~Buchheit and T.~Ke{\ss}ler, Singular Euler--Maclaurin expansion on multidimensional lattices, \href{https://arxiv.org/abs/2102.10941}{arXiv:2102.10941}.
\bibitem{AMW} A.~Alekseev, E.~Meinrenken, and C.~Woodward, Duistermaat--Heckman distributions for group valued moment maps, \emph{Geom. Funct. Anal.} \textbf{12} (2002), 1--31, \href{https://arxiv.org/abs/math/9903087}{arXiv:math/9903087}.
\bibitem{JeffreyMare} L.~C.~Jeffrey and A.-L.~Mare, Products of conjugacy classes in \(\mathrm{SU}(2)\), \href{https://arxiv.org/abs/math/0211424}{arXiv:math/0211424}.
\bibitem{Toffoli} E.~Toffoli, The Atiyah--Patodi--Singer rho invariant and signatures of links, \href{https://www.cambridge.org/core/journals/proceedings-of-the-edinburgh-mathematical-society/article/atiyahpatodisinger-rho-invariant-and-signatures-of-links/4BC45AB3FD65076115A82564AAAB23E0}{journal edition}.
\bibitem{Hansen} S.~K.~Hansen, Analytic asymptotic expansions of the Reshetikhin--Turaev invariants of Seifert 3-manifolds for \(\mathrm{SU}(2)\), \href{https://arxiv.org/abs/math/0510549}{arXiv:math/0510549}.
\bibitem{Auckly} D.~R.~Auckly, Topological methods to compute Chern--Simons invariants, \emph{Math. Proc. Cambridge Philos. Soc.} \textbf{115} (1994), 229--251, \href{https://www.math.ksu.edu/about/personal/dav/topmethCS.pdf}{paper}.
\bibitem{Adams} D.~H.~Adams, The semiclassical approximation for the Chern--Simons partition function, \emph{Phys. Lett. B} \textbf{417} (1998), 53--60, \href{https://arxiv.org/abs/hep-th/9709147}{arXiv:hep-th/9709147}.
\bibitem{GoldmanMillson} W.~M.~Goldman and J.~J.~Millson, The deformation theory of representations of fundamental groups of compact K\"ahler manifolds, \emph{Publ. Math. IH\'ES} \textbf{67} (1988), 43--96, \href{https://doi.org/10.1007/BF02699127}{doi:10.1007/BF02699127}.
\bibitem{Kirillov1996} A.~Kirillov Jr., On an inner product in modular tensor categories. I, \emph{J. Amer. Math. Soc.} \textbf{9} (1996), 1135--1169, \href{https://arxiv.org/abs/q-alg/9508017}{arXiv:q-alg/9508017}.
\bibitem{TaylorWoodward2006} Y.~Taylor and C.~T.~Woodward, 6j symbols for \(U_q(\mathfrak{sl}_2)\) and non-Euclidean tetrahedra, \emph{Selecta Math. (N.S.)} \textbf{11}, 539--571, \href{https://arxiv.org/abs/math/0305113}{arXiv:math/0305113}.
\bibitem{IsmailStanton1988} M.~E.~H.~Ismail and D.~Stanton, On the Askey--Wilson and Rogers polynomials, \emph{Canad. J. Math.} \textbf{40} (1988), 1025--1045, \href{https://doi.org/10.4153/CJM-1988-041-0}{doi:10.4153/CJM-1988-041-0}.
\bibitem{vdDMiller1996} L.~van den Dries and C.~Miller, Geometric categories and o-minimal structures, \emph{Duke Math. J.} \textbf{84} (1996), 497--540, \href{https://doi.org/10.1215/S0012-7094-96-08416-1}{doi:10.1215/S0012-7094-96-08416-1}.
\bibitem{FordFulkerson1956} L.~R.~Ford Jr. and D.~R.~Fulkerson, Maximal flow through a network, \emph{Canad. J. Math.} \textbf{8} (1956), 399--404, \href{https://doi.org/10.4153/CJM-1956-045-5}{doi:10.4153/CJM-1956-045-5}.
\bibitem{BierstoneMilman1988} E.~Bierstone and P.~D.~Milman, Semianalytic and subanalytic sets, \emph{Publ. Math. IH\'ES} \textbf{67} (1988), 5--42, \href{https://numdam.org/item/PMIHES_1988__67__5_0/}{Numdam}.
\bibitem{P48} Leading logarithms in closed SU(2) Witten--Reshetikhin--Turaev invariants of capped trees, manuscript 48, \href{https://github.com/g3173028019-beep/ai-automated-quantum-gravity-research/tree/main/papers/48-wrt-leading-logarithms}{repository edition}.
\bibitem{P49} Chern--Simons phases and degenerate sewing in SU(2) WRT invariants, manuscript 49, \href{https://github.com/g3173028019-beep/ai-automated-quantum-gravity-research/tree/main/papers/49-wrt-semiclassical-sewing}{repository edition}.
\bibitem{P49B} Exact Chern--Simons normal forms for paired-cap tree plumbings, manuscript 49B, \href{https://github.com/g3173028019-beep/ai-automated-quantum-gravity-research/tree/main/papers/49b-cs-normal-forms}{repository edition}.
\bibitem{P4849} Adjoint torsion and a coupled logarithmic cutoff for paired-cap Witten--Reshetikhin--Turaev invariants, manuscript 48+49, \href{https://github.com/g3173028019-beep/ai-automated-quantum-gravity-research/tree/main/papers/48-49-wrt-torsion-bridge}{repository edition}.
\end{thebibliography}

\end{document}
